% ==========================================================
%  MASTER SUMMARY TEMPLATE — AI / ML / MATHEMATICS
%  Design Philosophy:
%    1. Readable running text is the default. Prose over boxes.
%    2. Boxes are used sparingly as visual landmarks for skimming.
%    3. Visual hierarchy: chapter > section > subsection > prose.
%    4. Every box is breakable — no orphaned whitespace at page ends.
%    5. A single consistent accent color ties all structure together.
% ==========================================================

\documentclass[11pt, oneside]{book}


% ══════════════════════════════════════════════════════════
%  PAGE GEOMETRY
% ══════════════════════════════════════════════════════════
\usepackage[
  a4paper,
  top    = 2.6cm,
  bottom = 2.1cm,
  left   = 2.5cm,
  right  = 2.5cm,
  headheight = 15pt
]{geometry}


% ══════════════════════════════════════════════════════════
%  TYPOGRAPHY
% ══════════════════════════════════════════════════════════
\usepackage[T1]{fontenc}
\usepackage[utf8]{inputenc}
\usepackage{lmodern}         % Latin Modern: improved Computer Modern
\usepackage{microtype}       % Microtypographic refinements (better justification,
                             %   fewer overfull hboxes — invaluable for math text)

% If inconsolata is not installed, simply remove this line.
% It provides a crisper monospace font for code listings.
\usepackage{inconsolata}

\setlength{\parindent}{0pt}
\setlength{\parskip}{0.62em}


% ══════════════════════════════════════════════════════════
%  MATHEMATICS
% ══════════════════════════════════════════════════════════
\usepackage{amsmath, amssymb, amsthm, mathtools}
\usepackage{bm}        % Bold math symbols: \bm{\theta}
\usepackage{bbm}       % Blackboard bold for indicator: \mathbbm{1}
\usepackage{mathrsfs}  % Script letters: \mathscr{F}, \mathscr{H}
\usepackage{nicefrac}  % Compact inline fractions: \nicefrac{1}{2}
\usepackage{cancel}    % Strike-through in math: \cancel{x}


% ══════════════════════════════════════════════════════════
%  SCIENTIFIC NOTATION
% ══════════════════════════════════════════════════════════
% siunitx — consistent typesetting of numbers and SI units.
%   Usage: \SI{1.5}{\micro\meter}  \SI{50}{\kilo\hertz}  \num{6.022e23}
\usepackage{siunitx}
\sisetup{per-mode = symbol}   % render "per unit" as e.g. m/s rather than m s^{-1}

% mhchem — chemical formulas and reaction equations.
%   Usage: \ce{H2O}   \ce{ATP -> ADP + P_i}   \ce{^{14}C}
\usepackage[version=4]{mhchem}


% ══════════════════════════════════════════════════════════
%  FIGURES, TABLES, LISTS
% ══════════════════════════════════════════════════════════
\usepackage{graphicx}
\usepackage{booktabs}         % \toprule, \midrule, \bottomrule
\usepackage{array}            % Extended column types in tabular
\usepackage{float}            % [H] float placement
\usepackage{enumitem}

\setlist[itemize]{topsep=3pt, itemsep=2pt, leftmargin=2.2em}
\setlist[enumerate]{topsep=3pt, itemsep=2pt, leftmargin=2.4em}


% ══════════════════════════════════════════════════════════
%  ALGORITHMS / PSEUDOCODE
% ══════════════════════════════════════════════════════════
% Remove this block if you do not need pseudocode.
\usepackage[ruled, vlined, linesnumbered]{algorithm2e}
\SetAlFnt{\small}
\SetAlCapFnt{\small\bfseries}
\SetCommentSty{textit}


% ══════════════════════════════════════════════════════════
%  COLORS  (a refined, academic palette)
% ══════════════════════════════════════════════════════════
\usepackage{xcolor}

% Primary structural accent — deep navy used for headings, rules, ToC
\definecolor{accent}{HTML}{1B3A5C}

% Box accent colors: frame and background, designed to be legible but
% unobtrusive so running text remains the visual focus.
\definecolor{def@frame}{HTML}{1A5276}   \definecolor{def@bg}{HTML}{EBF5FB}
\definecolor{thm@frame}{HTML}{922B21}   \definecolor{thm@bg}{HTML}{FDEDEC}
\definecolor{ex@frame} {HTML}{1E8449}   \definecolor{ex@bg} {HTML}{EAFAF1}
\definecolor{int@frame}{HTML}{9A6A00}   \definecolor{int@bg}{HTML}{FEF9E7}
\definecolor{note@frame}{HTML}{5B2C6F}  \definecolor{note@bg}{HTML}{F5EEF8}


% ══════════════════════════════════════════════════════════
%  CODE LISTINGS
% ══════════════════════════════════════════════════════════
\usepackage{listings}
\lstset{
  basicstyle       = \ttfamily\small,
  breaklines       = true,
  numbers          = left,
  numberstyle      = \tiny\color{gray!70},
  frame            = tb,
  framerule        = 0.4pt,
  rulecolor        = \color{gray!35},
  tabsize          = 2,
  showstringspaces = false,
  keywordstyle     = \color{def@frame}\bfseries,
  commentstyle     = \color{gray!55!black}\itshape,
  stringstyle      = \color{ex@frame!90!black},
  backgroundcolor  = \color{gray!5},
  xleftmargin      = 2em,
  xrightmargin     = 0.4em,
  aboveskip        = 0.8em,
  belowskip        = 0.5em
}


% ══════════════════════════════════════════════════════════
%  CUSTOM BOXES  (tcolorbox)
%
%  KEY DESIGN DECISIONS:
%  1. enhanced   — enables advanced tcolorbox features (borderline etc.)
%  2. breakable  — box may split across a page break; eliminates the ugly
%                  white gap that appears when a long box would overflow.
%  3. Left-accent bar style — a thick colored stripe on the left margin.
%                  This is legible, modern, and unobtrusive relative to
%                  a full colored frame. The box background is a very
%                  light tint, keeping the reading experience calm.
%  4. sharp corners=west — left side is flat (squared), right is rounded.
%                  This mirrors the visual logic of the left accent bar.
%  5. parbox=false — allows display math to stretch properly inside boxes.
% ══════════════════════════════════════════════════════════
\usepackage[most]{tcolorbox}

\tcbset{
  enhanced,
  breakable,
  arc             = 2.5mm,
  boxrule         = 0.5pt,
  left            = 4.5mm,
  right           = 4mm,
  top             = 2.5mm,
  bottom          = 2.5mm,
  fonttitle       = \bfseries\small,
  sharp corners   = west,
  parbox          = false
}

\newtcolorbox{defbox}[1][]{
  title  = {Definition},
  colback = def@bg, colframe = def@frame,
  borderline west = {3.5pt}{0pt}{def@frame},
  #1
}

\newtcolorbox{thmbox}[1][]{
  title  = {Theorem / Key Formula},
  colback = thm@bg, colframe = thm@frame,
  borderline west = {3.5pt}{0pt}{thm@frame},
  #1
}

\newtcolorbox{exbox}[1][]{
  title  = {Example},
  colback = ex@bg, colframe = ex@frame,
  borderline west = {3.5pt}{0pt}{ex@frame},
  #1
}

\newtcolorbox{intbox}[1][]{
  title  = {Intuition \& Mental Model},
  colback = int@bg, colframe = int@frame,
  borderline west = {3.5pt}{0pt}{int@frame},
  #1
}

\newtcolorbox{notebox}[1][]{
  title  = {Note},
  colback = note@bg, colframe = note@frame,
  borderline west = {3.5pt}{0pt}{note@frame},
  #1
}


% ══════════════════════════════════════════════════════════
%  TABLE OF CONTENTS  (tocloft)
%
%  The [titles] option prevents tocloft from conflicting
%  with titlesec's redefinition of chapter/section titles.
% ══════════════════════════════════════════════════════════
\usepackage[titles]{tocloft}

% Chapter entries: accented, bold, with dots leading to page number
\renewcommand{\cftchapfont}     {\bfseries\color{accent}}
\renewcommand{\cftchappagefont} {\bfseries\color{accent}}
\renewcommand{\cftchapleader}   {\bfseries\color{accent!40}\cftdotfill{\cftdotsep}}
\setlength{\cftbeforechapskip} {7pt}

% Section entries: standard weight, small size
\renewcommand{\cftsecfont}      {\small}
\renewcommand{\cftsecpagefont}  {\small}
\setlength{\cftsecindent}      {1.8em}
\setlength{\cftbeforesecskip}  {1.5pt}

% Subsection entries: muted gray to visually recede
\renewcommand{\cftsubsecfont}      {\small\color{gray!65!black}}
\renewcommand{\cftsubsecpagefont}  {\small\color{gray!65!black}}
\setlength{\cftsubsecindent}      {3.4em}
\setlength{\cftbeforesubsecskip}  {0.5pt}


% ══════════════════════════════════════════════════════════
%  CHAPTER & SECTION HEADINGS  (titlesec)
%
%  Chapter: large faded chapter number floated right, title
%           left-aligned beneath it, then a double rule.
%           The oversized ghost number is a quiet nod to
%           modern textbook design without being gimmicky.
%
%  Section: accent-colored, with a thin rule beneath.
%           The rule draws the eye without shouting.
%
%  Subsection: slightly smaller accent color, no rule — the
%           rule at section level is already enough structure.
% ══════════════════════════════════════════════════════════
\usepackage{titlesec}

\titleformat{\chapter}[display]
  {\normalfont\bfseries\color{accent}\huge}
  {\hfill\fontsize{68}{68}\selectfont\color{accent!13}\thechapter}
  {-10pt}
  {}
  [{\vspace{4pt}%
    {\color{accent}\rule{\linewidth}{1.6pt}}\\[-2pt]%
    {\color{accent!35}\rule{\linewidth}{0.5pt}}}]
\titlespacing*{\chapter}{0pt}{-28pt}{22pt}

\titleformat{\section}
  {\normalfont\large\bfseries\color{accent}}
  {\thesection}
  {0.8em}
  {}
  [{\vspace{1pt}\color{accent!45}\titlerule[0.45pt]}]
\titlespacing*{\section}{0pt}{3.8ex plus 1ex minus .2ex}{2.2ex plus .2ex}

\titleformat{\subsection}
  {\normalfont\normalsize\bfseries\color{accent!80!black}}
  {\thesubsection}
  {0.7em}
  {}
\titlespacing*{\subsection}{0pt}{2.6ex plus 1ex minus .2ex}{1.4ex plus .2ex}

% Run-in paragraph label — useful for named sub-items in a derivation.
\titleformat{\paragraph}[runin]
  {\normalfont\small\bfseries}
  {}
  {0em}
  {}[.\enspace]
\titlespacing{\paragraph}{0pt}{1.2ex plus .5ex minus .2ex}{0.5em}


% ══════════════════════════════════════════════════════════
%  HEADER & FOOTER
% ══════════════════════════════════════════════════════════
\usepackage{fancyhdr}
\pagestyle{fancy}
\fancyhf{}
\fancyhead[L]{\small\color{gray!70}\nouppercase{\leftmark}}
\fancyhead[R]{\small\color{gray!70}\thepage}
\renewcommand{\headrulewidth}{0.4pt}

% Chapter-opening pages use LaTeX's 'plain' style — redefine it here
% to match the rest of the document (page number at foot, no header line).
\fancypagestyle{plain}{%
  \fancyhf{}
  \fancyfoot[C]{\small\color{gray!70}\thepage}
  \renewcommand{\headrulewidth}{0pt}%
}


% ══════════════════════════════════════════════════════════
%  LINKS & REFERENCES
% ══════════════════════════════════════════════════════════
\usepackage[hidelinks]{hyperref}
\usepackage[nameinlink, capitalize]{cleveref}


% ══════════════════════════════════════════════════════════
%  DOCUMENT METADATA  (fill in for each new subject)
% ══════════════════════════════════════════════════════════
\newcommand{\coursetitle}{Deep Reinforcement Learning}
\newcommand{\documenttype}{Master Summary}
\newcommand{\authorname}{Stefan Eder}
\newcommand{\basedon}{Based on Lecture \emph{Deep Reinforcement Learning} SS26, JKU Linz}


% ══════════════════════════════════════════════════════════
%  MATHEMATICAL MACROS
%  Organized by domain. Add your own at the end of each group.
% ══════════════════════════════════════════════════════════

% ── Number Sets & Spaces ──────────────────────────────────
\newcommand{\R}{\mathbb{R}}
\newcommand{\N}{\mathbb{N}}
\newcommand{\Z}{\mathbb{Z}}
\newcommand{\C}{\mathbb{C}}
\newcommand{\Q}{\mathbb{Q}}

% ── Probability & Statistics ──────────────────────────────
\newcommand{\E}{\mathbb{E}}              % Expectation
\newcommand{\Prob}{\mathbb{P}}           % Probability measure
\newcommand{\Var}{\mathrm{Var}}
\newcommand{\Cov}{\mathrm{Cov}}
\newcommand{\Corr}{\mathrm{Corr}}
\newcommand{\iid}{\overset{\mathrm{iid}}{\sim}}
% Indicator function. Usage: \ind_{\{x > 0\}}
\newcommand{\ind}{\mathbbm{1}}

% ── Information Theory ────────────────────────────────────
% KL divergence. Usage: \KL{q}{p}  →  D_KL(q ‖ p)
\newcommand{\KL}[2]{D_{\mathrm{KL}}\!\left(#1 \,\|\, #2\right)}
\newcommand{\Ent}{\mathrm{H}}            % Entropy:         \Ent(X)
\newcommand{\MI}{\mathrm{I}}             % Mutual information: \MI(X;Y)

% ── Calculus & Optimization ───────────────────────────────
\newcommand{\argmax}{\operatorname*{arg\,max}}
\newcommand{\argmin}{\operatorname*{arg\,min}}
\newcommand{\grad}{\nabla}               % Gradient
\newcommand{\hess}{\nabla^2}             % Hessian
\newcommand{\lapl}{\Delta}               % Laplacian
% Differential (upright d). Usage: \int f(x) \dd x
\newcommand{\dd}{\,\mathrm{d}}
% Partial derivatives. Usage: \pd{f}{x}  or  \pdd{f}{x}
\newcommand{\pd}[2]{\frac{\partial #1}{\partial #2}}
\newcommand{\pdd}[2]{\frac{\partial^2 #1}{\partial #2^2}}

% ── Norms, Abs, Inner Products ────────────────────────────
\newcommand{\norm}[1]{\left\lVert #1 \right\rVert}
\newcommand{\normsq}[1]{\left\lVert #1 \right\rVert^2}
\newcommand{\abs}[1]{\left\lvert #1 \right\rvert}
\newcommand{\inner}[2]{\left\langle #1,\, #2 \right\rangle}

% ── Linear Algebra ────────────────────────────────────────
\newcommand{\T}{^{\mathsf{T}}}           % Transpose:        \mathbf{A}\T
\newcommand{\inv}{^{-1}}                 % Inverse:          A\inv
\newcommand{\pinv}{^{\dagger}}           % Pseudo-inverse:   A\pinv
\newcommand{\tr}{\operatorname{tr}}
\newcommand{\rank}{\operatorname{rank}}
\newcommand{\diag}{\operatorname{diag}}
\newcommand{\spanop}{\operatorname{span}}
\newcommand{\eye}{\mathbf{I}}            % Identity matrix
\newcommand{\zeros}{\mathbf{0}}
\newcommand{\ones}{\mathbf{1}}

% ── ML / AI Core Notation ────────────────────────────────
\newcommand{\loss}{\mathcal{L}}          % Loss:          \loss(\params)
\newcommand{\reg}{\mathcal{R}}           % Regularizer:   \reg(\params)
\newcommand{\dataset}{\mathcal{D}}       % Dataset:       \dataset
\newcommand{\params}{\boldsymbol{\theta}}   % Parameters
\newcommand{\weights}{\mathbf{w}}
\newcommand{\xvec}{\mathbf{x}}
\newcommand{\yvec}{\mathbf{y}}
\newcommand{\zvec}{\mathbf{z}}
\newcommand{\hvec}{\mathbf{h}}           % Hidden state
\newcommand{\feature}{\boldsymbol{\phi}} % Feature map

% Activation functions
\DeclareMathOperator{\softmax}{softmax}
\DeclareMathOperator{\relu}{ReLU}
\DeclareMathOperator{\sign}{sign}
% sigmoid: just use \sigma (standard Greek letter)

% ── Complexity ────────────────────────────────────────────
\newcommand{\bigO}[1]{\mathcal{O}\!\left(#1\right)}
\newcommand{\bigOm}[1]{\Omega\!\left(#1\right)}
\newcommand{\bigTh}[1]{\Theta\!\left(#1\right)}

% ── Common Distributions ──────────────────────────────────
\newcommand{\Normal}{\mathcal{N}}
\newcommand{\Uniform}{\mathcal{U}}
\newcommand{\Bernoulli}{\mathrm{Ber}}
\newcommand{\Categorical}{\mathrm{Cat}}
\newcommand{\Dirichlet}{\mathrm{Dir}}
\newcommand{\Binomial}{\mathrm{Bin}}

% ── Miscellaneous ─────────────────────────────────────────
\newcommand{\given}{\,\vert\,}           % Conditional: \Prob(A \given B)
\newcommand{\st}{\;\text{s.t.}\;}        % "subject to"
\newcommand{\eqdef}{\coloneqq}           % Definitional equality  :=
\newcommand{\eqq}{\overset{?}{=}}        % "Does this equal?"

% ── Reinforcement Learning ────────────────────────────────
\newcommand{\Sset}{\mathcal{S}}          % State space
\newcommand{\Aset}{\mathcal{A}}          % Action space
\newcommand{\vpi}{v_{\pi}}               % State-value under policy \pi
\newcommand{\qpi}{q_{\pi}}               % Action-value under policy \pi
\newcommand{\vstar}{v_{*}}               % Optimal state-value
\newcommand{\qstar}{q_{*}}               % Optimal action-value


% ==========================================================
%  BEGIN DOCUMENT
% ==========================================================
\begin{document}

\frontmatter  % Roman-numeral page numbers for front matter


% ──────────────────────────────────────────────────────────
%  TITLE PAGE
% ──────────────────────────────────────────────────────────
\thispagestyle{empty}
\vspace*{1.8cm}

\begin{center}
  {\color{accent}\rule{\linewidth}{1.8pt}}
  \vspace{0.65cm}

  {\fontsize{28}{34}\selectfont\bfseries\color{accent}\coursetitle\par}
  \vspace{0.25cm}
  {\large\color{gray!80!black}\documenttype\par}

  \vspace{0.55cm}
  {\color{accent!35}\rule{\linewidth}{0.5pt}}
  \vspace{0.55cm}

  {\large\authorname\par}
  \vspace{0.2cm}
  {\small\color{gray!70}\basedon\par}
\end{center}

\vspace{1.6cm}

\begin{center}
\begin{minipage}{0.84\textwidth}
  \small\color{gray!80!black}
  \textbf{How to use this document.}\enspace
  This is a structured summary built around \emph{running text}.
  Explanations, derivations, and logical flow between ideas live
  in the prose — colored boxes are reserved for the few things
  that must stand out at a glance.

  \medskip
  \begin{itemize}
    \item \textbf{Running text} — read for understanding, derivations, and context.
    \item \textbf{Colored boxes} — scan for key definitions, theorems, examples, and exam notes.
  \end{itemize}

  \smallskip
  Boxes lose their impact when overused. When in doubt, write prose.
\end{minipage}
\end{center}

\vfill
{\centering\small\color{gray!65} Compiled \today\par}
\clearpage


% ──────────────────────────────────────────────────────────
%  TABLE OF CONTENTS
% ──────────────────────────────────────────────────────────
\setcounter{tocdepth}{2}  % Show: chapters, sections, subsections only

\begingroup
  \let\cleardoublepage\relax
  \let\clearpage\relax
  \tableofcontents
\endgroup

\mainmatter  % Switch to Arabic numerals


% ==========================================================
%  CHAPTER 1 — RL FOUNDATIONS
% ==========================================================
\chapter{RL Foundations}

Reinforcement learning (RL) is the third major paradigm of machine learning alongside
supervised and unsupervised learning.
In \textbf{supervised learning} we learn a mapping from labeled input--output pairs;
in \textbf{unsupervised learning} we discover structure in unlabeled data.
RL is different: an \textbf{agent} interacts with an \textbf{environment}, receives
feedback in the form of a scalar reward signal, and learns to act so as to maximize
cumulative reward.
The first lecture of this course is a structured recap of the core formalism ---
the machinery developed here underpins every deep-RL method covered later.

\section{What is Reinforcement Learning?}

Reinforcement learning is a \emph{computational approach of learning to make decisions
from interactions}.
The agent is not told what the correct action is; it must discover which actions lead
to high reward by trying them.
This is what separates RL from supervised learning, where target labels are provided.

\begin{intbox}
\textbf{Five defining characteristics of RL:}
\begin{itemize}
  \item \textbf{No labeled targets} --- only a scalar reward signal, not explicit
    supervision.
  \item \textbf{Delayed feedback} --- reward may arrive long after the action that
    caused it (the \emph{credit-assignment} problem).
  \item \textbf{Active data collection} --- the agent's actions influence which states
    are visited next, so the data distribution is non-stationary and policy-dependent.
  \item \textbf{Sequential interactions} --- decisions at time $t$ affect states at
    $t+1, t+2, \dots$, creating long-term dependencies.
  \item \textbf{Goal-directed} --- the objective is to maximize cumulative reward, not
    any single step's outcome.
\end{itemize}
\end{intbox}

The overarching assumption that makes the RL framework broadly applicable is the
\emph{reward hypothesis}.

\begin{defbox}[title={Reward Hypothesis}]
All goals and purposes of an agent can be described as the maximization of the
\emph{expected cumulative reward}.
\end{defbox}

This is a conjecture, not a theorem --- but it is a remarkably useful design principle
for defining objectives across robotics, games, language models, and beyond.

\section{The Agent--Environment Interface}

At each discrete time step $t = 0, 1, 2, \dots$ the agent observes the current state
$S_t$, selects an action $A_t$, and the environment responds with a scalar reward
$R_{t+1}$ and the next state $S_{t+1}$.
The agent executes $A_t$ and receives observation $O_t$ and reward $R_t$; the
environment receives $A_t$ and emits $O_{t+1}$ and $R_{t+1}$.

\begin{defbox}[title={Trajectory}]
The agent observes state $S_t \in \Sset$ and selects action $A_t \in \Aset(S_t)$.
The environment responds with reward $R_{t+1} \in \R$ and next state $S_{t+1}$.
The resulting sequence
\[
  S_0,\; A_0,\; R_1,\; S_1,\; A_1,\; R_2,\; \dots
\]
is called a \emph{trajectory} (or history).
\end{defbox}

A useful tool throughout RL is the conditional expectation.
Recall $\E[X] = \sum_i x_i \Pr\{X = x_i\}$ and
$\E[X \mid Y=y] = \sum_i x_i \Pr\{X = x_i \mid Y = y\}$.
The \emph{law of total expectation},
$\E[X \mid Y] = \sum_z \Pr\{Z=z \mid Y\}\,\E[X \mid Y, Z=z]$,
is used repeatedly when expanding value functions over actions and next states.

To formalize the environment we use the Markov Decision Process framework.
The slides present two equivalent definitions; the first is standard in theoretical
work, the second separates transition dynamics from expected reward.

\begin{defbox}[title={Markov Decision Process}]
A \textbf{Markov Decision Process (MDP)} is a tuple
$\langle \Sset, \Aset, \mathcal{P}, \gamma \rangle$, where
\begin{itemize}
  \item $\Sset$ is the finite set of \textbf{states},
  \item $\Aset$ is the finite set of \textbf{actions},
  \item $\mathcal{P}$ specifies the \textbf{joint transition probabilities} (Markov):
  \[
    p(s', r \mid s, a) \;=\; \Pr(S_{t+1}=s',\, R_{t+1}=r \mid S_t=s,\, A_t=a),
  \]
  \item $\gamma \in [0,1]$ is the \textbf{discount factor}.
\end{itemize}

\medskip
\textbf{Alternative formulation.}
Separate state transitions from expected reward:
\[
  p(s' \mid s, a) = \sum_r p(s', r \mid s, a),
\]
\[
  r(s, a) = \E[R_{t+1} \mid S_t=s, A_t=a] = \sum_r r \sum_{s'} p(s',r \mid s,a).
\]
Both formulations define the same dynamics; use whichever is more convenient.
\end{defbox}

The crucial assumption embedded in the MDP definition is the \emph{Markov property}:
the next state and reward depend on history only through the current state--action pair.

\begin{defbox}[title={Markov Property}]
A state signal is \emph{Markov} if
\begin{align*}
  &\Pr(S_{t+1}=s', R_{t+1}=r \mid S_t=s, A_t=a,\;\text{past}) \\
  &\quad=\; \Pr(S_{t+1}=s', R_{t+1}=r \mid S_t=s, A_t=a).
\end{align*}
``The future is independent of the past given the present.''
\end{defbox}

The state must capture all information relevant to future dynamics.
In a chess game the full board position is Markov --- knowing the board and the next
move uniquely determines the distribution over future boards.
Only tracking one piece's position would \emph{not} be Markov, since the hidden
pieces still influence legal moves.

\begin{exbox}[title={A Cleaning Robot as a tiny MDP}]
The classic toy example.
A robot vacuum has two battery states $\Sset = \{\texttt{high}, \texttt{low}\}$.
From \texttt{high} it can \texttt{wait} (sit still, low reward, battery unchanged) or
\texttt{search} (clean the room, high reward, but might drain to \texttt{low}).
From \texttt{low} it can again \texttt{wait} or \texttt{search}, but now \texttt{search}
risks running out of charge entirely --- which costs a big penalty (someone has to
rescue the robot) and resets it to \texttt{high}. From \texttt{low} it can also
\texttt{recharge}: deterministic transition back to \texttt{high}, zero reward.

Concretely: $r_{\texttt{search}} > r_{\texttt{wait}} > 0$, running out gives reward $-3$,
and stochastic transition probabilities $\alpha, \beta \in (0,1)$ control how often
searching drains the battery.

Tiny as it is, the example already contains every piece of the MDP framework: a finite
state space, a state-dependent action set ($\texttt{recharge}$ only available in
\texttt{low}), stochastic transitions, and a reward function. Once you can write the
problem down like this, you can apply any of the algorithms in the rest of this course.
\end{exbox}

\section{Goals and Returns}

The agent's objective is to maximize not just the immediate reward but the long-run
\emph{return} --- the cumulative discounted sum of future rewards.

\begin{defbox}[title={Discounted Return}]
The \textbf{discounted return} from time $t$ is
\[
  G_t \;=\; R_{t+1} + \gamma R_{t+2} + \gamma^2 R_{t+3} + \cdots
        \;=\; \sum_{k=0}^{\infty} \gamma^k R_{t+k+1},
\]
where $\gamma \in [0,1)$ is the \textbf{discount factor}.
Discounting makes the sum finite for continuing tasks and encodes preference for
sooner rewards.

\medskip
Recursive identity:
\[
  G_t \;=\; R_{t+1} + \gamma\, G_{t+1}.
\]
Starting from $G_T = 0$ and iterating backwards is far more efficient than summing
forward from each $t$.
\end{defbox}

The discount factor $\gamma$ controls the effective time horizon.
As $\gamma \to 0$ only the immediate reward matters (myopic agent); as $\gamma \to 1$
distant rewards count nearly as much as immediate ones (far-sighted agent).
For \textbf{episodic tasks} that terminate at time $T$ the sum is naturally finite
and $\gamma = 1$ is permissible.
For \textbf{continuing tasks} with no terminal time, $\gamma < 1$ is required.
A sanity check: for a continuing task with constant reward $r$ and $\gamma < 1$,
the geometric series gives $G_t = r/(1-\gamma)$.

\section{Policies and Value Functions}

The agent's behavior is fully described by its \emph{policy}.

\begin{defbox}[title={Policy}]
A \textbf{policy} $\pi$ is a mapping $\pi : \Sset \times \Aset \to [0,1]$ that assigns
each state--action pair the probability of selecting $a$ in state $s$:
\[
  \pi(a \mid s) = \Pr(A_t = a \mid S_t = s),
  \qquad \sum_{a' \in \Aset} \pi(a' \mid s) = 1 \;\;\forall s \in \Sset.
\]
A policy is \emph{deterministic} if it assigns probability 1 to a single action per
state. MDP policies depend only on the current state (Markov sufficiency).
\end{defbox}

The agent seeks a policy $\pi$ that maximizes expected return.
Two value functions are the central objects for evaluating and improving policies.

\begin{defbox}[title={State-Value and Action-Value Functions}]
Given policy $\pi$, the \textbf{state-value function} and
\textbf{action-value function} are
\[
  \vpi(s) = \E_{\pi}[G_t \mid S_t=s],
  \qquad
  \qpi(s,a) = \E_{\pi}[G_t \mid S_t=s,\, A_t=a].
\]
$\vpi(s)$ measures how good it is to \emph{be} in state $s$ under $\pi$;
$\qpi(s,a)$ measures how good it is to \emph{take action $a$} in $s$ and then
follow $\pi$. They are related by:
\[
  \vpi(s) = \sum_{a} \pi(a \mid s)\, \qpi(s,a)
           = \E_{\pi}[q_\pi(S_t,A_t) \mid S_t=s].
\]
\end{defbox}

Optimal policies are defined via a partial order on value functions: $\pi \ge \pi'$
iff $\vpi(s) \ge v_{\pi'}(s)$ for all $s \in \Sset$.

\begin{defbox}[title={Optimal Value Functions}]
\[
  \vstar(s) = \max_{\pi}\, \vpi(s),
  \qquad
  \qstar(s,a) = \max_{\pi}\, \qpi(s,a).
\]
The MDP is ``solved'' when the optimal value function is known.
\end{defbox}

\begin{thmbox}[title={Optimal Policies}]
For any finite MDP:
\begin{itemize}
  \item There exists an optimal policy $\pi^*$ such that $\pi^* \ge \pi$ for
    all~$\pi$.
  \item All optimal policies achieve $\vstar$ and $\qstar$ simultaneously.
  \item There is always a \emph{deterministic} optimal policy.
\end{itemize}
\end{thmbox}

Given $\qstar$, the optimal policy is recovered by a greedy step:
\[
  \pi^*(s, a) =
  \begin{cases}
    1 & a = \argmax_{a'}\, \qstar(s, a'), \\
    0 & \text{otherwise.}
  \end{cases}
\]
If multiple actions share the maximum, any (stochastic) selection among them is optimal.
Crucially, knowing $\qstar$ yields $\pi^*$ without needing a model of
$p(s',r \mid s,a)$ --- a key motivator for action-value methods.

\section{Exploration vs.\ Exploitation}

A learning agent must balance two competing objectives.
\textbf{Exploitation}: act on current knowledge to maximize performance.
\textbf{Exploration}: try actions that increase knowledge, possibly at short-term cost.
Neither extreme is optimal: pure exploitation risks committing to a suboptimal action
before adequate data are gathered; pure exploration never leverages what is already known.

The tension is cleanest in the \emph{stateless} setting of multi-armed bandits, where
no state transitions complicate the picture.
Bandits therefore serve as the canonical sandbox for studying exploration strategies
before extending them to the full MDP setting.

\section{Multi-Armed Bandits}

\subsection{Problem setup}

\begin{defbox}[title={$k$-Armed Bandit}]
A \textbf{$k$-armed bandit} has $k$ actions (``arms'') $a \in \{1,\dots,k\}$.
There is a single state (non-associative setting): actions have no long-term
consequences on the environment beyond the immediate reward.
At each step $t$ the agent selects arm $A_t$ and observes reward $R_t$.
The \textbf{true action value} is
\[
  q(a) \;=\; \E[R_t \mid A_t = a].
\]
The agent maintains an estimate $Q_t(a)$ and a count $N_t(a)$.
True values are unknown and must be learned from experience.
\end{defbox}

A problem is \textbf{stationary} if $q(a)$ is constant over time, and
\textbf{non-stationary} if the reward distributions drift.
This distinction drives the choice of estimation method.

\subsection{Action value estimation}

The natural unbiased estimator of $q(a)$ is the sample average of observed rewards:
\[
  Q_t(a) = \frac{\displaystyle\sum_{n=1}^{t} R_n \cdot \ind_{A_n = a}}{N_t(a)}.
\]
By the law of large numbers, $Q_t(a) \to q(a)$ as $N_t(a) \to \infty$.
Rather than recomputing from scratch each step, the incremental form processes one
observation at a time:

\begin{thmbox}[title={Incremental Update}]
After selecting $A_t$ and observing $R_t$, update only the chosen arm's estimate:
\[
  Q_{t+1}(A_t) = Q_t(A_t) + \alpha_t\bigl(R_t - Q_t(A_t)\bigr),
\]
where $(R_t - Q_t(A_t))$ is the \emph{prediction error}.
\begin{itemize}
  \item $\alpha_t = 1/N_t(A_t)$: reproduces the unbiased sample average.
  \item \textbf{Constant} $\alpha \in (0,1]$: exponentially recency-weighted average
    --- older rewards count less, ideal for non-stationary problems.
\end{itemize}
All other estimates $Q_t(a)$ with $a \ne A_t$ remain unchanged.
\end{thmbox}

\subsection{Regret}

Rather than tracking absolute reward, it is informative to measure how much is
\emph{lost} by not always selecting the optimal arm.

\begin{defbox}[title={Regret}]
The \textbf{optimal value} is $q^* = \max_a q(a)$.
The \textbf{regret} of selecting action $a$ is
\[
  \Delta_a \;=\; q^* - q(a) \;\ge\; 0.
\]
The \textbf{total regret} after $t$ steps is
\[
  L_t \;=\; \sum_{n=1}^{t} \bigl(q^* - q(A_n)\bigr) \;=\; \sum_{n=1}^{t} \Delta_{A_n}.
\]
Maximizing cumulative reward is equivalent to minimizing total regret $L_t$.
\end{defbox}

A good exploration strategy keeps $L_t$ sub-linear in $t$ --- the per-step regret
eventually goes to zero.
Pure greedy can suffer linear regret if it commits early to a suboptimal arm.

\subsection{Algorithms}

Bandit algorithms split into two families. \textbf{Value-based} methods estimate
$Q_t(a)$ and derive a policy from those estimates. \textbf{Policy-based} methods skip
estimation and tune action preferences directly. Below, we walk up the value-based
ladder from naive to clever, then look at the policy-based alternative.

\paragraph{Greedy --- the cautionary tale.}
$A_t \in \argmax_a Q_t(a)$. Pure exploitation: always pick the arm with the highest
current estimate. Trivial to implement, and almost guaranteed to fail.
A typical failure mode: the truly best arm gets unlucky on its first pull, drops to
the bottom of the leaderboard, and never gets tried again --- because greedy never
revisits an arm it thinks is bad. The agent locks in on whichever arm happened to
look good early. Linear regret.

\paragraph{$\varepsilon$-greedy --- a constant trickle of exploration.}
The simplest fix: most of the time act greedily, but with probability $\varepsilon$
pick uniformly at random. Now every arm gets sampled forever, so the estimates
eventually concentrate on their true values.

\begin{defbox}[title={$\varepsilon$-Greedy Action Selection}]
\[
  A_t = \begin{cases}
    \argmax_a Q_t(a), & \text{with probability } 1 - \varepsilon, \\
    \text{uniform random from } \{1, \dots, k\}, & \text{with probability } \varepsilon.
  \end{cases}
\]
\end{defbox}

A fixed $\varepsilon > 0$ guarantees infinite exploration --- but also wastes a constant
fraction of pulls on random actions forever, even after we have figured out which arm
is best. The compromise is a \textbf{GLIE} schedule (Greedy in the Limit with Infinite
Exploration), like $\varepsilon_k = 1/k$, that decays exploration over time while still
sampling each arm infinitely often. Sub-linear regret, asymptotic convergence.

\paragraph{UCB --- explore the arms you are unsure about, not the ones you have written off.}
$\varepsilon$-greedy treats all non-greedy arms equally. That is wasteful: an arm we
have pulled $1{,}000$ times and know is bad gets the same exploration probability as
one we have pulled twice and barely understand.
UCB fixes this by adding an explicit \emph{uncertainty bonus} to each arm:

\begin{thmbox}[title={UCB Action Selection}]
\[
  A_t \in \argmax_{a} \left[\, Q_t(a) + c\, \sqrt{\frac{\ln t}{N_t(a)}}\, \right],
\]
where $c > 0$ controls how aggressively we explore.
Pull each arm at least once before applying the formula (the bonus is undefined for $N_t(a) = 0$).
\end{thmbox}

The bonus $c\sqrt{\ln t / N_t(a)}$ shrinks as $N_t(a)$ grows, but never quite to zero
--- $\ln t$ keeps inching up, so even well-understood arms occasionally get pulled
again.
The $\sqrt{\ln t / N_t(a)}$ shape is not arbitrary: it comes from a confidence
interval on the empirical mean (Hoeffding's inequality), so the bonus is roughly
``a 95\%-confidence upper bound on this arm's true value, given the data so far.''
Acting on the upper bound is the principle of \textbf{optimism in the face of
uncertainty}: when in doubt, act as if the unknown thing might be great --- because
either it \emph{is} great (and you collect the reward), or it isn't (and you learn
that and stop exploring it).

\begin{intbox}[title={Why UCB beats $\varepsilon$-greedy}]
$\varepsilon$-greedy: ``with 10\% probability, do something random.''
UCB: ``with 100\% probability, do whatever has the highest value-plus-uncertainty.''

The key difference: UCB knows \emph{which} arms it is uncertain about. It will keep
sampling an arm that might be best, but it won't waste pulls on arms it has
established are clearly bad. $\varepsilon$-greedy explores blindly; UCB explores
strategically.
\end{intbox}

\paragraph{Thompson Sampling --- gambling with posteriors.}
A Bayesian alternative: maintain a posterior distribution over each $q(a)$ instead of
a point estimate. At each step, \emph{sample} one $\tilde q(a)$ from each arm's
posterior, then act greedily on the samples.
\begin{itemize}
  \item Arms with high posterior variance (we are uncertain about them) produce widely
    varying samples, sometimes large enough to win the argmax. Those arms get pulled.
  \item Arms whose posterior is sharply peaked around a low value almost never produce
    a high sample, so they are ignored.
\end{itemize}

\begin{intbox}[title={UCB vs.\ Thompson Sampling}]
Both methods are uncertainty-aware, just with different decision rules.
\textbf{UCB} is deterministic and optimistic: ``I'll pick the arm whose \emph{best
plausible} value is highest.''
\textbf{Thompson} is probabilistic and Bayesian: ``I'll pretend one specific value
function is true (sampled from my beliefs) and act greedily on that.''
Both achieve near-optimal regret. Thompson tends to outperform UCB on real problems
because its randomization is naturally calibrated to the posterior, while UCB's
$c\sqrt{\ln t / N_t}$ bonus is just a worst-case bound.
This is the same UCB-vs.-Thompson choice we will see again in Chapter~6 when discussing
deep-RL exploration.
\end{intbox}

\medskip
\textbf{Gradient Bandits (policy-based).}

\begin{defbox}[title={Preference Parameters}]
Instead of estimating action values, gradient bandits maintain a real-valued
\textbf{preference} $H_t(a) \in \R$ for each arm.
Only \emph{differences} matter: adding a constant to all preferences leaves the
policy unchanged.
\end{defbox}

Preferences are converted to a sampling distribution via the softmax:
\[
  \pi_t(a) = \frac{\exp(H_t(a))}{\displaystyle\sum_b \exp(H_t(b))}.
\]
After observing reward $R_t$ with running baseline $\bar{R}_t$ (often the running
average reward), preferences are updated by stochastic gradient ascent on expected
reward:
\[
  H_{t+1}(a) =
  \begin{cases}
    H_t(a) + \alpha(R_t - \bar{R}_t)(1 - \pi_t(a)) & a = A_t, \\[4pt]
    H_t(a) - \alpha(R_t - \bar{R}_t)\,\pi_t(a) & a \ne A_t.
  \end{cases}
\]

\begin{intbox}[title={What the Gradient Update Does}]
Think of $\pi_t$ as a weighted lottery.
If $R_t > \bar{R}_t$ (better than average): increase $H(A_t)$, making the selected
arm more likely; slightly decrease the rest.
If $R_t < \bar{R}_t$: decrease $H(A_t)$ and increase the rest.
The factors $(1-\pi_t(A_t))$ and $\pi_t(a)$ scale the update: a surprising action
(low $\pi_t$) triggers a large adjustment; a near-certain action changes little.
Unlike value-based methods, gradient bandits have no explicit exploration rate ---
exploration emerges from the softmax sampling itself.
\end{intbox}



% ==========================================================
%  CHAPTER 2 — TABULAR VALUE-BASED METHODS
% ==========================================================
\chapter{Tabular Value-Based Methods}
\label{chap:tabular}

Lecture 02 is organized around a single central question: \emph{how do we
actually solve for an optimal policy?}
The answer depends on two things --- whether we have access to a model of the
environment, and whether we compute expectations exactly or estimate them from
samples.
Those two choices carve out the three algorithm families covered here:
Dynamic Programming, Monte Carlo, and Temporal-Difference learning.

\section{RL Taxonomy}

Before diving into algorithms, it helps to have a mental map of where each one sits.
The RL landscape can be carved up along three roughly independent axes: \emph{what}
the agent learns, \emph{whose data} it learns from, and \emph{whether it has a model}
of the environment.

\paragraph{What gets learned.}
A method is \textbf{value-based} if it learns a value function and reads the policy
off it (``in this state, the action with the highest Q-value is best, so do that'').
\textbf{Policy-based} methods skip the value function and directly tune a policy;
they answer ``what should I do?'' without ever asking ``how good is this state?''
\textbf{Actor-critic} methods do both --- they keep a policy (the actor that picks
actions) and a value function (the critic that judges them), and use the critic's
opinions to improve the actor.

\paragraph{Whose data.}
\textbf{On-policy} methods are puritans: the policy they evaluate has to be the policy
that collected the data. As soon as you change the policy, the old data is officially
stale.
\textbf{Off-policy} methods are pragmatists: the \emph{target} policy $\pi$ (what we
want to learn) and the \emph{behavior} policy $\mu$ (what actually collected the data)
can be different. That gap is what lets us reuse old experience, learn from human
demonstrations, or learn a greedy policy while still exploring.
\textbf{Offline} methods go to the extreme: a fixed dataset, no environment interaction
at all.

\paragraph{Model or no model.}
\textbf{Model-based} methods use a model of the dynamics $p(s', r \mid s, a)$ ---
either given to them or learned --- and can plan by simulating transitions.
\textbf{Model-free} methods just learn from interaction, never bothering to predict
what the environment will do next.

\begin{intbox}[title={Composing the axes}]
The three axes are largely independent, so any combination is possible:
\begin{itemize}
  \item \textbf{DQN}: value-based, off-policy, model-free.
  \item \textbf{REINFORCE}: policy-based, on-policy, model-free.
  \item \textbf{PPO / SAC}: actor-critic, on/off-policy, model-free.
  \item \textbf{AlphaZero}: actor-critic, on-policy, model-based (it has a perfect
    chess simulator and uses it to plan).
\end{itemize}
This chapter lives in the value-based + model-free corner of the cube.
\end{intbox}

\section{Bellman Equations}

The Bellman equations are the mathematical backbone of tabular RL.
They express a simple recursive insight: the value of being in a state equals
the immediate reward you expect plus the discounted value of where you end up.
This sounds almost trivial, but it is the engine that drives almost every algorithm
in this chapter.

\begin{intbox}[title={Where the recursion comes from}]
Take the definition of the return, $G_t = R_{t+1} + \gamma R_{t+2} + \gamma^2 R_{t+3} + \cdots$,
and pull out the first reward:
\[
  G_t = R_{t+1} + \gamma\, \underbrace{(R_{t+2} + \gamma R_{t+3} + \cdots)}_{= G_{t+1}}.
\]
Take expectations on both sides given $S_t = s$ and a policy $\pi$, and the value
function ``unrolls itself by one step'':
\[
  \vpi(s) = \E_\pi\!\bigl[\, R_{t+1} + \gamma\, \vpi(S_{t+1})\, \mid S_t = s\, \bigr].
\]
That is the entire content of the Bellman expectation equation. Everything below
is just spelling out the expectation as a sum over actions and next states.
\end{intbox}

\begin{thmbox}[title={Bellman Expectation and Optimality Equations}]
For any policy $\pi$, the value functions satisfy the \emph{expectation equations}:
\begin{align*}
  \vpi(s) &= \sum_a \pi(a\mid s)\!\left[r(s,a) + \gamma\sum_{s'} p(s'\mid s,a)\,\vpi(s')\right], \\[4pt]
  \qpi(s,a) &= r(s,a) + \gamma\sum_{s'} p(s'\mid s,a)\sum_{a'} \pi(a'\mid s')\,\qpi(s',a').
\end{align*}
The \emph{optimal} value functions satisfy the \emph{optimality equations}:
\begin{align*}
  \vstar(s) &= \max_a \!\left[r(s,a) + \gamma\sum_{s'} p(s'\mid s,a)\,\vstar(s')\right], \\[4pt]
  \qstar(s,a) &= r(s,a) + \gamma\sum_{s'} p(s'\mid s,a)\,\max_{a'}\qstar(s',a').
\end{align*}
\end{thmbox}

The expectation equations are \emph{linear} in the value function --- in principle
you could solve them as a linear system in $O(|\Sset|^3)$ operations.
For anything but tiny MDPs that is computationally prohibitive, so iterative methods
take over.

The optimality equations look almost the same but replace the policy-weighted
average with a $\max$.
That $\max$ makes them nonlinear, so there is no closed-form solution --- but it also
gives them a beautiful interpretation: the optimal value of a state is what you get
if you act greedily forever after this step.

\begin{intbox}[title={Why $\qstar$ is more useful than $\vstar$}]
Once you have $\qstar$, the optimal policy is free: $\pi^*(s) = \argmax_a \qstar(s, a)$.
Just look up which action has the highest action-value and take it.

If you only have $\vstar$, you have to compute
$\argmax_a [\, r(s,a) + \gamma \sum_{s'} p(s' \mid s, a)\, \vstar(s')\, ]$ ---
which requires the environment model.
This is why action-value methods are so attractive in model-free settings:
$\qstar$ tells you what to do without consulting any dynamics.
\end{intbox}

\section{Planning with Dynamic Programming}

Dynamic programming computes optimal policies \emph{given a perfect model} of
the environment.
DP uses the model to form exact expected-value backups rather than sampling.
Every DP algorithm alternates between two steps:
\textbf{policy evaluation} (compute how good the current policy is) and
\textbf{policy improvement} (make the policy better).

\subsection{Policy evaluation and improvement}

Given a fixed policy $\pi$, we want to compute its value function $\vpi$
exactly.
The trick is to treat the Bellman expectation equation not as a system of
equations to solve, but as an \emph{update rule}.
Starting from any initialization $V_0$, iterating the update
\[
  V_{k+1}(s) \;\leftarrow\; \sum_a \pi(a\mid s)\!\left[r(s,a)
    + \gamma\sum_{s'}p(s'\mid s,a)\,V_k(s')\right]
\]
for all states $s$ simultaneously converges to $\vpi$.
In practice we stop when the maximum change across states falls below a small
threshold $\varepsilon$.

Once we have $V \approx \vpi$, policy improvement is straightforward.
The \textbf{policy improvement theorem} guarantees that choosing the greedy
action at every state --- $\pi'(s) = \argmax_a\,q_\pi(s,a)$ --- always yields
a policy at least as good as $\pi$.
If $v_{\pi'}(s) = \vpi(s)$ for all $s$, the policy was already optimal.

\subsection{Policy iteration and value iteration}

\textbf{Policy iteration} repeats evaluation and improvement in alternation:
run iterative policy evaluation to convergence, extract the greedy policy,
repeat.
It is guaranteed to reach the optimal policy after a finite number of
improvements, but running evaluation fully to convergence on each iteration
is expensive.

\textbf{Value iteration} is a smarter shortcut: do not wait for evaluation to
converge.
Instead, apply a \emph{single} backup using the optimality operator --- taking
the $\max$ rather than the policy-weighted sum:
\[
  V_{k+1}(s) \;\leftarrow\; \max_a\!\left[r(s,a)
    + \gamma\sum_{s'} p(s'\mid s,a)\,V_k(s')\right].
\]
This is equivalent to interleaving one step of evaluation with one step of
improvement: every iteration both evaluates and improves.
Value iteration converges asymptotically (no finite termination guarantee) but is
typically more efficient per iteration than running full policy evaluation.

\begin{exbox}[title={Value iteration unrolling, by hand}]
Imagine a tiny three-state corridor $s_1 - s_2 - s_3$, with $s_3$ a terminal goal
giving reward $+1$ and $s_1, s_2$ giving reward $0$. Actions move deterministically
left or right. Discount $\gamma = 0.9$. Initialize $V_0 = (0, 0, 0)$.

\textbf{Iteration 1.} For each non-terminal state, take the best action.
\begin{itemize}
  \item $V_1(s_2) = \max\{\, 0 + 0.9 \cdot V_0(s_1),\; 1 + 0.9 \cdot V_0(s_3)\, \} = \mathbf{1}$ (go right).
  \item $V_1(s_1) = \max\{\, 0 + 0.9 \cdot V_0(s_2),\; \dots\, \} = 0$.
\end{itemize}
\textbf{Iteration 2.}
\begin{itemize}
  \item $V_2(s_2) = 1 + 0.9 \cdot 0 = 1$ (still optimal to step right and collect).
  \item $V_2(s_1) = 0 + 0.9 \cdot 1 = 0.9$ (one step toward $s_2$, then collect).
\end{itemize}
The reward propagates one step backward per iteration --- after $k$ iterations every
state within $k$ steps of the goal has correct value. The greedy policy ``always
move toward $s_3$'' is recovered immediately.

That is the entire algorithm: forward Bellman backups, reward seeping backward through
the state graph one step at a time, until the value function stops changing.
\end{exbox}

\begin{intbox}[title={Limitation of DP}]
Both algorithms require knowing $p(s',r\mid s,a)$ for every transition.
In most real-world settings this model is unavailable.
Model-free methods, covered next, replace the model-based expectation with samples
drawn from the environment --- giving up exactness for the ability to learn from
raw experience.
\end{intbox}

\section{Model-Free Prediction}

The prediction task is: given a policy $\pi$, estimate $\vpi$ from experience
alone, without any model.
Two fundamentally different approaches exist.

\subsection{Monte Carlo: learn from complete episodes}

Monte Carlo (MC) methods estimate $\vpi(s)$ by averaging the actual returns
observed when starting from $s$.
Since $\vpi(s) = \E_\pi[G_t \mid S_t = s]$ and $G_t$ is observable at episode
end, we can use the sample average as the estimator, updated incrementally:
\[
  V(S_t) \;\leftarrow\; V(S_t) + \alpha\bigl[G_t - V(S_t)\bigr].
\]
$G_t$ is an \emph{unbiased} estimate of $\vpi(S_t)$ --- on average it equals
the true value.
The downside is variance: $G_t$ is a sum of many random rewards, so individual
estimates can vary wildly, and the method requires waiting until the episode ends.

\subsection{Temporal-difference: learn online, one step at a time}

TD methods trade unbiasedness for immediacy.
Instead of waiting for the actual return $G_t$, TD uses the Bellman equation
to construct an estimate after just one step --- the \textbf{TD(0) target}
$R_{t+1} + \gamma V(S_{t+1})$.
The update is:

\begin{thmbox}[title={TD(0) Update Rule}]
\[
  V(S_t) \;\leftarrow\; V(S_t)
  + \alpha\underbrace{\bigl[R_{t+1} + \gamma V(S_{t+1}) - V(S_t)\bigr]}_{\text{TD error } \delta_t}.
\]
TD(0) is model-free like MC, bootstraps like DP, and can update after
\emph{every step} without waiting for episode end.
\end{thmbox}

The TD target is \emph{biased} because $V(S_{t+1})$ is only an approximation ---
we are plugging an estimate back into its own update.
But it has far lower variance than $G_t$, because only a single reward is
random rather than a whole trajectory's worth.
In practice, TD often learns faster even though it introduces this bias.

Information propagation differs too.
In MC, a surprising outcome at the end of an episode immediately updates all
states visited along the way.
In TD, that information trickles backward only one step at a time.

\subsection{Multi-step returns and TD($\lambda$)}

$n$-step TD provides a smooth interpolation between TD(0) ($n=1$) and MC
($n=\infty$) by collecting $n$ actual rewards before bootstrapping:
\[
  G_t^{(n)} = R_{t+1} + \gamma R_{t+2} + \cdots + \gamma^{n-1}R_{t+n} + \gamma^n V(S_{t+n}).
\]
Rather than committing to one $n$, TD($\lambda$) mixes \emph{all} $n$-step
returns simultaneously using geometrically decaying weights:
\[
  G_t^\lambda = (1-\lambda)\sum_{n=1}^{T-t-1}\lambda^{n-1}G_t^{(n)}
  + \lambda^{T-t-1}G_t.
\]
Setting $\lambda=0$ recovers TD(0); $\lambda=1$ recovers MC.
Intermediate values get the best of both worlds: lower variance than MC while
propagating information faster than one-step TD.
In practice the $\lambda$-return is implemented online via
\textbf{eligibility traces}, which maintain a decaying memory of recently
visited states and spread the TD error backward proportionally.

\section{Model-Free Control}

For control, we need action-value estimates $Q(s,a)$ rather than state values.
The reason is simple: to be greedy over $\vpi$, you need to compute
$\argmax_a\bigl[r(s,a) + \gamma\sum_{s'} p(s'\mid s,a)V(s')\bigr]$,
which requires a model.
Being greedy over $Q(s,a)$ is model-free: just pick $\argmax_a\,Q(s,a)$.

The general framework is \textbf{generalized policy iteration}: alternate
between approximate policy evaluation (using MC or TD to estimate $Q \approx \qpi$)
and policy improvement (act greedily with respect to $Q$).
The catch is exploration: a purely greedy policy never tries suboptimal actions,
so $Q(s,a)$ for those actions never gets updated and the policy can get stuck.
The standard fix is $\varepsilon$-greedy: with probability $\varepsilon$ take a
random action, otherwise act greedily.
To guarantee convergence to the optimal policy, we need the exploration schedule
to satisfy GLIE --- \emph{Greedy in the Limit with Infinite Exploration}:
every state-action pair must be visited infinitely often, yet the policy must
eventually converge to greedy.
An $\varepsilon$-greedy schedule with $\varepsilon_t \to 0$ and
$\sum_t \varepsilon_t = \infty$ satisfies GLIE.

\subsection{SARSA: on-policy TD control}

SARSA applies the TD idea directly to action values.
At each step we observe the transition $(S_t, A_t, R_{t+1}, S_{t+1})$,
then sample the next action $A_{t+1}$ from the current policy and update:
\[
  Q(S_t,A_t) \leftarrow Q(S_t,A_t) + \alpha\bigl[R_{t+1} + \gamma Q(S_{t+1},A_{t+1}) - Q(S_t,A_t)\bigr].
\]
The five quantities $(S_t, A_t, R_{t+1}, S_{t+1}, A_{t+1})$ give the algorithm
its name.
SARSA is \emph{on-policy}: the action $A_{t+1}$ used in the target is actually
sampled from the current $\varepsilon$-greedy policy, so the estimate converges
to $q_\pi$ for the policy being followed.

A variant, \textbf{Expected SARSA}, replaces the sampled $A_{t+1}$ with an
expectation over the policy:
\[
  Q(S_t,A_t) \leftarrow Q(S_t,A_t)
  + \alpha\!\left[R_{t+1} + \gamma\sum_{a'}\pi(a'\mid S_{t+1})Q(S_{t+1},a') - Q(S_t,A_t)\right].
\]
This reduces variance because it averages over the randomness in $A_{t+1}$
rather than sampling it.

\subsection{Q-learning: off-policy TD control}

Q-learning breaks the link between behavior and learning target.
Regardless of which action the agent \emph{actually} takes next, the update uses
the \emph{greedy} action at the next state:
\[
  Q(S_t,A_t) \leftarrow Q(S_t,A_t)
  + \alpha\!\left[R_{t+1} + \gamma\max_{a'} Q(S_{t+1},a') - Q(S_t,A_t)\right].
\]
This is a sample-based approximation of the Bellman optimality equation for $\qstar$,
and under standard conditions $Q$ converges to $\qstar$ regardless of the behavior
policy --- the behavior just needs to ensure every state-action pair gets visited.
No importance-sampling correction is needed because the $\max$ target does not
depend on the behavior policy at all.

In one sentence: SARSA evaluates what the agent \emph{actually does}; Q-learning
evaluates what the agent \emph{would do if it were greedy}.

\begin{exbox}[title={Cliff Walking: SARSA vs.\ Q-learning}]
A classic gridworld where the bottom row is a cliff: stepping into it gives a big
negative reward and resets the agent to the start. The shortest path runs along the
edge of the cliff; a longer, safer path goes through the top of the grid.

Train both algorithms with $\varepsilon$-greedy exploration ($\varepsilon = 0.1$).
\begin{itemize}
  \item \textbf{Q-learning} learns the \emph{optimal} (cliff-edge) policy: its target
    uses $\max_{a'}$, ignoring the fact that during exploration the agent will
    occasionally take a random action --- and falling off the cliff while exploring
    along the edge hurts.
    Q-learning's $Q$-values reflect the value of \emph{playing perfectly}, even when
    the actual behavior is noisy.
  \item \textbf{SARSA} learns the \emph{safe} policy: its target uses the
    next action $A_{t+1}$ that the policy actually picks --- which, with $\varepsilon$
    chance, is random. Walking the cliff edge with random actions costs a lot in
    expectation, so SARSA's $Q$-values penalize it and the policy detours through the
    safer top row.
\end{itemize}
The lesson: SARSA optimizes the policy you're currently using (including its noise);
Q-learning optimizes the greedy policy you'd play \emph{if} you stopped exploring.
Once $\varepsilon$ is annealed to zero, both converge to the same optimal solution ---
but their behavior during training is different in a way that actually matters in
risky environments.
\end{exbox}

\subsection{Maximization bias and Double Q-learning}

There is a subtle problem with the $\max$ in Q-learning.
When $Q$-values are estimated with noise --- as they always are from finite data ---
taking the maximum over noisy estimates produces a value that is systematically
\emph{too high}. The noise pushes the max upward.
This is called \textbf{maximization bias}.

\begin{intbox}[title={Why max gets it wrong (concrete numbers)}]
Suppose three actions in some state really have value $0$, but each $Q$-estimate
is true value plus noise $\Normal(0, 1)$. So on any given update:
\begin{itemize}
  \item $Q(s, a_1), Q(s, a_2), Q(s, a_3)$ are all centered at zero on average.
  \item $\E[\max_a Q(s, a)] \approx 0.85 > 0$ --- the maximum of three independent
    standard Gaussians is biased high.
\end{itemize}
$Q$-learning's TD target uses this maximum, so it inflates value estimates by
roughly $0.85$. The error compounds: inflated values get bootstrapped into
neighboring states, which inflates them too, which inflates the original state more
on the next update, and so on.
The bias is worst exactly when it hurts most --- early in training, when noise is
high and the agent commits to whichever action got lucky on the first few samples.
\end{intbox}

The fix is to \emph{decouple action selection from action evaluation}, so the noise
in selection doesn't propagate into the evaluated value.

\begin{defbox}[title={Double Q-Learning}]
Maintain two independent estimates $Q^A$ and $Q^B$.
Use one to \emph{select} the action and the other to \emph{evaluate} it:
\[
  Q^A(S_t,A_t) \leftarrow Q^A(S_t,A_t)
    + \alpha\!\left[R_{t+1} + \gamma\, Q^B\!\bigl(S_{t+1},\, \argmax_a Q^A(S_{t+1},a)\bigr) - Q^A(S_t,A_t)\right].
\]
At each step, flip a coin and update either $Q^A$ or $Q^B$; the other plays the
evaluation role.
Decoupling selection from evaluation breaks the positive correlation that drives
the bias upward: $Q^A$ might overestimate $a_1$, but $Q^B$'s independent estimate
of $a_1$ is just as likely to be too low.
\end{defbox}

\subsection{Off-policy learning and importance sampling}

In off-policy learning, the agent collects data with a \emph{behavior policy} $\mu$
(maybe $\varepsilon$-greedy, or human demonstrations, or a previous version of itself)
but wants to learn about a different \emph{target policy} $\pi$ (maybe the greedy one).
The catch is that any expectation we compute from $\mu$-collected data is, by default,
an expectation under $\mu$ --- not the $\pi$ we care about.

\textbf{Importance sampling} corrects this distribution mismatch by reweighting each
sample according to how likely it would be under the target policy:
\[
  \rho_t \;=\; \frac{\pi(A_t \mid S_t)}{\mu(A_t \mid S_t)}.
\]
A transition where $\pi$ would have taken this action much more often than $\mu$
gets a large weight; one where $\pi$ would rarely have chosen this action gets a small
weight.

\begin{intbox}[title={Reading the importance ratio}]
$\rho_t > 1$: ``the target would have done this more often than the behavior --- count
this transition for more.''

$\rho_t < 1$: ``the target would have done this less often --- count it for less.''

$\rho_t = 0$: ``the target would never have done this --- ignore the transition.''
\end{intbox}

Different methods need different amounts of correction.
\textbf{Off-policy MC} reweights the entire return by the product of ratios over the
whole trajectory, $\rho_t \rho_{t+1} \cdots \rho_{T-1}$ --- one mismatch anywhere in
the future contaminates the credit at $t$.
\textbf{Off-policy SARSA} only needs a single-step ratio $\rho_{t+1}$, because the
TD target only looks one step ahead.
\textbf{Q-learning} doesn't need any importance correction: its $\max_{a'}$ target
ignores which action the behavior policy actually took, so there is no mismatch to
correct.
This is one reason Q-learning is the standard off-policy method --- it avoids the
high-variance ratio products that haunt MC-style off-policy learning. We will see
the price of those products again in Chapter~5 when we look at off-policy policy
gradients.


% ==========================================================
%  CHAPTER 3 — VALUE FUNCTION APPROXIMATION
% ==========================================================
\chapter{Value Function Approximation}

Every tabular method we have seen so far stores a separate number for each
state or state-action pair.
That is fine for toy problems, but it breaks at scale.
Backgammon has around $10^{20}$ states; Go has around $10^{170}$; robotic
manipulation lives in a continuous space with infinitely many states.
No table can hold all of those, and no algorithm could visit each state enough
times to learn a reliable estimate.

The solution is to give up on exact representation and \emph{approximate}
the value function with a parameterized function $v_\mathbf{w}(s) \approx \vpi(s)$,
where $\mathbf{w}$ is a weight vector with far fewer entries than there are states.
The approximation generalizes: updating $\mathbf{w}$ from one state implicitly
updates the estimated values of nearby states that share features.
When neural networks play the role of this approximator, the whole framework
becomes \textbf{deep reinforcement learning}.

\section{Learning Value Functions by Gradient Descent}

The natural objective is to minimize the mean-squared difference between the
approximate and the true value function over states visited under $\pi$:
\[
  J(\mathbf{w}) = \E_\pi\!\left[\bigl(\vpi(S) - v_\mathbf{w}(S)\bigr)^2\right].
\]
Stochastic gradient descent gives the parameter update at each visited state $S$:
\[
  \Delta\mathbf{w} = \alpha\bigl(\vpi(S) - v_\mathbf{w}(S)\bigr)\nabla_\mathbf{w} v_\mathbf{w}(S).
\]
The problem is that $\vpi(S)$ is unknown --- we can never query the oracle.
So we substitute it with a \textbf{target} derived from observed data.
The MC target is the actual return $G_t$ (unbiased but noisy).
The TD(0) target is the one-step bootstrapped estimate $R_{t+1} + \gamma v_\mathbf{w}(S_{t+1})$
(lower variance, but biased because it depends on $\mathbf{w}$ itself).
The TD($\lambda$) target $G_t^\lambda$ sits between the two extremes.

This dependency is important: the TD target changes as $\mathbf{w}$ updates, so
we are not performing true gradient descent --- we are ignoring the gradient
through the target.
This is called a \textbf{semi-gradient} method.
It works well empirically, but convergence guarantees are weaker than for full
gradient descent, and there are regimes where it can diverge
(see the Deadly Triad, below).

\section{Function Approximators}

\textbf{Linear approximation} represents each state by a feature vector
$\mathbf{x}(s) \in \mathbb{R}^d$ and computes the value as a dot product:
$v_\mathbf{w}(s) = \mathbf{x}(s)^\top \mathbf{w}$.
The squared-error objective is then quadratic in $\mathbf{w}$, guaranteeing a
unique global minimum, and the SGD update simplifies to
$\Delta\mathbf{w} = \alpha\bigl(\text{target} - v_\mathbf{w}(S)\bigr)\mathbf{x}(S)$.
Tabular lookup is actually a special case of linear approximation: represent each
state with a one-hot vector, and the weights $\mathbf{w}$ directly encode the value
table.
The difference is that one-hot features provide no generalization between states,
whereas a shared feature representation does.

\begin{intbox}[title={Challenges Specific to RL}]
Unlike supervised learning, function approximation in RL faces three
compounding difficulties:
\textbf{non-i.i.d.\ data} (consecutive transitions are strongly correlated,
not independent samples);
\textbf{policy-dependent data distribution} (the states the agent visits
change as the policy improves); and
\textbf{non-stationary targets} (in TD methods, the regression target changes
every time $\mathbf{w}$ is updated, because the target network depends on the
same weights being trained).
Together these can destabilize learning, especially with nonlinear approximators.
\end{intbox}

In practice, the choice of approximator involves a three-way trade-off.
\emph{Tabular} methods have the best theory but cannot scale.
\emph{Linear} methods have reasonably strong theory but require hand-crafted
features.
\emph{Nonlinear} methods --- especially deep neural networks --- lack tight
theoretical guarantees but scale well and eliminate the need to design features
by hand.
Neural networks are the dominant choice in modern deep RL.

\section{Action-Value Approximation and Control}

For control we approximate $\qpi(s,a)$ rather than $\vpi(s)$, because
the greedy improvement step $\pi'(s) = \argmax_a q_\mathbf{w}(s,a)$ then
requires no model.
There are two natural network architectures.
\emph{Action-in} takes $(s, a)$ as input and produces a single Q-value ---
you need one forward pass per action to find the best one.
\emph{Action-out} takes $s$ as input and outputs a Q-value for every action
simultaneously --- much more efficient when the action space is discrete.

The same MC / TD(0) / TD($\lambda$) targets used for state values apply
directly to action values:
the oracle $\qpi(S,A)$ is replaced by $G_t$, by
$R_{t+1} + \gamma q_\mathbf{w}(S_{t+1},A_{t+1})$, or by $G_t^\lambda$.
The control loop then mirrors GPI: approximate policy evaluation
followed by $\varepsilon$-greedy improvement.

\begin{notebox}[title={The Deadly Triad}]
Combining \textbf{function approximation}, \textbf{bootstrapping} (TD-style
targets), and \textbf{off-policy learning} in the same algorithm can cause
divergence.
Each ingredient alone is fine; together they can form a feedback loop where
weights drift to infinity.
MC avoids it by not bootstrapping; on-policy TD with linear approximation
avoids it by restricting the data distribution.
Non-linear, off-policy TD --- the core of DQN-style algorithms --- sits right
inside the danger zone, which is why DQN needs special stabilization tricks.
\end{notebox}

\section{Batch Learning and Experience Replay}

Plain online SGD is sample-inefficient: each transition is seen once and then
discarded.
Worse, consecutive transitions are highly correlated, which makes the gradient
estimates noisy and can cause oscillations.

\textbf{Experience replay} addresses both problems.
The agent stores all transitions $(s_t,a_t,r_{t+1},s_{t+1})$ in a
\textbf{replay buffer} $\mathcal{D}$.
At each training step, a random mini-batch is sampled from $\mathcal{D}$ and
used to compute a gradient update:
\[
  \Delta\mathbf{w} = \alpha\bigl(r + \gamma v_\mathbf{w}(s') - v_\mathbf{w}(s)\bigr)\nabla_\mathbf{w} v_\mathbf{w}(s).
\]
Random sampling breaks the temporal correlations.
Each transition can be replayed many times, improving sample efficiency.
The buffer also acts as a stabilizing data distribution: because it contains
transitions from many past policies, the training distribution changes slowly
even as the current policy evolves quickly.

\section{Deep Q-Networks}

Applying Q-learning directly with a neural network runs into trouble.
Consecutive transitions are correlated (SGD assumes i.i.d.\ data).
The TD target $r + \gamma\max_{a'}Q_\mathbf{w}(s',a')$ depends on $\mathbf{w}$,
so every weight update shifts the target --- like trying to hit a moving goalpost.
And the combination of off-policy data, function approximation, and bootstrapping
puts us squarely in the deadly triad.

DQN (Mnih et al., 2015) tames all three with two techniques.
\textbf{Experience replay} (described above) breaks correlations and enables
data reuse.
\textbf{Fixed target networks} address the moving-goalpost problem: maintain a
second, frozen network $Q_{\mathbf{w}^-}$ and compute targets from it,
$y = r + \gamma\max_{a'}Q_{\mathbf{w}^-}(s',a')$,
only periodically copying $\mathbf{w}$ into $\mathbf{w}^-$.
Between copies the targets are stable, making training resemble supervised learning.

\begin{thmbox}[title={DQN Training Loop}]
For each step: collect $(s_t,a_t,r_{t+1},s_{t+1})$ with $\varepsilon$-greedy
and store in $\mathcal{D}$.
Sample a mini-batch $\{(s_j,a_j,r_j,s_j')\}_{j=1}^B$.
Compute targets $y_j = r_j + \gamma\max_{a'}Q_{\mathbf{w}^-}(s_j',a')$
(or $y_j = r_j$ if terminal).
Minimize $L(\mathbf{w}) = \frac{1}{B}\sum_j (y_j - Q_\mathbf{w}(s_j,a_j))^2$
w.r.t.\ $\mathbf{w}$ only (target $y_j$ treated as a constant).
Every $C$ steps: $\mathbf{w}^- \leftarrow \tau\,\mathbf{w} + (1-\tau)\,\mathbf{w}^-$.
\end{thmbox}

DQN was evaluated on 49 Atari 2600 games with \emph{the same architecture and
hyperparameters} across all games.
Raw pixels were preprocessed to $84\times84$ grayscale; four frames were stacked
to give the network a sense of motion; a CNN processed these into Q-values for
each action.
DQN outperformed all prior RL methods on 43 of 49 games and reached
human-level performance on 29 of them --- the first demonstration of
general-purpose deep RL at scale.

\section{Beyond DQN: Key Extensions}

DQN was a breakthrough but not the end of the story --- six targeted improvements
plug six different weaknesses, and \textbf{Rainbow} (Hessel et al., 2018) combines all
of them into one agent that reaches DQN's final performance in roughly $7\times$ fewer
environment steps.
The interesting thing about Rainbow is that no single ingredient is responsible for
the gain --- ablating any one of them noticeably hurts. Each extension fixes a
different bottleneck.

\paragraph{Double DQN --- fixing maximization bias.}
DQN's target $y = r + \gamma \max_{a'} Q_{\mathbf{w}^-}(s', a')$ uses the \emph{same}
frozen network to pick the best action and to evaluate its value. Whichever action
got an inflated estimate from noise gets selected, and that inflated estimate is
exactly the value used in the target. Same maximization bias as tabular Q-learning,
just with a neural network making the noise.
Double DQN decouples the two roles: the \emph{online} network selects the action,
$a^* = \argmax_{a'} Q_\mathbf{w}(s', a')$, and the \emph{target} network evaluates it,
$y = r + \gamma\, Q_{\mathbf{w}^-}(s', a^*)$.
The two networks have correlated noise but not identical noise, so the worst of the
upward bias is removed --- and DQN's notorious early-training overestimation largely
disappears.

\paragraph{Prioritized experience replay --- mining the buffer for surprises.}
Uniform replay treats every transition the same. But some transitions are
\emph{boring} --- the network already predicts them perfectly --- and some are
\emph{informative} --- the network still gets them wrong.
Sampling everything uniformly wastes most of the gradient budget on transitions the
network already understands.
Prioritized replay samples each transition $i$ with probability proportional to
$|\delta_i|^\alpha / \sum_k |\delta_k|^\alpha$, where $\delta_i$ is the TD error
(``how surprised was the network?''). Surprises get replayed more often. Because this
biases the data distribution, an importance-sampling weight $w_i \propto (1/p_i)^\beta$
corrects the resulting bias in the gradient. Result: faster learning, especially in
sparse-reward Atari games where successful trajectories are rare.

\paragraph{Dueling networks --- separating ``where am I'' from ``what should I do''.}
A Q-value bundles two questions into one number: how good is this state, and how
much better is this action than typical?
Sometimes the first question dominates (you're winning, every reasonable action
keeps the lead) and sometimes the second does (the position is neutral, but one
specific move wins). Mixing both into a single Q-head means the network spends
capacity learning the same ``state value'' redundantly across all actions.
A dueling architecture has a shared backbone that splits into two heads --- one for
the state value $V(s)$ and one for the per-action advantage $A(s,a)$ --- and
recombines them as $Q(s,a) = V(s) + A(s,a) - \tfrac{1}{|\mathcal{A}|}\sum_{a'} A(s,a')$.
The mean-subtraction term enforces an identifiability constraint and avoids the
two streams drifting in lockstep.
The big practical win: the value head can still be trained from \emph{every}
transition even when a particular action wasn't taken --- so generalization across
states improves, especially when many actions have similar value.

\paragraph{Multi-step returns --- propagating reward faster.}
1-step TD ($n = 1$) is conservative: a reward at time $t$ takes one update to reach
$t-1$, two to reach $t-2$, and so on. In sparse-reward games (where reward might be
hundreds of steps away from the action that caused it) this is glacial.
$n$-step returns $y_t^{(n)} = \sum_{k=0}^{n-1} \gamma^k R_{t+k+1} + \gamma^n \max_{a'} Q_{\mathbf{w}^-}(S_{t+n}, a')$
collect $n$ actual rewards before bootstrapping, propagating credit $n$ steps in one
go. The trade-off is variance (more random rewards summed in) and the fact that the
target becomes slightly off-policy (the $n$-step rollout was generated by the old
policy). Intermediate values like $n = 3$ to $5$ are usually a sweet spot.

\paragraph{Distributional RL --- modeling the whole distribution, not just the mean.}
Standard Q-learning learns $\E[G_t \mid s, a]$ --- a single scalar.
But two actions can have the same expected return while one is consistently mediocre
and the other is a high-variance gamble between disaster and triumph. The mean throws
that away.
Distributional RL learns the entire \emph{distribution} of returns $Z(s, a)$, satisfying
the distributional Bellman equation $Z(s,a) \overset{d}{=} R + \gamma Z(s', a')$.
The agent ends up modeling outcomes as probability distributions over a set of
``atoms'' (e.g.\ $51$ discrete return levels, hence the original method's name C51).
Empirically, training with a cross-entropy loss on these distributions produces much
richer representations --- and improves performance even though we still extract a
deterministic policy by taking the mean's argmax.

\paragraph{Noisy Networks --- learning to explore.}
$\varepsilon$-greedy is uniform: the same noise level everywhere, regardless of
whether the agent is confident or confused. That is a bad use of exploration budget.
Noisy Networks replace deterministic linear layers with stochastic ones:
\[
  y = (\mu^w + \sigma^w \odot \varepsilon^w)\, x + (\mu^b + \sigma^b \odot \varepsilon^b),
  \quad \varepsilon \sim \Normal(0, I).
\]
The noise scales $\sigma^w, \sigma^b$ are \emph{learned} alongside the means. Where
the network is confident it shrinks $\sigma$ toward zero (less exploration); where it
is uncertain it keeps $\sigma$ large (more exploration).
This produces \emph{state-dependent} exploration without a hand-tuned $\varepsilon$
schedule --- and on hard-exploration Atari games it makes a real difference.

\begin{intbox}[title={Why six fixes, not one?}]
Each Rainbow ingredient targets a separate failure mode of vanilla DQN: maximization
bias (Double DQN), inefficient sampling (prioritized replay), wasted capacity in the
value head (dueling), slow credit propagation (multi-step), information loss in the
mean (distributional), and dumb exploration (noisy nets).
None of these is a magic bullet by itself. Together they cover most of the obvious
ways DQN leaks performance, which is why Rainbow's gain is greater than the sum of
the individual contributions.
\end{intbox}


% ==========================================================
%  CHAPTER 4 — POLICY-BASED METHODS
% ==========================================================
\chapter{Policy-Based Methods}

So far every method in this course has learned a value function and then derived a
policy from it --- usually $\varepsilon$-greedy w.r.t.\ a learned $q(s,a)$.
This chapter takes the opposite route: we \emph{directly parameterize the policy
itself} and learn its parameters by gradient ascent on the expected return.
A value function may still appear (as a baseline, or as a critic) but it is no
longer the primary object of learning.

The motivation for this shift is concrete: there are tasks where value-based RL is a
poor fit, and policy-based methods handle them naturally.

\section{Why Bother With Policies Directly?}

Recall the three families of RL methods:
\begin{itemize}
  \item \textbf{Value-based:} learn $q(s,a)$, derive $\pi(s) = \argmax_a q(s,a)$.
  \item \textbf{Policy-based:} learn $\pi_\theta(a \mid s)$ directly --- no value function.
  \item \textbf{Actor-critic:} learn both --- a parameterized policy (the actor) and a
    value function (the critic) that helps train the actor.
\end{itemize}

Value-based methods are great when actions are discrete and the optimal policy is
essentially deterministic. They run into trouble in two specific situations.

\subsection{Continuous action spaces}

Many real control problems have \emph{continuous} actions: a robot arm's joint
torques live in $\R^d$, a car's steering and throttle are real numbers, locomotion
involves dozens of simultaneous motor commands.
Deriving a greedy policy in these settings means solving
\[
  a^*(s) \;=\; \argmax_{a \in \R^d}\, q(s,a)
\]
at every step.
For a handful of discrete actions this is trivial --- compute $q(s,a)$ for each one
and pick the largest.
In a continuous space it is a global optimization problem inside the inner loop of
training.
There is no easy way to do it.

A parameterized policy sidesteps the issue entirely: instead of finding the argmax
of a learned function, we just \emph{output} the action (or a distribution over
actions) directly from a network.

\subsection{Deterministic policies are not always enough}

Value-based control with greedy improvement converges to a \emph{deterministic}
policy --- one action per state.
For most problems, that is fine. For some, it is fatal.

\begin{exbox}[title={Rock--Paper--Scissors}]
The optimal policy in rock-paper-scissors against a savvy opponent is to play each
move with probability $\tfrac{1}{3}$ --- the uniform Nash equilibrium.
Any deterministic policy is trivially exploitable: play rock every time and your
opponent will eventually figure out to play paper.
Value-based RL cannot represent the optimal policy here, because greedy improvement
always collapses to a single action.
\end{exbox}

\begin{exbox}[title={Aliased Gridworld}]
A 1D gridworld where the agent's state representation is the binary feature vector
\[
  \phi(s) = [\text{wall up},\, \text{wall right},\, \text{wall down},\, \text{wall left}].
\]
The treasure sits in the middle, between two skull squares.
The two ``corridor'' states on either side of the treasure look identical to the
agent: same walls in the same directions. They are \emph{aliased}.

Any deterministic policy on this aliased feature must take the same action in both
corridor states. If it picks ``move left,'' the agent gets stuck on the left side;
``move right'' and it gets stuck on the right. Either way the agent never reliably
finds the treasure.

A \emph{stochastic} policy that flips a coin in those aliased states ($\pi(\text{left}) =
\pi(\text{right}) = 0.5$) reaches the treasure quickly with high probability,
because the random walk is short and bounded. Direct policy learning can represent
exactly this kind of policy.
\end{exbox}

\subsection{Interlude: POMDPs}

The aliased gridworld is a small instance of a much bigger phenomenon.
Up to now we have assumed the agent observes the true state $S_t$ --- a
\textbf{fully observable MDP}.
In practice the agent often sees only a partial \textbf{observation}
$O_t \sim p(o \mid S_t)$, with the actual state $S_t$ hidden.
This is a \textbf{Partially Observable MDP (POMDP)}.

Once observations no longer contain the full state, the optimal action may depend
on history (memory) or randomization (a stochastic policy).
This course mostly ignores memory --- but allowing stochastic policies is essentially
free with policy-based methods, and it solves a real subset of POMDP problems.

\begin{intbox}
The two reasons to want a stochastic policy are subtly different:
\begin{itemize}
  \item \textbf{Game-theoretic:} the optimal policy is genuinely mixed (rock-paper-scissors).
  \item \textbf{Aliasing:} our state representation is incomplete, so randomization is
    a cheap proxy for the missing information.
\end{itemize}
Both are out of reach for value-based control. Both are easy for policy-based methods.
\end{intbox}

\subsection{Trade-offs}

Policy-based methods are not strictly better than value-based ones --- they buy
flexibility at a cost.

\paragraph{Advantages.}
They optimize the true objective (expected return) directly, with no value-function
intermediary.
They handle continuous actions naturally.
They can represent stochastic policies.
Sometimes the policy is much simpler than the value function (think ``always go left''
versus the maze of values that justifies it), and learning the simpler thing is
easier.
The choice of policy class can also encode prior knowledge about the problem's structure.

\paragraph{Disadvantages.}
They typically have higher \emph{variance} than value-based methods --- which we'll
see directly when we look at the policy gradient.
They can be less sample-efficient, get stuck in local optima, and the knowledge they
learn tends to be specific to the policy class chosen --- it generalizes less than a
value function would.

\section{Setting Up the Optimization}

A parameterized policy $\pi_\theta(a \mid s)$ is just a function from states to action
distributions, with parameters $\theta$ (e.g.\ the weights of a neural network).
The job is to find the parameters that maximize policy quality --- but first we need
to define what ``quality'' means.
There are two natural choices, depending on whether the task is episodic or continuing.

\subsection{Episodic objective: average return per episode}

For episodic tasks we measure the policy by the expected return from the
starting-state distribution $d_0$:
\[
  J(\theta) \;=\; \E_{S_0 \sim d_0,\,\pi_\theta}\!\left[ \sum_{t=0}^{T} \gamma^t R_t \right]
            \;=\; \E_{S_0 \sim d_0}\!\bigl[ v_{\pi_\theta}(S_0) \bigr].
\]
In words: \emph{the value of the start state, averaged over where the start state can be}.

\subsection{Continuing objective: average reward per step}

For continuing tasks the discounted episodic objective doesn't quite work --- the sum
runs to infinity and the boundary effects of $S_0$ vanish.
We use the average reward per step instead:
\[
  J(\theta) \;=\; \E_{\pi_\theta}[R_{t+1}]
  \;=\; \sum_s d_{\pi_\theta}(s) \sum_a \pi_\theta(a \mid s) \sum_r p(r \mid s, a)\, r,
\]
where $d_{\pi_\theta}(s)$ is the \emph{stationary distribution} over states under
$\pi_\theta$ --- the long-run fraction of time the agent spends in $s$.
This is the expected reward in the states the policy visits.

For the rest of this chapter, treat both objectives as ``the expected return $J(\theta)$.''
The mathematical machinery that follows applies to both.

\subsection{Optimization choices}

Once we have $J(\theta)$, finding $\theta^* = \argmax_\theta J(\theta)$ is just
optimization. There are two broad classes of methods.
\begin{itemize}
  \item \textbf{Gradient-free:} hill climbing, simplex / Nelder-Mead, genetic algorithms.
    Useful when the policy isn't differentiable, but not competitive in the
    high-dimensional regime where neural-network policies live.
  \item \textbf{Gradient-based:} gradient ascent, conjugate gradient, quasi-Newton.
    These require a way to compute or estimate $\nabla_\theta J(\theta)$ --- and
    that is where things get interesting.
\end{itemize}

The rest of the chapter is about the gradient route.

\section{Computing the Policy Gradient}

The vanilla update is
\[
  \theta \;\leftarrow\; \theta + \alpha\, \nabla_\theta J(\theta),
\]
identical to ordinary gradient ascent.
The challenge is that $J(\theta)$ is an \emph{expectation over trajectories}, and the
trajectories themselves depend on $\theta$. We can't write down a closed form ---
we need a way to estimate the gradient from samples.

Three approaches:
\begin{itemize}
  \item \textbf{Finite differences:} perturb each parameter individually and look at
    the change in $J$.
  \item \textbf{Score function / log-likelihood trick:} a clever algebraic identity
    that converts a gradient of an expectation into an expectation of a gradient.
  \item \textbf{Policy gradient theorem:} the score-function trick, applied to the
    full multi-step MDP setting.
\end{itemize}

\subsection{Finite differences}

The crudest approach is to estimate each partial derivative numerically by perturbing
the parameter:
\[
  \pd{J(\theta)}{\theta_i} \;\approx\; \frac{J(\theta + \varepsilon e_i) - J(\theta - \varepsilon e_i)}{2\varepsilon},
\]
where $e_i$ is the unit vector in the $i$-th direction.

This is simple and works for any policy --- even non-differentiable ones.
The bad news is that it requires $\bigO{n}$ rollouts of the entire policy for $n$
parameters.
For a neural net with millions of parameters this is hopeless, and the noise from
sampling makes the estimate unreliable.

\subsection{The score-function (log-likelihood) trick}

For differentiable policies there is a much better approach. Use the identity
\[
  \nabla_\theta \pi_\theta(a \mid s)
  \;=\; \pi_\theta(a \mid s)\, \frac{\nabla_\theta \pi_\theta(a \mid s)}{\pi_\theta(a \mid s)}
  \;=\; \pi_\theta(a \mid s)\, \nabla_\theta \log \pi_\theta(a \mid s).
\]
The gradient of the log-probability $\nabla_\theta \log \pi_\theta(a \mid s)$ is called
the \textbf{score function}.

\begin{intbox}
The score-function identity reads ``$\nabla \pi = \pi \cdot \nabla \log \pi$.'' Its
purpose is to convert a $\nabla$ that sits \emph{outside} an expectation into a
$\nabla \log \pi$ that sits \emph{inside} the expectation, where it can be sampled.
This is the same trick that powers REINFORCE, variational inference, and policy
gradient methods generally --- it is one of the most useful identities in modern
ML.
\end{intbox}

To see why this matters, consider a one-step MDP (a contextual bandit): the agent
sees $S_t$, picks $A_t$, gets reward $R_{t+1}$, and the episode ends.
The objective is $J(\theta) = \E_{\pi_\theta}[R_{t+1}]$.
Naively you might try ``sample $R$ and take the gradient,'' but $R$ depends on
$\theta$ only through the action the policy happened to choose --- once sampled, $R$
is just a number with no $\theta$ inside it.

Apply the trick to push the gradient inside the expectation:
\begin{align*}
  \nabla_\theta J(\theta)
  &= \nabla_\theta \sum_s d(s) \sum_a \pi_\theta(a \mid s)\, r_{s,a} \\
  &= \sum_s d(s) \sum_a r_{s,a}\, \nabla_\theta \pi_\theta(a \mid s) \\
  &= \sum_s d(s) \sum_a r_{s,a}\, \pi_\theta(a \mid s)\, \nabla_\theta \log \pi_\theta(a \mid s) \\
  &= \E_{\pi_\theta}\!\bigl[\, R_{t+1}\, \nabla_\theta \log \pi_\theta(A_t \mid S_t)\, \bigr].
\end{align*}
Now the gradient is an expectation we can sample.
The stochastic update is
\[
  \theta \;\leftarrow\; \theta + \alpha\, R_{t+1}\, \nabla_\theta \log \pi_\theta(A_t \mid S_t).
\]
This is the \textbf{REINFORCE update for one-step MDPs}.

\subsection{Example: softmax policy}

The standard discrete-action policy class is the softmax:
\[
  \pi_\theta(a \mid s) \;=\; \frac{e^{h_\theta(s,a)}}{\sum_b e^{h_\theta(s,b)}},
\]
where $h_\theta(s,a)$ is a learnable preference (logit) for action $a$ in state $s$.
The score function is
\[
  \nabla_\theta \log \pi_\theta(A_t \mid S_t)
  \;=\; \underbrace{\nabla_\theta h_\theta(S_t, A_t)}_{\text{preference of chosen action}}
       \;-\; \underbrace{\sum_a \pi_\theta(a \mid S_t)\, \nabla_\theta h_\theta(S_t, a)}_{\text{expected preference under policy}}.
\]
Compare this to the gradient of softmax cross-entropy in supervised learning ---
exactly the same shape: ``increase the preference of the action you took, decrease
the average preference.''
The reward-weighted version then says: ``do this more strongly when the reward was high.''

In the linear case $h_\theta(s,a) = \theta\T \phi(s,a)$ the score function simplifies to
\[
  \nabla_\theta \log \pi_\theta(A_t \mid S_t)
  \;=\; \phi(S_t, A_t) - \E_{\pi_\theta}\!\bigl[\phi(S_t,\cdot)\bigr]
  \;=\; \phi(S_t, A_t) - \sum_a \pi_\theta(a \mid S_t)\, \phi(S_t, a),
\]
i.e.\ the difference between the chosen action's features and the policy's average
features.

\section{The Policy Gradient Theorem}

The score-function trick gives us the gradient for one-step problems.
The full multi-step case looks harder: an action $A_t$ affects all future rewards
through state transitions, so naively you'd need to differentiate through the
environment dynamics --- which the agent does not know.

The remarkable result of the \textbf{policy gradient theorem} is that you don't have to.
The gradient takes essentially the same shape as the bandit case, with $R_{t+1}$
replaced by the action-value $q_{\pi_\theta}(S_t, A_t)$.

\begin{thmbox}[title={Policy Gradient Theorem (Episodic)}]
For any differentiable policy $\pi_\theta(a \mid s)$ and any starting-state
distribution $d_0$, the gradient of $J(\theta) = \E_{d_0,\pi_\theta}[G_0]$ is
\[
  \nabla_\theta J(\theta)
  \;=\; \E_{d_0,\, \pi_\theta}\!\left[\, \sum_{t=0}^{T}
            \gamma^t\, q_{\pi_\theta}(S_t, A_t)\, \nabla_\theta \log \pi_\theta(A_t \mid S_t)\, \right].
\]
The gradient does \emph{not} require differentiating through the MDP dynamics.
\end{thmbox}

The proof is short enough to spell out --- and the structure is worth seeing once, because
the same pattern reappears every time you derive a policy gradient.

\paragraph{Step 1: write $J$ as an expectation over trajectories.}
Let $\tau = (S_0, A_0, R_1, S_1, A_1, \dots)$ and $G(\tau) = \sum_t \gamma^t R_{t+1}$.
Then
\[
  \nabla_\theta J(\theta)
  = \nabla_\theta \E_{\pi_\theta}[G(\tau)]
  = \nabla_\theta \sum_\tau p_{\pi_\theta}(\tau)\, G(\tau).
\]

\paragraph{Step 2: apply the score-function trick.}
Push $\nabla_\theta$ inside via $\nabla p = p \nabla \log p$:
\[
  \nabla_\theta J(\theta)
  = \sum_\tau p_{\pi_\theta}(\tau)\, G(\tau)\, \nabla_\theta \log p_{\pi_\theta}(\tau)
  = \E_{\pi_\theta}\!\bigl[\, G(\tau)\, \nabla_\theta \log p_{\pi_\theta}(\tau)\, \bigr].
\]

\paragraph{Step 3: expand $\log p_{\pi_\theta}(\tau)$.}
The trajectory probability factorizes as
\[
  p_{\pi_\theta}(\tau)
  = p(S_0) \prod_{t=0}^{T} \underbrace{p(S_{t+1} \mid S_t, A_t)}_{\text{environment dynamics}}\;
                          \underbrace{\pi_\theta(A_t \mid S_t)}_{\text{policy}}.
\]
Take the log and differentiate w.r.t.\ $\theta$.
The initial-state distribution $p(S_0)$ and the transition probabilities
$p(S_{t+1} \mid S_t, A_t)$ \textbf{do not depend on $\theta$}, so their gradients
vanish:
\[
  \nabla_\theta \log p_{\pi_\theta}(\tau)
  \;=\; \sum_{t=0}^{T} \nabla_\theta \log \pi_\theta(A_t \mid S_t).
\]

\begin{intbox}[title={The key step}]
This is the moment the unknown environment dynamics drops out of the gradient.
We never need a model of $p(s' \mid s, a)$. The gradient of the policy's
log-probability is enough --- and we have that, because we wrote down the policy
ourselves.
\end{intbox}

\paragraph{Step 4: causality / clean-up.}
Substituting back gives
\[
  \nabla_\theta J(\theta)
  = \E_{\pi_\theta}\!\left[\, G(\tau) \sum_t \nabla_\theta \log \pi_\theta(A_t \mid S_t)\, \right].
\]
A causality argument tightens this further: rewards $R_{k+1}$ with $k < t$ cannot
depend on $A_t$ (the action hasn't happened yet), so in expectation those terms
contribute zero.
After splitting $G(\tau) = \sum_k \gamma^k R_{k+1}$ and dropping the zero-mean
$k < t$ terms, the remaining reward-tail at step $t$ is the discounted return from
$t$ onwards --- which in expectation equals $\gamma^t q_{\pi_\theta}(S_t, A_t)$.
That gives the theorem. $\blacksquare$

\begin{intbox}[title={What the policy gradient is doing}]
At every visited $(S_t, A_t)$, the gradient says: take the score function
$\nabla \log \pi_\theta(A_t \mid S_t)$ (which points in the parameter direction
that increases the probability of $A_t$ in $S_t$), and weight it by how good the
action turned out to be ($q_{\pi_\theta}(S_t, A_t)$).
Good actions get their probability pushed up; bad actions get pushed down.
The amazing part is that we never had to know how the environment works.
\end{intbox}

\subsection{From theorem to algorithm: REINFORCE}

The theorem tells us \emph{what} the gradient is, but $q_{\pi_\theta}(S_t, A_t)$ is
not directly observable.
The simplest unbiased estimate is the actually observed return $G_t$ from a sampled
episode:
\[
  q_{\pi_\theta}(S_t, A_t) = \E_{\pi_\theta}[G_t \mid S_t, A_t] \;\approx\; G_t.
\]
Substituting gives the \textbf{REINFORCE} algorithm (Williams, 1992):
\[
  \theta \;\leftarrow\; \theta + \alpha\, \gamma^t\, G_t\, \nabla_\theta \log \pi_\theta(A_t \mid S_t).
\]

\begin{thmbox}[title={REINFORCE (Monte-Carlo Policy Gradient)}]
For each episode $S_0, A_0, R_1, \dots, S_T$ generated by $\pi_\theta$:
\begin{enumerate}
  \item Compute returns $G_t = R_{t+1} + \gamma R_{t+2} + \cdots + \gamma^{T-t-1} R_T$
    for every $t$.
  \item For each $t$, update
  $\theta \leftarrow \theta + \alpha\, \gamma^t\, G_t\, \nabla_\theta \log \pi_\theta(A_t \mid S_t).$
\end{enumerate}
\emph{Monte-Carlo} because we use the actual sampled return as the target.
\end{thmbox}

REINFORCE is unbiased --- in expectation it tracks the true policy gradient ---
but it is \textbf{noisy}.
$G_t$ is a sum of many random rewards, so its variance grows with the episode length,
and the resulting gradient estimates can swing wildly between updates.

\begin{exbox}[title={CartPole as a sanity check}]
The textbook policy-gradient toy: balance a pole on a moving cart by pushing left or
right. State: $S_t = (\text{position}, \text{speed}, \text{angle}, \text{angle-speed})$.
Action: $A_t \in \{\text{left}, \text{right}\}$. Reward: $+1$ for every step the pole
stays up; episode ends if it falls or after $200$ steps.
Policy: a two-layer net with $32$ hidden tanh units and a softmax over the two actions.

REINFORCE learns to balance, but the training curve is a noisy zigzag --- typical
high-variance behavior. Adding a baseline (next section) smooths it out dramatically.
\end{exbox}

\section{Reducing Variance: Baselines}

REINFORCE works, but its variance is brutal.
The good news: we can reduce variance dramatically without introducing any bias.

\subsection{The baseline identity}

Pick any function $b(s)$ that depends only on the state, not on the action. Then
\begin{align*}
  \E_{\pi_\theta}\!\bigl[\, b(S_t)\, \nabla_\theta \log \pi_\theta(A_t \mid S_t)\, \bigr]
  &= \E\!\left[\, b(S_t) \sum_a \pi_\theta(a \mid S_t)\, \nabla_\theta \log \pi_\theta(a \mid S_t)\, \right] \\
  &= \E\!\left[\, b(S_t) \sum_a \nabla_\theta \pi_\theta(a \mid S_t)\, \right] \\
  &= \E\!\left[\, b(S_t)\, \nabla_\theta \underbrace{\sum_a \pi_\theta(a \mid S_t)}_{=\, 1}\, \right]
  \;=\; 0.
\end{align*}
The expected ``baseline term'' is exactly zero, regardless of what $b$ is.
So we can subtract $b(S_t)$ from $G_t$ inside the REINFORCE update without changing the
gradient in expectation:
\[
  \theta \;\leftarrow\; \theta + \alpha\, \gamma^t\,\bigl(G_t - b(S_t)\bigr)\, \nabla_\theta \log \pi_\theta(A_t \mid S_t).
\]
The expected gradient is unchanged.
The variance, on the other hand, can be enormously reduced if we choose $b(S_t)$ to
track $\E[G_t \mid S_t]$.

\subsection{Why a state-value baseline?}

The natural choice is the state-value function: $b(s) = v_\mathbf{w}(s)$, learned
alongside the policy.
The advantage interpretation makes this concrete:
\[
  G_t - v_\mathbf{w}(S_t) \;\approx\; q_{\pi_\theta}(S_t, A_t) - v_{\pi_\theta}(S_t)
                                 \;=\; A_{\pi_\theta}(S_t, A_t),
\]
the \textbf{advantage} of action $A_t$ over the policy's average behavior in $S_t$.

\begin{intbox}
A baseline doesn't change \emph{which way} the gradient pushes; it changes \emph{how
emphatically}.
If every action in some state earns reward around $+100$, plain REINFORCE strongly
increases the probability of \emph{all} of them --- but only one of them was actually
above average.
Subtracting the baseline cancels the common offset so only genuine differences
in action quality move the policy.
The right question shifts from ``was the cumulative reward high or low?'' to
``was this action better or worse than typical?''
On CartPole the practical effect is a much cleaner training curve.
\end{intbox}

In REINFORCE-with-baseline the value function is itself learned by regression on
the same returns:
\[
  \mathbf{w} \;\leftarrow\; \mathbf{w} + \alpha_\mathbf{w}\bigl(G_t - v_\mathbf{w}(S_t)\bigr)\nabla_\mathbf{w} v_\mathbf{w}(S_t).
\]
But it is \emph{only} a baseline --- it does not provide the learning signal for the
policy.
The target for $\theta$ is still the Monte-Carlo return $G_t$.

\section{Actor-Critic Methods}

REINFORCE-with-baseline still uses the full Monte-Carlo return as its target.
This is unbiased but high-variance, and it requires waiting for the episode to
terminate before any update is possible.
The natural next step is to \emph{bootstrap} the return --- exactly what TD does for
value-based methods.

\subsection{From baseline to critic}

Replace $G_t$ with the one-step bootstrapped estimate
$G_{t:t+1} = R_{t+1} + \gamma v_\mathbf{w}(S_{t+1})$ in the policy update:
\begin{align*}
  \theta &\leftarrow \theta + \alpha_\theta\, \gamma^t\bigl(G_{t:t+1} - v_\mathbf{w}(S_t)\bigr)\, \nabla_\theta \log \pi_\theta(A_t \mid S_t) \\
         &= \theta + \alpha_\theta\, \gamma^t\bigl(R_{t+1} + \gamma v_\mathbf{w}(S_{t+1}) - v_\mathbf{w}(S_t)\bigr)\, \nabla_\theta \log \pi_\theta(A_t \mid S_t) \\
         &= \theta + \alpha_\theta\, \gamma^t\, \delta_t\, \nabla_\theta \log \pi_\theta(A_t \mid S_t).
\end{align*}
The quantity in parentheses is just the TD error
$\delta_t = R_{t+1} + \gamma v_\mathbf{w}(S_{t+1}) - v_\mathbf{w}(S_t)$ --- which doubles
as a one-step advantage estimate.
The value function is now doing two jobs at once: providing the baseline \emph{and}
bootstrapping the return.

\begin{defbox}[title={Actor-Critic (One-Step)}]
Two parameter sets are learned simultaneously:
\begin{itemize}
  \item the \textbf{actor} $\pi_\theta(a \mid s)$ --- the parameterized policy.
  \item the \textbf{critic} $v_\mathbf{w}(s)$ --- a value function used to estimate
    advantages and to provide the policy's learning signal.
\end{itemize}
At each step compute
$\delta_t = R_{t+1} + \gamma v_\mathbf{w}(S_{t+1}) - v_\mathbf{w}(S_t)$, then update
\[
  \mathbf{w} \,\leftarrow\, \mathbf{w} + \alpha_\mathbf{w}\, \delta_t\, \nabla_\mathbf{w} v_\mathbf{w}(S_t),
  \qquad
  \theta \,\leftarrow\, \theta + \alpha_\theta\, \gamma^t \delta_t\, \nabla_\theta \log \pi_\theta(A_t \mid S_t).
\]
\end{defbox}

The big change relative to REINFORCE-with-baseline is that the critic now
\emph{provides} the learning signal, not just a variance reduction.
Updates can happen online, every step --- no need to wait for the episode to finish.

\begin{intbox}[title={Bias--variance trade-off again}]
\textbf{Pure Monte-Carlo} (REINFORCE) is unbiased but high-variance.
\textbf{One-step bootstrap} (actor-critic) is low-variance but biased --- $v_\mathbf{w}$
is only an approximation, and bootstrapping propagates its errors.
You get faster, smoother learning but pay with some asymptotic correctness.
This is exactly the same dial as $\lambda$ in TD($\lambda$) from the value-based
chapter; the policy-gradient world has its own version of it (GAE, below).
\end{intbox}

\subsection{n-step actor-critic}

The natural interpolation between one-step bootstrap and full Monte-Carlo is the
$n$-step return:
\[
  G_t^{(n)} \;=\; \sum_{l=0}^{n-1} \gamma^l\, R_{t+l+1} \;+\; \gamma^n v_\mathbf{w}(S_{t+n}).
\]
The corresponding $n$-step advantage is $\hat\delta_t^{(n)} = G_t^{(n)} - v_\mathbf{w}(S_t)$,
and the actor update uses it in place of $\delta_t$:
\[
  \theta \;\leftarrow\; \theta + \alpha_\theta\, \gamma^t\, \hat\delta_t^{(n)}\, \nabla_\theta \log \pi_\theta(A_t \mid S_t).
\]
Larger $n$ uses more of the actual reward signal (lower bias, higher variance);
smaller $n$ relies more on the critic (higher bias, lower variance).
Picking $n$ becomes a per-environment tuning problem.

\subsection{Generalized Advantage Estimation (GAE)}

Why pick a single $n$ at all? GAE (Schulman et al., 2015) averages over all of them
with an exponential weighting.

Let $\delta_t = R_{t+1} + \gamma v_\mathbf{w}(S_{t+1}) - v_\mathbf{w}(S_t)$ be the
per-step TD error. The GAE estimator is
\[
  \hat A_t^{\text{GAE}(\lambda)}
  \;=\; \sum_{l=0}^{\infty} (\gamma\lambda)^l\, \delta_{t+l}
  \;=\; \delta_t + \gamma\lambda\, \delta_{t+1} + (\gamma\lambda)^2\, \delta_{t+2} + \cdots.
\]
This is exactly an exponentially-weighted combination of the $n$-step advantages
$\hat\delta_t^{(n)}$ for all $n$, with $\lambda$ controlling how aggressively longer
horizons are downweighted.

\begin{intbox}
GAE is the policy-gradient analogue of TD($\lambda$): same exponential weighting,
same $\lambda$ knob.
\begin{itemize}
  \item $\lambda = 0$: pure one-step TD, $\hat A_t = \delta_t$. Maximally biased,
    minimally noisy.
  \item $\lambda = 1$: full Monte-Carlo advantage. Unbiased, maximally noisy.
  \item $0 < \lambda < 1$: smooth interpolation; in practice
    $\lambda \in [0.9, 0.97]$ is typical.
\end{itemize}
$\lambda$ and $\gamma$ play different roles --- $\gamma$ is the agent's actual
discount on future reward, $\lambda$ is purely an estimator-side bias/variance knob.
\end{intbox}

GAE will reappear in the next chapter as the standard advantage estimator for modern
actor-critic algorithms like A2C and PPO.

\section{Recap}

Policy-based methods learn the policy directly, by gradient ascent on the expected
return. Three big ideas make this work.
\begin{itemize}
  \item The \textbf{score-function trick} converts $\nabla \E[\cdot]$ into
    $\E[\nabla \log \pi \cdot \,\cdot\,]$, so the gradient becomes a sample-able
    expectation.
  \item The \textbf{policy gradient theorem} extends the trick to multi-step MDPs
    and shows that we never need to differentiate through the environment.
  \item \textbf{Baselines and critics} keep the gradient unbiased (or only mildly
    biased) while drastically reducing its variance --- without them REINFORCE is
    almost too noisy to be useful.
\end{itemize}
The next chapter pushes this further with modern actor-critic algorithms: deep,
batched, and engineered to actually be stable on hard problems.


% ==========================================================
%  CHAPTER 5 — MODERN ACTOR-CRITIC ALGORITHMS
% ==========================================================
\chapter{Modern Actor-Critic Algorithms}

Plain actor-critic with GAE is the algorithmic skeleton.
Making it actually work on hard problems --- continuous control, sparse rewards,
real-world robotics, language-model fine-tuning --- requires several more layers.
This chapter walks through the modern landscape: NPG, TRPO, PPO on the on-policy side;
DDPG, TD3, SAC on the off-policy / continuous-control side.

Three concrete weaknesses of vanilla actor-critic motivate everything below.
First, it is \textbf{on-policy} --- data collected with the current $\pi_\theta$ cannot
be reused after $\theta$ changes, which is wasteful.
Second, its updates are \textbf{unstable} --- small steps in parameter space can produce
huge changes in action probabilities, and one bad step contaminates all future data.
Third, even with baselines and GAE the gradient estimates remain \textbf{noisy}.
Modern algorithms attack each of these weaknesses with a different tool.

\section{A Unified View of Policy Gradients}

Every policy-gradient method we have seen --- and most we are about to see ---
shares the same structural form:
\[
  \nabla_\theta J(\theta) \;=\; \E\!\left[\, \sum_{t=0}^{T} \Psi_t\, \nabla_\theta \log \pi_\theta(A_t \mid S_t)\, \right],
\]
where $\Psi_t$ is the \emph{learning signal} for the action taken at time $t$.
Different algorithms differ only in the choice of $\Psi_t$:

\begin{center}
\begin{tabular}{ll}
\toprule
\textbf{Choice of $\Psi_t$} & \textbf{Algorithm} \\
\midrule
$G_t$ (Monte-Carlo return)                                         & REINFORCE \\
$G_t - v_\mathbf{w}(S_t)$                                          & REINFORCE with baseline \\
$\delta_t = R_{t+1} + \gamma v_\mathbf{w}(S_{t+1}) - v_\mathbf{w}(S_t)$ & TD(0) actor-critic / A2C \\
$G_t^{(n)} - v_\mathbf{w}(S_t)$                                    & $n$-step actor-critic \\
$\sum_{l=0}^\infty (\gamma\lambda)^l \delta_{t+l}$                 & GAE (used in PPO) \\
$q_\mathbf{w}(S_t, A_t)$                                           & Q-based actor-critic (DDPG, SAC) \\
\bottomrule
\end{tabular}
\end{center}

Reading top-to-bottom, this is a spectrum from \textbf{maximum-bias-zero / variance-high}
(REINFORCE) to \textbf{low-bias / low-variance} (actor-critic with a learned critic).
The job of each algorithm is to pick a point on this spectrum and pair it with a way
of stabilising the resulting updates.

\subsection{Bias and variance, again}

Every $\Psi_t$ sitting below the first row uses a learned value function and is
therefore \emph{biased} --- the critic is only an approximation, and bootstrapping its
errors propagate.

\begin{intbox}[title={Compatible function approximation (briefly)}]
There \emph{is} a way to pair the critic with the actor so that the policy-gradient
estimate stays unbiased: choose features for the critic that are ``compatible'' with
the policy's score function (Sutton et al., 2000). The compatible features are
$\nabla_\theta \log \pi_\theta(s,a)$ themselves, and the critic is restricted to a
linear function of those features.
This is mathematically beautiful but in deep RL almost nobody uses it: the linearity
constraint is way too restrictive for neural networks.
The pragmatic deep-RL stance is ``low variance plus an expressive critic beats unbiasedness.''
\end{intbox}

\subsection{Why $\gamma^t$ usually disappears}

Looking back at the policy gradient theorem, every term carried a $\gamma^t$ in front
of the score function. In the slides and in most practical code that $\gamma^t$ is
silently dropped. The reason is pragmatic: with $\gamma^t$, late-step gradients are
exponentially down-weighted, and most of the learning signal comes from the first few
steps of each episode --- which works poorly in practice.
Dropping it gives a (slightly) biased estimator of the discounted return objective but
a much better one for the average-reward objective, and it is what every modern
implementation actually does.
We follow that convention from now on.

\section{Off-Policy Policy Gradients}

The on-policy nature of REINFORCE / actor-critic is a real annoyance.
Every gradient step requires fresh trajectories from the current $\pi_\theta$, and as
soon as we update, all the previously collected data is technically off-policy.
Throwing away data after one update is wasteful --- can we keep using it?

\subsection{Importance sampling}

If data was collected by some \emph{behavior policy} $\mu(a \mid s)$ but we want
expectations under the \emph{target policy} $\pi(a \mid s)$, the standard fix is
importance sampling:
\[
  \E_{a \sim \pi}[f(s,a)] \;=\; \E_{a \sim \mu}\!\left[\, \frac{\pi(a \mid s)}{\mu(a \mid s)}\, f(s,a)\, \right].
\]
The ratio $w = \pi/\mu$ is the per-action \textbf{importance weight}; large $w$ means
the target policy would have taken this action much more often than the behavior
policy did, and we re-weight accordingly.

Applied to the policy gradient, the off-policy version is
\[
  \nabla_\theta J(\theta) \;=\; \E_\mu\!\left[\, \frac{\pi_\theta(a \mid s)}{\mu(a \mid s)}\, \nabla_\theta \log \pi_\theta(a \mid s)\, G_t\, \right],
\]
or, if we want to reweight an entire trajectory $\tau$ collected under $\mu$,
\[
  \nabla_\theta J(\theta) \;=\; \E_{\tau \sim \mu}\!\left[\, \prod_{t=0}^T \frac{\pi_\theta(a_t \mid s_t)}{\mu(a_t \mid s_t)}\, \sum_{t=0}^T \nabla_\theta \log \pi_\theta(a_t \mid s_t)\, G(\tau)\, \right].
\]

\begin{intbox}[title={Why per-trajectory IS is doomed}]
A product of $T$ ratios is a recipe for variance disaster.
If even one ratio is $0.1$ or $10$, the whole trajectory's contribution swings by an
order of magnitude.
Over hundreds of timesteps the products either explode or collapse to zero, and the
gradient estimate becomes useless.
\textbf{Per-decision} importance sampling --- correcting only up to time $t$, not the
whole future --- helps somewhat, but the underlying problem (unbounded variance from
ratio products) remains.
\end{intbox}

In practice, deep-RL methods deal with this in one of two ways.
\textbf{Constrain} the policy to stay close to the data-collecting policy, so the
ratios stay bounded (TRPO, PPO --- next section).
Or \textbf{avoid} importance sampling entirely and rely on a replay buffer + bootstrapped
TD targets, which sidesteps the distribution-mismatch problem at the cost of additional
bias (DDPG, SAC --- continuous-control section).

\section{Stable Updates: NPG, TRPO, PPO}

Even on-policy, large updates are dangerous.
A small step in $\theta$ can move action probabilities a lot --- especially when the
policy is near-deterministic and the score function blows up.
A bad update produces a worse policy, which produces worse data, which produces an
even worse update. Easy to spiral.

The cure is to control updates in \emph{policy space}, not parameter space.

\subsection{Measuring policy change with KL divergence}

The natural way to measure ``how different are these two policies?'' is the
KL divergence over action distributions:
\[
  \KL{\pi_{\theta_{\text{old}}}}{\pi_\theta}
  \;=\; \E_{s \sim d_{\pi_{\theta_{\text{old}}}}}\!\left[\, \sum_a \pi_{\theta_{\text{old}}}(a \mid s) \log \frac{\pi_{\theta_{\text{old}}}(a \mid s)}{\pi_\theta(a \mid s)}\, \right].
\]
Small KL means similar action distributions in the states the old policy actually visits.
The constrained update is then
\[
  \max_\theta J(\theta) \quad \text{s.t.} \quad \KL{\pi_{\theta_{\text{old}}}}{\pi_\theta} \le \epsilon.
\]
The trust region $\epsilon$ caps how far the policy can drift in one update --- which
keeps both the gradient estimates and the importance ratios well-behaved.

\subsection{Natural Policy Gradient (NPG)}

The plain gradient $\nabla_\theta J$ is the steepest direction in \emph{parameter}
space, not in policy space.
The Natural Policy Gradient (Kakade, 2002) rescales it by the inverse of the local
curvature of KL:
\[
  \theta \;\leftarrow\; \theta + \alpha\, F^{-1}\, \nabla_\theta J(\theta),
\]
where $F$ is the \textbf{Fisher information matrix},
$F = \E[\nabla \log \pi_\theta\, \nabla \log \pi_\theta\T]$, the local Hessian of the
KL in $\theta$.

\begin{intbox}
You can think of NPG as gradient ascent in a coordinate system where unit distance
corresponds to a unit of KL divergence.
A step of fixed length always changes the policy by roughly the same amount, regardless
of how the parameters happen to be scaled.
The price is the Fisher matrix, which is expensive to invert for neural networks
(usually approximated with conjugate gradient).
And NPG only uses a \emph{local} approximation to KL --- nothing prevents a single
large step from violating the trust region.
\end{intbox}

\subsection{Trust Region Policy Optimization (TRPO)}

TRPO (Schulman et al., 2015) makes the KL constraint explicit, optimizing a
\textbf{surrogate objective} that's equivalent to $J(\theta)$ for small KL:
\[
  \max_\theta\; \E_{(s,a) \sim \pi_{\theta_{\text{old}}}}\!\left[\, \underbrace{\frac{\pi_\theta(a \mid s)}{\pi_{\theta_{\text{old}}}(a \mid s)}}_{r_t(\theta)}\, \hat A_t\, \right]
  \quad \text{s.t.} \quad \KL{\pi_{\theta_{\text{old}}}}{\pi_\theta} \le \epsilon.
\]

Three things are happening here.
The expectation is over data collected with $\pi_{\theta_{\text{old}}}$ (so we can reuse
recent rollouts, not just the latest one).
The importance ratio $r_t(\theta)$ corrects for the distribution mismatch.
The KL constraint keeps the ratio close to $1$, which keeps the surrogate a faithful
local approximation of the true objective.

\begin{intbox}[title={Why a surrogate objective?}]
The true objective $J(\theta) = \E_{\tau \sim \pi_\theta}[G(\tau)]$ has two practical
problems:
\begin{itemize}
  \item We can't evaluate it without fresh data from $\pi_\theta$.
  \item The data we have was collected under a different policy.
\end{itemize}
The surrogate solves both: it uses old data (plus a ratio correction) and it's
\emph{provably} a local approximation to $J$ when KL stays small. Solve the surrogate
with the trust region; you've effectively done a controlled step on the true objective.
\end{intbox}

TRPO is robust and sample-efficient, but algorithmically heavy --- second-order
optimization, conjugate gradient, and a line search to enforce the KL bound.

\subsection{Proximal Policy Optimization (PPO)}

PPO (Schulman et al., 2017) replaces TRPO's hard KL constraint with a much simpler
trick: \textbf{clip the ratio} so the optimization can never benefit from a too-large
update:
\[
  L^{\text{CLIP}}(\theta)
  = \E_{(s,a) \sim \pi_{\theta_{\text{old}}}}\!\left[\, \min\!\Bigl(\, r_t(\theta)\, \hat A_t,\; \mathrm{clip}\bigl(r_t(\theta),\, 1-\varepsilon,\, 1+\varepsilon\bigr)\, \hat A_t\, \Bigr)\, \right],
\]
with $r_t(\theta) = \pi_\theta(a_t \mid s_t)\,/\,\pi_{\theta_{\text{old}}}(a_t \mid s_t)$
and $\varepsilon$ typically around $0.1$--$0.2$.

\begin{thmbox}[title={PPO's clipped objective in plain English}]
For a sample with positive advantage ($\hat A_t > 0$, ``this action was better than average''):
\begin{itemize}
  \item If $r_t(\theta) \le 1 + \varepsilon$, the gradient pushes up the action's
    probability normally.
  \item If $r_t(\theta) > 1 + \varepsilon$, the clip kicks in and the gradient becomes
    zero --- you've already increased the probability enough; stop.
\end{itemize}
Symmetrically for $\hat A_t < 0$: probability is decreased until $r_t(\theta) = 1 - \varepsilon$,
then no further.
The minimum with the unclipped term ensures the optimizer cannot exploit the clip
to undo a prior update.
\end{thmbox}

\begin{intbox}
PPO is essentially TRPO with all the heavy machinery thrown away.
No Fisher matrix, no conjugate gradient, no line search --- just first-order SGD on a
clipped objective.
The trust-region effect now comes from the clip itself: the surrogate stops changing
once the ratio leaves $[1-\varepsilon, 1+\varepsilon]$, so further gradient steps in
that direction do nothing useful.
The result is that PPO can do \emph{multiple epochs of updates over the same batch}
of trajectories --- vanilla actor-critic can only do one before the data goes stale.
That alone is a huge win in sample efficiency.
\end{intbox}

PPO uses GAE for the advantage estimate, often adds an entropy bonus to keep the
policy from collapsing prematurely, and is the workhorse for large-scale on-policy
deep RL: Atari, MuJoCo, robotics simulators, RLHF for language models.

\subsection{NPG, TRPO, PPO at a glance}

\begin{center}
\begin{tabular}{llll}
\toprule
\textbf{Method} & \textbf{Constraint} & \textbf{Cost} & \textbf{Key idea} \\
\midrule
NPG  & implicit KL (Fisher) & second-order & natural gradient / KL geometry \\
TRPO & hard KL bound        & second-order & explicit trust region \\
PPO  & soft via clipping    & first-order  & cheap clipped surrogate \\
\bottomrule
\end{tabular}
\end{center}

All three are on-policy in the sense that data must come from $\pi_{\theta_{\text{old}}}$,
but the trust region lets us \emph{reuse} that data for several update steps.
The trade-off across the three is essentially complexity vs.\ rigor: NPG/TRPO have
cleaner theory; PPO is what people actually run.

\section{Continuous Control}

Everything so far has been algorithm-agnostic about action spaces.
Continuous actions just require a different policy parameterization --- the same
gradient machinery still works.

\subsection{The Gaussian policy}

The standard stochastic policy for continuous control is a diagonal Gaussian:
\[
  A_t \sim \pi_\theta(\cdot \mid S_t) = \Normal(\mu_\theta(S_t),\, \Sigma_\theta(S_t)),
\]
where the network outputs the mean $\mu_\theta(s)$ and (optionally) a state-dependent
covariance $\Sigma_\theta(s)$.
Often the covariance is a single learned vector independent of $s$.
The score function for fixed scalar variance $\sigma^2$ is
\[
  \nabla_\theta \log \pi_\theta(A_t \mid S_t) \;=\; \frac{A_t - \mu_\theta(S_t)}{\sigma^2}\, \nabla_\theta \mu_\theta(S_t).
\]
Read this as: \emph{nudge the mean toward the chosen action $A_t$, weighted by how
unusual the action was given the current mean}.
Combined with a positive advantage, the policy update moves the mean toward
high-advantage samples; combined with a negative advantage, it moves the mean away
from them.

\subsection{Deterministic Policy Gradient (DPG)}

For continuous actions there is also a non-stochastic alternative.
If the critic $q_\mathbf{w}(s,a)$ is differentiable in $a$, we can take gradients
\emph{through} the critic to improve a deterministic policy $a = \mu_\theta(s)$:
\[
  \nabla_\theta J(\theta)
  \;=\; \E_{s \sim d_\mu}\!\left[\, \nabla_a q_\mathbf{w}(s,a)\big|_{a = \mu_\theta(s)}\, \nabla_\theta \mu_\theta(s)\, \right].
\]

\begin{intbox}
The deterministic policy gradient avoids the high-variance ``sample then weight by
score function'' loop.
Instead, it asks the critic directly: ``which way in action space increases value?''
and adjusts the actor's mean accordingly.
The price: we need a critic that is accurate \emph{and smooth} in $a$. If the critic
has sharp peaks or is wrong about the action gradient in some region, the actor
chases bad signals --- which is why DPG-family methods need extra stabilization.
\end{intbox}

\subsection{Deep Deterministic Policy Gradient (DDPG)}

DDPG (Lillicrap et al., 2015) is essentially \emph{DQN for continuous actions}:
DPG, plus a replay buffer for off-policy learning, plus target networks for stability.
\begin{itemize}
  \item \textbf{Critic update} (TD regression on tuples sampled from the buffer):
  \[
    \min_\mathbf{w}\; \E_{(s,a,r,s') \sim \mathcal{D}}\!\left[\bigl(q_\mathbf{w}(s,a) - y\bigr)^2\right],
    \quad y = r + \gamma\, q_{\mathbf{w}^-}\!\bigl(s',\, \mu_{\theta^-}(s')\bigr).
  \]
  \item \textbf{Actor update} (deterministic policy gradient through the critic):
  \[
    \max_\theta\; \E_{s \sim \mathcal{D}}\bigl[\, q_\mathbf{w}(s,\, \mu_\theta(s))\, \bigr].
  \]
  \item Target networks $\mathbf{w}^-, \theta^-$ track the live ones with a slow Polyak
    average $\mathbf{w}^- \leftarrow \tau \mathbf{w} + (1-\tau)\mathbf{w}^-$.
  \item Exploration is added by hand: $a_t = \mu_\theta(s_t) + \mathcal{N}_t$ with some
    Gaussian or Ornstein--Uhlenbeck noise.
\end{itemize}

\begin{intbox}[title={DQN $\leftrightarrow$ DDPG correspondence}]
Both algorithms are off-policy + replay + target nets + bootstrapped TD targets, with
the same critic structure $q(s,a)$. The only difference is how the greedy action is
obtained: DQN takes $\argmax_a q(s,a)$ explicitly (cheap for discrete actions);
DDPG learns a parameterized maximizer $\mu_\theta(s) \approx \argmax_a q_\mathbf{w}(s,a)$
(necessary for continuous actions). DDPG = DQN + a learned argmax.
\end{intbox}

DDPG works, but it is brittle. The critic tends to \emph{overestimate} Q-values, and
those errors feed straight into the actor.

\subsection{TD3}

TD3 (Fujimoto et al., 2018) is DDPG with three concrete fixes for that brittleness:
\begin{itemize}
  \item \textbf{Clipped double-Q learning:} keep two critics $q_{\mathbf{w}_1}, q_{\mathbf{w}_2}$
    and use $\min(q_1, q_2)$ in the target.
    Two independent overestimates rarely peak at the same action, so taking the minimum
    suppresses the bias.
  \item \textbf{Delayed policy updates:} update the critic on every step but the actor
    only every $d$ steps. Lets the critic catch up before chasing it.
  \item \textbf{Target policy smoothing:} add a small clipped Gaussian noise to the
    target action $\mu_{\theta^-}(s') + \tilde\varepsilon$ in the critic target. Prevents
    the actor from exploiting sharp peaks in the (still imperfect) critic.
\end{itemize}
The result is the standard ``clean'' off-policy continuous-control baseline --- TD3
on MuJoCo benchmarks is what every new paper is compared against.

\subsection{Maximum-entropy RL and SAC}

DDPG and TD3 still use \emph{deterministic} policies, with exploration tacked on as
external action noise. That noise is brittle (state-independent, often poorly tuned)
and the policy can collapse to a single suboptimal mode early in training.

The cleaner approach is to learn a \textbf{stochastic} policy whose stochasticity is
itself part of the objective. \textbf{Maximum-entropy RL} adds an entropy bonus:
\[
  J_{\text{ent}}(\theta)
  \;=\; \E_{\pi_\theta}\!\left[\, \sum_{t=0}^{T} \gamma^t \bigl( R_{t+1} + \alpha\, \Ent(\pi_\theta(\cdot \mid S_t)) \bigr)\, \right],
\]
where $\Ent(\pi(\cdot \mid s)) = -\E_{a \sim \pi}[\log \pi(a \mid s)]$ is the entropy of
the action distribution and $\alpha > 0$ is a temperature.

\begin{intbox}[title={Why entropy as objective, not regularizer?}]
Adding entropy to the objective changes \emph{which policy is optimal}, not just how we
search for it.
The optimal policy for the augmented objective has high reward \emph{and} high
stochasticity.
This means: sustained exploration (the policy refuses to become deterministic too
early), robustness to value errors (a stochastic policy hedges across action choices),
and smoother updates.
A small entropy bonus is sometimes added to PPO too, but in SAC it is the principled
core of the method.
\end{intbox}

Soft Actor-Critic (SAC) (Haarnoja et al., 2018) puts the maximum-entropy objective
together with the off-policy machinery of DDPG/TD3:
\begin{itemize}
  \item Replay buffer (off-policy).
  \item Two critics with target networks, trained with the entropy-augmented TD target
  \[
    y \;=\; r + \gamma\, \E_{a' \sim \pi_\theta}\!\bigl[\, q_{\mathbf{w}^-}(s', a') - \alpha \log \pi_\theta(a' \mid s')\, \bigr].
  \]
  \item Stochastic Gaussian (typically squashed-Gaussian) policy, updated to maximize
    $\E_{a \sim \pi_\theta}[\, q_\mathbf{w}(s,a) - \alpha \log \pi_\theta(a \mid s)\, ]$
    via the reparameterization trick.
  \item Optionally, learn $\alpha$ itself by gradient descent against a target entropy.
\end{itemize}

\subsection{When to use what}

Modern deep RL has roughly three workhorse algorithms, each with a distinct sweet spot:

\begin{center}
\begin{tabular}{llll}
\toprule
                & \textbf{PPO}        & \textbf{TD3}       & \textbf{SAC} \\
\midrule
Policy          & stochastic           & deterministic       & stochastic \\
Data regime     & on-policy            & off-policy          & off-policy \\
Best fit        & stable large-scale   & clean continuous    & robust, sample-efficient \\
                & training             & control baseline    & control \\
Use cases       & simulators, LLMs,    & benchmark control,  & robotics, real-world \\
                & games (RLHF)         & precise actuation   & control \\
Main trade-off  & data hungry          & weaker exploration  & more moving parts \\
\bottomrule
\end{tabular}
\end{center}

\begin{intbox}
Choosing between PPO, TD3, and SAC is a \emph{systems decision}, not just an
algorithmic one.
\textbf{PPO}: when you can run thousands of parallel environments cheaply and want
something simple that almost always works.
\textbf{TD3}: when you want the cleanest baseline for continuous benchmarks.
\textbf{SAC}: when each environment step is expensive (real robots) and you need
sample efficiency plus robust exploration.
\end{intbox}

\section{Recap}

The progression from REINFORCE to modern actor-critic is a story of \textbf{controlled
trade-offs}.

\begin{itemize}
  \item All policy-gradient methods share the form
    $\nabla J = \E[\Psi_t \nabla \log \pi]$.
    Different choices of $\Psi_t$ trade bias for variance.
  \item Vanilla actor-critic has three weaknesses --- on-policy, unstable, noisy ---
    and modern algorithms each address them differently.
  \item \textbf{Trust-region methods (TRPO/PPO)} stabilize on-policy updates by
    constraining the policy in KL space. PPO is what people actually run.
  \item \textbf{Off-policy actor-critics (DDPG/TD3/SAC)} reuse data via a replay buffer,
    use a parameterized maximizer for continuous actions, and stabilize the critic with
    target networks, double-Q learning, and (in SAC) maximum-entropy policies.
\end{itemize}

The next chapter zooms out: even with the best of these algorithms, deep RL still fails
on many problems. The bottlenecks are usually not algorithmic --- they're about
exploration, credit assignment, and state representation.


% ==========================================================
%  CHAPTER 6 — CORE CHALLENGES IN DEEP RL
% ==========================================================
\chapter{Core Challenges in Deep RL}

By now we have an algorithmic toolbox: DQN and its variants for value-based discrete
control, PPO for stable on-policy optimization, DDPG / TD3 / SAC for off-policy continuous
control. So why does deep RL still fail on so many real problems?

The honest answer is that the algorithms are usually not the bottleneck. The same
PPO that solves a benchmark in a few hours can fail completely on a problem that
\emph{looks} similar but has a sparser reward, a higher-dimensional observation, or a
longer horizon. Performance depends on the environment, the reward design, and the
state representation almost as much as on the algorithm itself.

Three core challenges sit underneath most failure modes:
\begin{itemize}
  \item \textbf{Exploration} --- collecting data that is informative and reward-relevant,
    not just data.
  \item \textbf{Credit assignment} --- figuring out which past actions caused the current
    outcome.
  \item \textbf{State representation} --- learning compact, informative latent states from
    high-dimensional, possibly partial observations.
\end{itemize}
These are deeply interconnected: better representations help both exploration and
credit assignment; better exploration provides the data needed for accurate credit
signals; accurate credit assignment guides the representation toward decision-relevant
features.

This chapter walks through each in turn.

\section{Exploration}

Exploration in RL means collecting data that is \emph{informative} and \emph{reward-relevant}.
Too little exploration and the agent locks in on a bad policy; too much random
exploration and it wastes samples on irrelevant states.
The goal isn't ``more randomness'' --- it is exploring the right states at the right
time.

\subsection{Why exploration is hard}

Three structural features make exploration genuinely difficult in deep RL:
\begin{itemize}
  \item \textbf{Sparse rewards.} Useful feedback is rare. The agent could stumble around
    for thousands of steps without seeing a single non-zero reward.
  \item \textbf{High-dimensional state and action spaces.} The combinatorial explosion of
    possibilities makes random exploration vanishingly unlikely to hit anything good.
  \item \textbf{Long horizons.} A correct decision early on may pay off only much later
    --- and the agent has to keep making the right calls in between.
\end{itemize}

\begin{exbox}[title={Montezuma's Revenge}]
The classic poster child for hard exploration. The agent navigates an Atari maze of
rooms, ladders, traps, and keys. Reward is given only after long, precise sequences
of actions (pick up a key, descend, dodge skulls, climb back up, open a door).
For years, plain DQN scored \emph{zero} --- the policy never reached the first
reward and so had nothing to learn from.
Algorithms only made progress by adding directed, exploration-aware bonuses.
\end{exbox}

\subsection{Baseline exploration: random and noise-based}

The exploration tools we have already met are basic and undirected.
\begin{itemize}
  \item \textbf{$\varepsilon$-greedy} (value-based RL): act greedily with probability
    $1 - \varepsilon$, else uniform random.
    Simple, widely used, fine in small or dense-reward problems --- hopeless when good
    behavior is more than a few random actions away.
  \item \textbf{Action-space noise} (DDPG/TD3): $a_t = \mu_\theta(s_t) + \varepsilon$
    with $\varepsilon \sim \Normal(0, \sigma^2)$.
    Easy to add for continuous actions but the noise is independent across time and
    state, leading to jittery, undirected exploration.
  \item \textbf{Entropy regularization} (PPO/SAC): add an entropy bonus to the objective
    so the policy doesn't collapse too soon.
    Preserves stochasticity, but ``stochastic'' isn't the same as ``directed.''
  \item \textbf{Noisy networks} (Rainbow): replace deterministic linear layers with
    stochastic ones whose noise scales are \emph{learned}, giving state-dependent
    exploration. A clear improvement over $\varepsilon$-greedy --- but still implicit,
    no explicit notion of what is unknown.
\end{itemize}

\begin{intbox}[title={What's missing from these methods}]
All of the above are essentially answers to ``how should I randomize my actions?''
None of them answer the actually important question: \emph{where should I explore next?}
That question requires a notion of \textbf{uncertainty}, \textbf{novelty}, or
\textbf{surprise} --- something that distinguishes ``state I have already seen
plenty of'' from ``state I know nothing about.''
\end{intbox}

\subsection{Three guiding principles for directed exploration}

The literature on directed exploration revolves around three ideas:
\begin{itemize}
  \item \textbf{Uncertainty:} explore where the value estimate is uncertain.
  \item \textbf{Novelty:} explore states the agent has rarely visited.
  \item \textbf{Surprise:} explore transitions that the agent's model fails to predict.
\end{itemize}
These overlap in practice and most modern methods combine more than one.

\subsection{Uncertainty-based exploration}

The principled view: in early training the agent has \textbf{epistemic uncertainty}
about the true value function (the kind of uncertainty that shrinks as you collect more
data). Exploration should target state-action pairs where that uncertainty is high.

The bandit literature gives two canonical strategies that scale to deep RL:
\begin{itemize}
  \item \textbf{UCB (optimism in the face of uncertainty):} act on the upper confidence bound,
    $a_t = \argmax_a \bigl[\mu_t(a) + \beta\, \sigma_t(a)\bigr]$. Combines the
    estimated value with a bonus for high uncertainty.
  \item \textbf{Thompson sampling:} sample a value-function hypothesis from the posterior,
    act greedily under that sample. Probabilistic and naturally exploration--exploitation
    balanced.
\end{itemize}

Both require a way to estimate uncertainty over $Q$-values. Deep RL has no tractable
posterior over neural-network parameters, so we approximate.

\paragraph{Ensembles.}
Maintain $K$ different $Q$-networks $Q^{(1)}, \dots, Q^{(K)}$ trained on the same data.
The variance across the ensemble at $(s,a)$ is a proxy for epistemic uncertainty.
Use it for UCB-style action selection or Thompson-style sampling of an acting head.

\paragraph{Bootstrapped DQN.}
Train each ensemble member on a different bootstrap subsample of the replay buffer
(via per-transition masks). Sample one acting head per episode and act greedily under
it for the whole episode. This produces \textbf{deep exploration}: temporally
consistent, multi-step exploration, rather than the per-step noise of $\varepsilon$-greedy.

\paragraph{Randomized priors.}
A subtle failure mode of ensembles: in regions with little data, all heads can drift to
the same uninformative value, becoming overconfidently identical. The fix is to add a
\emph{fixed random prior function} $p^{(j)}(s,a)$ to each head:
$\hat Q^{(j)}(s,a) = Q^{(j)}(s,a) + \lambda\, p^{(j)}(s,a)$.
Train only $Q^{(j)}$; keep $p^{(j)}$ fixed. The frozen priors guarantee head diversity
in unvisited regions, sustaining exploration even under sparse rewards.

\subsection{Novelty-based exploration: intrinsic motivation}

A different strategy: shape the reward so that visiting novel states is intrinsically
rewarding. Combine the environment reward with an \textbf{intrinsic} bonus:
\[
  r_t \;=\; r_t^{\text{ext}} + \beta\, r_t^{\text{int}},
\]
where $r_t^{\text{int}}$ measures how new or surprising the current transition is.

\paragraph{Count-based exploration.}
The clean tabular version: $r_t^{\text{int}}(s_t) = 1 / \sqrt{N(s_t)}$ where $N(s_t)$ is
the visitation count. Rare states get a big bonus that decays as they are visited
more.
In deep RL exact counts are impossible (no two pixel observations are ever identical),
so we approximate via \textbf{pseudo-counts} from a density model, or via state hashing
to map continuous observations to discrete bins.

\paragraph{Curiosity / prediction error.}
Instead of counting, measure how much the agent's model \emph{fails} to predict.
\textbf{Random Network Distillation (RND)} fixes a random target network $f_{\text{target}}$
and trains a predictor $f_\theta$ to mimic it on visited states.
The intrinsic reward is the prediction error
$r_t^{\text{int}} = \norm{f_\theta(s_t) - f_{\text{target}}(s_t)}^2$. Unseen states give
high error; familiar states have low error. Simple, no environment model required, and
remarkably effective on hard-exploration Atari games.

The \textbf{Intrinsic Curiosity Module (ICM)} learns a feature encoder $\phi(s)$ together
with a forward model $\hat\phi(s_{t+1}) = f(\phi(s_t), a_t)$ and rewards prediction error
in feature space, $r_t^{\text{int}} = \norm{\hat\phi(s_{t+1}) - \phi(s_{t+1})}^2$.

\paragraph{Disagreement-based.}
Train an ensemble of forward models, use the variance across their predictions as
the intrinsic reward. High disagreement = uncertain dynamics = worth exploring.
This isolates \emph{epistemic} uncertainty (which more data can fix) from
\emph{aleatoric} noise (which it cannot).

\paragraph{Combining signals (NGU).}
Modern methods combine multiple novelty signals.
\textbf{Never Give Up} (Badia et al., 2020) uses both \emph{episodic novelty} (how
often the state was visited within the current episode) and \emph{lifelong novelty}
(RND-style: how unfamiliar across all training so far).
The episodic component encourages within-episode exploration; the lifelong component
maintains long-term exploration pressure.

\begin{notebox}[title={The noisy-TV problem}]
Curiosity-style intrinsic rewards have a famous failure mode.
Imagine an Atari-like environment containing a TV that displays random noise. The
prediction error of any model on the TV's pixels is permanently high --- the noise
is unpredictable by definition. A naive curiosity-driven agent will sit in front of
the TV forever, collecting maximum intrinsic reward and learning nothing useful.
The fundamental issue: prediction error doesn't distinguish ``unfamiliar but learnable''
from ``random and unlearnable.''
Practical mitigations: normalize/clip intrinsic rewards, anneal $\beta$ over training,
and combine multiple novelty signals (e.g.\ NGU's episodic + lifelong).
\end{notebox}

\subsection{Where exploration fits in the pipeline}

Exploration can be injected at three different levels of the RL pipeline:
\begin{itemize}
  \item \textbf{Action selection:} $\varepsilon$-greedy, UCB, Thompson sampling, action noise.
  \item \textbf{Model / value estimation:} noisy networks, ensembles, randomized priors.
  \item \textbf{Reward shaping:} count-based bonuses, RND, ICM, disagreement.
\end{itemize}
The right level depends on the problem.
For dense-reward control, action noise is usually enough.
For sparse-reward, long-horizon, or partially-observable problems, you almost always
need something at the value-estimation or reward-shaping level.

\section{Credit Assignment}

The \textbf{credit assignment problem}: when an outcome arrives, which past actions
caused it?
This is hard because most real problems have delayed and sparse feedback, with many
irrelevant decisions sandwiched in between. The agent needs to disentangle informed
choices from serendipitous outcomes --- and credit assignment is what does that.

\begin{intbox}[title={Credit vs.\ value}]
Value asks ``what return do I expect from this state?''
Credit asks ``how much did this specific action contribute to that outcome?''
Two related but distinct questions. A good value estimate doesn't automatically give
good credit signals --- if the value drifts equally with both helpful and irrelevant
actions, every action gets the same credit.
\end{intbox}

\subsection{Three root causes}

Hard credit assignment usually comes from one of three structural features of the MDP:
\begin{itemize}
  \item \textbf{Depth} --- the temporal distance between an action and its outcome is large.
    TD errors shrink exponentially as they propagate backward, so important actions far
    from the reward get tiny gradients. Short-sighted methods develop a \emph{recency
    bias}, attributing credit to whichever action happened to be near the reward.
  \item \textbf{Density} --- few of the actions actually influence the outcome
    (sparse-reward setting). Rare goal states mean most TD errors carry no information,
    so updates are dominated by zero-information samples.
  \item \textbf{Breadth} --- many alternative paths lead to the same outcome.
    Credit gets diluted across all the equivalent trajectories, with no clear bottleneck
    action to learn.
\end{itemize}
Knowing which of the three is dominant tells you which technique to reach for. No single
method solves all three.

\subsection{Time-contiguity methods (depth)}

The default toolbox for moderate delays:
\begin{itemize}
  \item \textbf{TD($1$):} reward propagates backward one step per update. Slow on long horizons.
  \item \textbf{$n$-step returns:} reward propagates back $n$ steps in one shot. Faster, but
    higher variance.
  \item \textbf{Eligibility traces / TD($\lambda$):} exponentially-weighted combination of
    all $n$-step returns. The same idea on the policy-gradient side is \textbf{GAE}
    from Chapter~4, which is now the default advantage estimator.
\end{itemize}
Larger $n$ or larger $\lambda$ accelerates credit propagation at the cost of variance.
These methods address depth directly --- but if the depth is hundreds of steps and the
reward is sparse, even $n$-step methods struggle.

\subsection{Reducing breadth noise with advantages}

Many actions in a trajectory simply don't matter --- they don't change the outcome.
Naively, all of them get blamed/credited for the eventual reward, contributing noise.
The advantage function $A(s,a) = q(s,a) - v(s)$ already does the right thing here:
only \emph{above-} or \emph{below-average} actions get nonzero credit, irrelevant ones
average out.
TD handles \emph{when} reward is propagated; advantage handles \emph{whether the action
mattered}. They are complementary, which is why GAE-style advantage estimation has
become standard.

\subsection{Return decomposition (RUDDER) for very long delays}

When the reward is genuinely far from its cause, even $n$-step / $\lambda$-returns are
not enough.
\textbf{RUDDER} (Arjona-Medina et al., 2019) takes a different approach: \emph{move the
reward to the important event}.

The idea: train a sequence model (e.g.\ an LSTM) to predict the trajectory's future
return $\hat G_t$ from $(s_0, a_0, \dots, s_t)$. Then redistribute reward so that each
step gets the \emph{change} in predicted return:
\[
  \bar r_{t+1} \;=\; \hat G_{t+1} - \hat G_t.
\]
A large jump in predicted return marks an important event (``ah, this is when I
realized I'd win''). After redistribution, the long-horizon, sparse-reward problem
becomes a much easier local-credit problem, because reward now arrives at the actual
causal moment.

\subsection{Hindsight Experience Replay (HER) for sparse goals}

For \textbf{goal-conditioned} problems with binary reward (you reached the goal or you
didn't), most early trajectories fail and provide no learning signal --- the density
challenge in pure form.

\textbf{Hindsight Experience Replay} (Andrychowicz et al., 2017) flips the question:
every failed trajectory \emph{succeeded at some other goal}, namely the state it
actually reached. Relabel transitions as
\[
  (s_t, a_t, r_t, g) \;\longrightarrow\; (s_t, a_t, r'_t, g'),
\]
where $g' = s_T$ (or some achieved state) and $r'_t = 1$ iff $s_t = g'$. Train on
both the original and the relabeled experience.

\begin{intbox}
HER turns the sparse failure signal into dense training data: even when the agent
never reaches the original goal, it can still learn a useful goal-conditioned policy
by reaching whatever it actually reached. Over many episodes the agent learns to
reach a wide variety of subgoals --- and as that capability accumulates, reaching the
intended goal becomes feasible.
The catch: the MDP must be \emph{goal-conditioned}, with reward structure that
generalizes across goals. Not every problem has that shape.
\end{intbox}

\subsection{Practical takeaways}

\begin{itemize}
  \item \textbf{Mild delay:} TD, $n$-step, GAE. The default tools, almost always present.
  \item \textbf{Long delay:} return decomposition (RUDDER) when delays are extreme and
    causal events are sparse.
  \item \textbf{Sparse goal-conditioned reward:} hindsight relabeling (HER).
  \item Diagnose first: depth vs.\ density vs.\ breadth. The right method depends on
    which one is binding.
\end{itemize}
Credit assignment is also tightly coupled to representation: bad features make accurate
credit assignment impossible, no matter the algorithm.

\section{State Representation Learning}

Raw observations are usually high-dimensional (image pixels, lidar scans), noisy, and
often partial.
A good \textbf{state representation} compresses these into a latent vector that keeps
the control-relevant information and discards the rest.

\subsection{End-to-end RL vs.\ SRL + RL}

\paragraph{End-to-end.}
Train a single policy/value network straight from observations to actions. The early
layers implicitly learn features. This is the default in Atari-style benchmarks.
The downsides: features are entirely shaped by the (often weak) RL gradient, generalization
across tasks is poor, and sample efficiency is low.

\paragraph{SRL + RL.}
Decouple the encoder from the policy. First (or jointly) train an encoder
$\phi: o \to x$ that produces a useful latent state, then run RL on top of $x$.
Modular, often more sample-efficient, and the encoder transfers across tasks.

The trade-off is essentially \emph{generality of the representation vs.\ alignment with the
control objective}. Pure end-to-end is maximally aligned but data-hungry; pure SRL is
maximally general but the auxiliary objective may not match what the policy needs.

\subsection{What makes a representation good}

Three desiderata are usually invoked:
\begin{itemize}
  \item \textbf{Structure:} the latent space should reflect the causal, task-relevant
    structure of the environment, not arbitrary visual variation.
  \item \textbf{Continuity:} nearby states in the world should map to nearby latents.
    Small observation changes should produce small representation changes.
  \item \textbf{Sparsity / disentanglement:} only a small number of latent dimensions
    should be active per state. Reduces spurious correlations and makes downstream
    learning easier.
\end{itemize}

\subsection{A taxonomy of SRL methods}

The course groups online model-free SRL into six (overlapping) classes.

\paragraph{Metric-based / bisimulation.}
Shape the latent metric so that two states are close iff they produce similar rewards
and similar transition dynamics. Bisimulation is the principled formalization
(Ferns et al., 2004): states are \emph{bisimilar} iff for every action they yield the
same reward and transition to bisimilar states. Strong abstraction, expensive to compute.

\paragraph{Auxiliary tasks.}
Add side losses to the encoder: \emph{reconstruction} (recover the observation from the
latent), \emph{forward dynamics} (predict $\phi(s_{t+1})$ from $\phi(s_t)$ and $a_t$),
\emph{inverse dynamics} (predict $a_t$ from $\phi(s_t)$ and $\phi(s_{t+1})$),
\emph{reward prediction}. Provides a richer signal than the RL gradient alone, especially
under sparse reward. The risk: an auxiliary loss may not align with what the controller
needs (reconstruction wastes capacity on visually salient but task-irrelevant pixels).

\paragraph{Data augmentation.}
Apply visual augmentations (random crops, color jitter, masking, noise injection) to
observations and require the encoder to be invariant. Cheap, and it consistently
improves both sample efficiency and generalization in pixel-based RL. The caveat: too
strong an augmentation can destroy task-relevant cues.

\paragraph{Contrastive SRL.}
Train the encoder so that positive pairs (same state under different augmentations,
or temporally close observations) are close in latent space and negative pairs (random
other states) are far. The standard NCE-style loss
\[
  \mathcal{L}_{\text{NCE}}
  \;=\; -\E\!\left[\, \log \frac{\exp(\mathrm{sim}(\phi(o), \phi(o^+)))}
                          {\exp(\mathrm{sim}(\phi(o), \phi(o^+))) + \sum_i \exp(\mathrm{sim}(\phi(o), \phi(o_i^-)))}\, \right]
\]
encourages a clustered, well-structured latent space without needing reconstructions
or labels.

\paragraph{Non-contrastive.}
Same goal as contrastive (similar states close in latent space) but without explicit
negative pairs --- using e.g.\ predictor networks and target encoders to prevent
representational collapse (BYOL-style methods).

\paragraph{Attention-based.}
Learn attention masks over spatial or temporal features so that the encoder focuses on
decision-relevant regions and suppresses distractors. Improves interpretability and
robustness; adds architectural complexity.

\subsection{Why representation matters in practice}

A weak representation makes everything downstream harder.
Bad features mean noisy value estimates, which produce noisy gradients, which produce
noisy policies --- and the same algorithm that succeeds on a clean benchmark fails
on a visually messy real-world task.
Most progress in pixel-based deep RL over the last few years has come from better
encoders (data augmentation, contrastive losses) rather than better RL algorithms.

\section{Recap}

Modern deep RL is bottlenecked by three intertwined problems, none of which is solved
by ``pick a better algorithm.''

\begin{itemize}
  \item \textbf{Exploration:} how to discover rewarding behavior in large, sparse, or
    deceptive environments.
    Tools: noisy networks, uncertainty ensembles, bootstrap with randomized priors,
    intrinsic motivation (RND, ICM, disagreement, NGU).
  \item \textbf{Credit assignment:} how to attribute outcomes to past actions across long,
    sparse, or branching horizons.
    Tools: TD / $n$-step / GAE for moderate delays, advantages for breadth, RUDDER for
    long delays, HER for sparse goal-conditioned reward.
  \item \textbf{State representation:} how to learn compact, informative latents from
    high-dimensional observations.
    Tools: bisimulation, auxiliary tasks, data augmentation, contrastive and non-contrastive
    SRL, attention.
\end{itemize}

These three challenges feed into one another. Better representations make exploration
and credit assignment tractable. Better exploration provides the data needed for
learning representations and assigning credit. Better credit assignment guides the
representation toward features that actually matter for control. Real-world deep RL
is usually about getting all three right simultaneously --- not finding a single magic
algorithm.


% ==========================================================
%  CHAPTER 7 — DATA-CENTRIC RL: IMITATION LEARNING & OFFLINE RL
% ==========================================================
\chapter{Imitation Learning and Offline RL}

Every algorithm so far has been \emph{online}: the agent acts, collects fresh
transitions, updates its policy, and then collects more data with the improved policy.
Even off-policy methods with replay buffers still gather new environment experience
as they go. In many real domains this loop is impossible. You cannot test an untrained
policy on patients, explore freely in live traffic, wear out a real robot, or experiment
on millions of live users. The data already exists; the question is whether we can learn
a good policy \emph{from data alone}.

This chapter studies the \textbf{data-centric} regime, where the data set is fixed.
Two settings sit at opposite ends of a spectrum. \emph{Imitation learning} has
demonstrations but no reward and merely tries to reproduce expert behavior.
\emph{Offline RL} has reward-labelled logs and tries to do \emph{better} than the data
that generated them. Both face the same underlying obstacle: a policy trained on one
distribution of states and actions must generalize to a different distribution at
deployment. Controlling that \textbf{distributional shift} is the unifying theme.

\begin{intbox}
Three learning paradigms, distinguished by where the data comes from.
\textbf{Online RL:} the agent collects its own data with the current policy.
\textbf{Off-policy RL:} a replay buffer mixes old and new data, but the agent still
acts in the environment. \textbf{Offline (batch) RL:} a fixed data set $\dataset$
collected once by some behavior policy $\pi_\beta$, with \emph{no} further interaction.
Offline RL can only learn behaviors that are sufficiently represented in $\dataset$ ---
you cannot recover what was never tried.
\end{intbox}

\section{Imitation Learning}

The simplest fixed-data problem is to learn from \emph{expert actions with no rewards}.
We are given a demonstration data set
\[
  \dataset = \{(s_i, a_i)\}_{i=1}^N, \qquad a_i \sim \tilde\pi(\cdot \mid s_i),
\]
where $\tilde\pi$ is an (unknown) expert policy, and we want a policy
$\pi_\theta(a \mid s) \approx \tilde\pi(a \mid s)$. We use imitation rather than RL
from scratch because no reward function is required, expert knowledge shortcuts long
exploration, and the imitator is a powerful initialization for subsequent RL.

\subsection{Behavioral Cloning}

\textbf{Behavioral cloning (BC)} treats imitation as ordinary supervised learning:
predict the expert's action from the state.
\[
  \hat{\params} = \argmin_{\params}\; \E_{s \sim p_{\text{data}}}\,
    \loss\bigl(\tilde\pi(s),\, \pi_\theta(s)\bigr).
\]
The loss depends on the action type: mean-squared error for continuous actions and
cross-entropy for discrete ones,
\[
  \loss_{\text{BC}} = \frac{1}{N}\sum_{k=1}^N \norm{\pi_\theta(s_k) - a_k}^2
  \qquad\text{or}\qquad
  \loss_{\text{BC}} = -\frac{1}{N}\sum_{k=1}^N \log \pi_\theta(a_k \mid s_k).
\]
ALVINN (1989), an early self-driving system, was a BC network mapping camera and
range-finder input to a steering direction; it drove real highways but its core
weakness was already visible --- distribution shift between training and deployment.

\subsection{The distribution shift problem}

BC learns from states the \emph{expert} visits, $s \sim p_{\text{data}}(s)$, but at
deployment it visits states \emph{it itself} reaches, $s \sim p_{\pi_\theta}(s)$.
A small prediction error nudges the agent into a state slightly off the expert manifold;
in that unfamiliar state it errs more; errors compound. The training and deployment
distributions diverge:
\[
  p_{\text{data}}(s) \neq p_{\pi_\theta}(s).
\]

\begin{thmbox}[title={Compounding errors (Ross \& Bagnell, 2010)}]
If the learned policy makes a mistake with probability $\epsilon$ under the expert's
state distribution, then over a horizon $T$ its total error grows as
\[
  \text{policy error} \;\le\; \text{expert error} + T^2\epsilon.
\]
The error is \emph{quadratic} in the horizon: there are $T$ opportunities to leave the
expert trajectory, and each early mistake can corrupt up to $T-t$ future states, giving
$\bigO{T^2}$. Even tiny per-step error becomes severe on long trajectories.
\end{thmbox}

One partial fix is to augment the data with recovery behavior. NVIDIA's end-to-end
driving system (2016) mounted left/right cameras in addition to the center one and added
corrective steering labels, effectively teaching the network how to return to the lane
center --- injecting off-distribution recovery states into the data set.

\subsection{DAgger: dataset aggregation}

The principled fix is to make the training distribution match the deployment
distribution. \textbf{DAgger} (Dataset Aggregation) rolls out the \emph{current learner},
visits the states it actually reaches, and asks the expert to label those states.
The data set grows over iterations, and the learner is repeatedly retrained on the
union. This drives the error growth down from $\bigO{T^2}$ to $\bigO{T}$.

\begin{thmbox}[title={DAgger}]
Given expert $\tilde\pi$ and a mixing schedule $\beta_i$, initialize $\dataset$ with
expert demonstrations and train $\pi_{\theta_1}$. For $i = 1, \dots, N$:
\begin{enumerate}
  \item Define the rollout policy $\pi_i = \beta_i\,\tilde\pi + (1-\beta_i)\,\pi_{\theta_i}$
    (mostly expert early on, mostly learner later).
  \item Sample trajectories with $\pi_i$ and collect the states they visit.
  \item Query the expert for labels: $\dataset_i = \{(s, \tilde\pi(s))\}$ on those states.
  \item Aggregate $\dataset \leftarrow \dataset \cup \dataset_i$ and retrain
    $\pi_{\theta_{i+1}}$ on all of $\dataset$.
\end{enumerate}
\end{thmbox}

DAgger's cost is that the expert must be \emph{queryable at training time}. A human
labeller is expensive, a simulator expert needs a simulator, and a real-world expert
(a patient, an unsafe environment) is often unavailable. When the expert cannot be
queried, we must fall back to purely offline methods.

\subsection{Failure modes beyond distribution shift}

\begin{notebox}[title={Three classic imitation traps}]
\textbf{Non-Markovian behavior.} Humans condition on history, not just the current state
($\tilde\pi(a_t \mid o_t, o_{t-1}, \dots) \neq \tilde\pi(a_t \mid s_t)$): checking mirrors
before a lane change, turning after an earlier landmark. A policy that sees only $s_t$
misses these temporal dependencies. \emph{Fix:} recurrent or transformer policies over
the observation history.

\textbf{Causal confusion.} The agent learns spurious correlations instead of causes.
A brake light correlates with braking, so the policy learns ``brake light $\Rightarrow$
brake'' --- but the light appears \emph{because} the agent brakes; the causality is
reversed. Any feature correlated with expert actions can become a false cause.
\emph{Fix:} causal reasoning, disentangled representations.

\textbf{Multimodal behavior.} The expert demonstrates several valid responses to the
same state (pass the obstacle left \emph{or} right). An MSE loss averages the modes,
$\pi_\theta(s) \approx \tfrac{1}{2}(a_{\text{left}} + a_{\text{right}})$ ---
driving straight into the obstacle. \emph{Fix:} mixture-of-Gaussians output heads,
latent-variable models (CVAE, flows), or autoregressive discretization of actions.
\end{notebox}

When we can no longer query the expert, two roads open. If rewards are available, use
\emph{offline RL}; if rewards are unavailable, use \emph{inverse RL}
(\cref{chap:irl}).

\section{Offline Reinforcement Learning}

\textbf{Offline RL} learns a near-optimal policy from a fixed data set without any new
interaction. The data set $\dataset = \{(s^i, a^i, s'^i, r^i)\}$ was collected by a
behavior policy $\pi_\beta$ (possibly unknown, possibly a mixture); write $d^{\pi_\beta}(s,a)$
for its state-action visitation. The RL objective is unchanged,
\[
  \max_\pi\; J(\pi) = \E_{\tau \sim p_\pi(\tau)}\!\left[\sum_{t=0}^{H} \gamma^t r(s_t, a_t)\right],
\]
but we cannot draw new samples from $p_\pi(\tau)$ during training.

\subsection{Why standard off-policy RL is not enough}

In principle DQN or SAC could train on a fixed buffer. In practice they fail. Standard
off-policy methods still collect fresh experience that gradually tracks the current
policy; removing that adaptation mechanism is exactly what makes the setting hard.
With $\dataset$ frozen, the learned policy $\pi$ inevitably drifts away from $\pi_\beta$,
and it begins to ask the value function about actions that were never observed.

The damage enters through the bootstrapped maximization. Q-learning improves the policy
by $\pi(s) = \argmax_a Q(s,a)$ and backs up
\[
  Q(s,a) \leftarrow r(s,a) + \gamma \max_{a'} Q(s', a').
\]
The $\max_{a'}$ queries $Q$ at actions \emph{not in $\dataset$}. A neural network
extrapolates unpredictably and can assign spuriously high values to out-of-distribution
(OOD) actions; the error is then propagated backward by bootstrapping, and the greedy
policy happily selects actions with high predicted but low true value. Offline RL must
estimate values beyond the data support, and naive bootstrapping makes this catastrophic.

\begin{intbox}
The common principle behind every offline RL method: \emph{do not trust the model at
out-of-distribution inputs.} The three families differ only in how they enforce it ---
constrain the policy, penalize the values, or avoid the OOD query altogether.
\end{intbox}

\subsection{Behavior regularization}

The most direct idea is to keep the learned policy close to the behavior policy, so it
never selects unsupported actions:
\[
  \max_\pi\; J(\pi) \quad\text{subject to}\quad D(\pi \,\|\, \pi_\beta) \le \delta.
\]
Implementations add a KL penalty $\loss = -J(\pi) + \lambda\,\KL{\pi}{\pi_\beta}$,
a behavioral-cloning auxiliary loss $\loss = -J(\pi) + \lambda\,\loss_{\text{BC}}$, or a
hard support constraint $\pi(a\mid s) = 0$ whenever $(s,a) \notin \dataset$. The trade-off
is sharp: a tight constraint is safe but limits improvement (the policy collapses toward
BC), while a loose one may still extrapolate dangerously.

\subsection{Conservative Q-Learning (CQL)}

CQL keeps standard Bellman training but explicitly \emph{pushes down} $Q$-values for OOD
actions while pulling up the values of actions actually in the data:
\[
  \min_Q\; \Bigl(\underbrace{\E_{s\sim\dataset,\, a\sim\mu(a\mid s)}[Q(s,a)]}_{\text{minimize OOD }Q}
  \;-\; \underbrace{\E_{(s,a)\sim\dataset}[Q(s,a)]}_{\text{maximize in-data }Q}\Bigr)
  + \loss_{\text{TD}},
\]
where $\mu$ is a broad action distribution (e.g.\ uniform). The resulting $Q_{\text{CQL}}$
is a \emph{lower bound} on the true $Q^\pi$, so a policy maximizing the conservative
$Q$ still improves while overestimation of unseen actions is suppressed. CQL is strong
across locomotion, manipulation, and real-world benchmarks.

\subsection{Implicit Q-Learning (IQL)}

IQL avoids OOD queries entirely. Instead of the OOD-querying target
$Q(s,a) \leftarrow r + \gamma \max_{a'} Q(s', a')$, it replaces the $\max_{a'}$ with an
\emph{expectile regression} on a value function over \emph{dataset actions only}:
\[
  V(s) = \argmin_V\; \E_{a\sim\dataset(s)}\bigl[\loss^\tau\bigl(Q(s,a) - V(s)\bigr)\bigr],
  \qquad
  \loss^\tau(u) = \abs{\tau - \ind[u < 0]}\, u^2.
\]
The loss is asymmetric. When $Q(s,a) > V(s)$ ($u>0$) and $\tau \approx 1$, the penalty
is $\approx u^2$ (large), so high-$Q$ actions strongly push $V$ \emph{up}; when
$Q(s,a) < V(s)$ ($u<0$), the penalty is $\approx (1-\tau)u^2 \approx 0$, so low-$Q$
actions barely pull $V$ down. As $\tau \to 1$, $V(s)$ approaches a soft
$\max_{a\in\dataset} Q(s,a)$ --- the best value \emph{achievable within the data} ---
without ever evaluating an unseen action. A separate policy-extraction step then recovers
the policy.

\subsection{Decision Transformer: RL as sequence modeling}

A radically different route bypasses Bellman backups altogether. The
\textbf{Decision Transformer} treats a trajectory as a sequence of
(return-to-go, state, action) triples,
\[
  \hat R_t = \sum_{t'=t}^{H} r_{t'}, \qquad
  (\hat R_1, s_1, a_1,\; \hat R_2, s_2, a_2,\; \dots),
\]
and trains a GPT-style causal transformer to predict the next action given the
return-to-go. At test time you \emph{condition} on a desired return $\hat R_1 = R_{\text{target}}$
and let the model autoregressively generate actions that achieve it. There are no Bellman
backups and no OOD extrapolation, and it scales with data and model size. Its weakness is
\emph{trajectory stitching}: it cannot easily combine fragments of suboptimal trajectories
into a better-than-any-single-trajectory policy.

\begin{notebox}[title={Stitching: who can and who cannot}]
\emph{Stitching} means assembling a better-than-dataset policy by combining sub-trajectories
from different episodes. Methods built on Bellman backups (CQL, IQL) propagate value
through bootstrapping and \emph{can} stitch. Behavior-regularized methods stitch only
weakly (they are tied to $\pi_\beta$). Sequence models like the Decision Transformer
\emph{cannot} stitch --- they imitate returns they have seen, not values they can compose.
\end{notebox}

\section{Two Sides of the Same Coin}

Imitation learning and offline RL answer different questions but share one disease.

\begin{center}
\begin{tabular}{@{}lll@{}}
\toprule
\textbf{Aspect} & \textbf{Imitation Learning} & \textbf{Offline RL} \\
\midrule
Objective       & match the expert        & maximize reward \\
Data quality    & expert demonstrations   & any $\pi_\beta$ \\
Reward required & no                      & yes \\
Core challenge  & distribution shift      & $Q$-value extrapolation \\
Key algorithm   & DAgger                  & CQL, IQL \\
Upper bound     & expert level            & above the dataset possible \\
\bottomrule
\end{tabular}
\end{center}

In imitation learning, training is on $p_{\text{data}}(s)$ but deployment is on
$p_{\pi_\theta}(s)$; errors compound as $J(\pi_\theta) \le J(\tilde\pi) + T^2\epsilon$,
and DAgger's on-policy labels reduce this to $T\epsilon$. In offline RL, the $Q$-function
is evaluated at OOD $(s,a)$ pairs and overestimation propagates through Bellman backups;
the fixes constrain the policy (behavior regularization), penalize OOD values (CQL), or
avoid the query (IQL). The central lesson is the same throughout: \textbf{learning from a
fixed distribution while generalizing to a new one requires explicit design choices to
control the resulting shift.}


% ==========================================================
%  CHAPTER 8 — INVERSE RL AND PREFERENCE-BASED RL
% ==========================================================
\chapter{Inverse RL and Preference-Based Reward Learning}
\label{chap:irl}

Imitation learning copies the expert's \emph{actions}; offline RL assumes the reward is
given. But in most real tasks the objective is neither demonstrated cleanly nor written
down. What is ``success'' for a manipulation task, ``helpful'' for a dialogue system,
or the right trade-off between comfort and speed for a car? When the reward is implicit,
ambiguous, or multi-objective, \textbf{reward design becomes part of the learning
problem}. This chapter asks: can the data itself reveal the objective?

Different signals answer different questions. Demonstrations show what good behavior
\emph{looks like} (inverse RL). Pairwise preferences say which of two behaviors is
\emph{better} (preference-based RL, RLHF). Scalar feedback rates how good a behavior was
(reward modeling). And AI feedback asks whether the labelling itself can be scaled
(RLAIF). The unifying risk arrives at the end: any learned reward is a \emph{proxy} for
the true objective, and proxies can be exploited.

\section{Why Learn the Reward?}

Imitation captures \emph{what} the expert does but not \emph{why}. A cloned policy
breaks when the environment changes, when the agent's capabilities differ from the
expert's, or when it leaves the demonstrated distribution. \textbf{Inverse RL (IRL)}
instead infers the reward function the expert appears to be optimizing. A reward
\emph{transfers}: a new agent can re-optimize it in new situations. Where forward RL maps
reward $\to$ behavior, inverse RL maps behavior $\to$ reward.

The formal ingredients are an MDP \emph{without} reward,
$\mathcal{M}\setminus r = (\Sset, \Aset, p, \gamma)$, expert demonstrations
$\dataset = \{\tau_i\}$, and a parametric reward family $r_\psi(s,a)$. We seek $r_\psi$
under which the expert is (near-)optimal.

\begin{notebox}[title={IRL is ill-posed}]
Many reward functions explain the same behavior. A constant reward $r(s,a)=c$ makes
\emph{every} policy optimal, and infinitely many non-trivial rewards are also consistent
with a given set of demonstrations. We cannot simply invert the optimality conditions;
we need an extra \emph{principle} to select among feasible rewards. The two classical
principles are maximum margin (Abbeel \& Ng, 2004) and maximum entropy (Ziebart et al.,
2008).
\end{notebox}

\section{Feature Matching and Maximum-Margin IRL}

Assume the reward is linear in hand-designed features,
\[
  r_\psi(s,a) = \psi\T f(s,a),
\]
where $f(s,a)$ describes \emph{what happens} and the weights $\psi$ describe \emph{what
matters}. The value of a policy then depends only on its expected discounted features:
\[
  J(\pi) = \E_\pi\!\left[\sum_{t=0}^\infty \gamma^t r_\psi(s_t,a_t)\right]
         = \psi\T \underbrace{\E_\pi\!\left[\sum_{t=0}^\infty \gamma^t f(s_t,a_t)\right]}_{\mu(\pi)\;\text{feature expectation}}.
\]
Two policies with identical feature expectations earn identical reward. Matching reward
therefore reduces to \emph{matching feature expectations}: find $\psi$ such that the
expert's $\mu(\tilde\pi)$ scores at least as high as any other policy's:
$\psi\T \mu(\tilde\pi) \ge \psi\T \mu(\pi)$ for all $\pi$.

\subsection{The maximum-margin principle}

Among the many $\psi$ that satisfy the constraint, prefer those that make the expert
\emph{clearly} better than alternatives, by a margin $m$:
\[
  \psi\T \mu(\tilde\pi) > \max_{\pi\in\Pi}\, \psi\T \mu(\pi) + m.
\]
This is exactly the support-vector-machine idea. Since $\psi$ and $10\psi$ both satisfy
the constraint, we fix the margin to $1$ and minimize the weight norm instead
($m \propto 1/\norm{\psi}$):
\[
  \min_\psi\; \norm{\psi}^2 \quad\text{s.t.}\quad
  \psi\T \mu(\tilde\pi) \ge \psi\T \mu(\pi) + 1 \quad \forall \pi.
\]
Not all wrong policies are equally wrong. \textbf{Structured} max-margin IRL replaces the
fixed margin with a similarity-dependent one, requiring a larger margin against policies
that differ more from the expert:
\[
  \min_\psi\; \norm{\psi}^2 \quad\text{s.t.}\quad
  \psi\T \mu(\tilde\pi) \ge \psi\T \mu(\pi) + D(\tilde\pi, \pi) \quad \forall \pi,
\]
where $D(\tilde\pi,\pi)$ measures policy dissimilarity. A slack variable can relax these
constraints to allow a suboptimal expert.

\begin{thmbox}[title={Maximum-margin IRL loop}]
Collect demonstrations $\dataset$ and estimate the expert feature expectation
$\tilde\mu$. For $i = 1,\dots,K$:
\begin{enumerate}
  \item Set the reward $r_\psi(s,a) = \psi\T f(s,a)$.
  \item Solve the RL problem under $r_\psi$ to get policy $\pi_\psi$.
  \item Compute its feature expectation $\mu(\pi_\psi)$.
  \item Update $\psi$ so that $\psi\T\tilde\mu > \psi\T\mu(\pi_\psi) + D(\tilde\pi,\pi_\psi)$.
\end{enumerate}
\end{thmbox}

Its limitations are intrinsic: it repeatedly solves a full RL problem, assumes
near-optimal demonstrations, leaves many feasible rewards, and is hard to scale to
nonlinear rewards.

\section{Maximum-Entropy IRL}

Rather than assuming the expert is perfectly optimal, model demonstrations
\emph{probabilistically}. Introduce a binary optimality variable $\mathcal{O}_t$ with
\[
  p(\mathcal{O}_t = 1 \mid s_t, a_t) \propto \exp\bigl(r_\psi(s_t, a_t)\bigr),
\]
so that high-reward actions are \emph{more likely} to be optimal, not strictly required.
Trajectories with larger total reward then become more probable:
\[
  p(\tau \mid \mathcal{O}_{1:T}) \propto p(\tau)\,\exp\!\left(\sum_{t} r_\psi(s_t, a_t)\right).
\]

\textbf{Maximum-entropy IRL} chooses, among all rewards consistent with the data, the one
inducing the \emph{most uniform} (maximum-entropy) trajectory distribution:
\[
  p_\psi(\tau) = \frac{1}{Z_\psi}\, p(\tau)\,\exp\bigl(r_\psi(\tau)\bigr),
  \qquad Z_\psi = \int p(\tau)\,\exp\bigl(r_\psi(\tau)\bigr)\,\dd\tau.
\]
Here $p(\tau)$ encodes the environment dynamics --- how likely the raw trajectory is,
\emph{independent of reward} --- $\exp(r_\psi(\tau))$ up-weights high-reward trajectories,
and $Z_\psi$ normalizes. Maximizing the likelihood of the expert trajectories gives a
gradient with a clean interpretation:
\[
  \grad_\psi \loss
  = \underbrace{\E_{\tau\sim\tilde\pi}\bigl[\grad_\psi r_\psi(\tau)\bigr]}_{\text{increase expert reward}}
  - \underbrace{\E_{\tau\sim p_\psi(\tau)}\bigl[\grad_\psi r_\psi(\tau)\bigr]}_{\text{decrease current-policy reward}}.
\]
Raise the reward of expert trajectories, lower the reward of trajectories the current
model prefers; the gradient is zero exactly when the model's trajectory distribution
\emph{matches} the expert's. This resolves the ambiguity with a unique, principled
solution.

\begin{notebox}[title={MaxEnt IRL: strengths and the partition-function cost}]
Strengths: a principled probabilistic model, a soft optimality assumption that handles
noisy or suboptimal demonstrations, and compatibility with flexible reward
parameterizations. Costs: it repeatedly solves the RL problem, computing state-visitation
frequencies is expensive, it often assumes known dynamics, and it does not fully remove
ambiguity (reward-shaping transformations remain). With high-dimensional function
approximators the key bottleneck is \emph{the partition function $Z_\psi$} --- a sum over
all possible trajectories --- which forces differentiating through $Z$ or running repeated
forward--backward passes.
\end{notebox}

\section{Generative Adversarial Imitation Learning (GAIL)}

MaxEnt IRL alternates expensively between reward learning and policy optimization.
\textbf{GAIL} short-circuits this by matching expert behavior \emph{adversarially}, in the
style of a GAN. A discriminator $D_\psi(s,a)$ learns to tell expert state-action pairs
from policy ones, while the policy $\pi_\theta$ learns to produce pairs that fool it.
The discriminator's output \emph{is} the reward signal:
\[
  r(s,a) = -\log\bigl(1 - D_\psi(s,a)\bigr).
\]
The policy is trained with ordinary policy gradients on this reward; at convergence its
state-action visitation distribution is indistinguishable from the expert's. GAIL needs
no known dynamics and avoids explicit reward recovery --- it learns a policy directly,
treating ``recover reward, then optimize'' as a single adversarial game.

\section{Preference-Based Reward Learning}

Demonstrations require a human to actually \emph{perform} the task --- infeasible when the
expert lacks the robot's physical capability, when the action space is high-dimensional,
or when demonstration is simply tedious. \emph{Comparing} two behaviors is far easier than
producing one. Preference-based RL shows two trajectory clips and asks ``which is better?''
Critiquing scales to domains where demonstration is impossible.

Given a human judgment $\tau_w \succ \tau_l$ (segment $w$ preferred to $l$), the
\textbf{Bradley--Terry} model estimates the preference probability from the segments'
total predicted reward:
\[
  P(\tau_w \succ \tau_l) = \sigma\bigl(r_\theta(\tau_w) - r_\theta(\tau_l)\bigr),
  \qquad r_\theta(\tau) = \sum_{(s,a)\in\tau} r_\theta(s,a).
\]
We learn the reward by maximizing the log-likelihood of observed preferences,
\[
  \max_\theta\; \E_{\tau_w \succ \tau_l}\bigl[\log \sigma\bigl(r_\theta(\tau_w) - r_\theta(\tau_l)\bigr)\bigr],
\]
which is just binary cross-entropy with soft targets; the reward model can be any network.
The overall loop alternates sampling trajectory segments from the current policy, querying
annotators for preferences, updating $r_\theta$ with the Bradley--Terry loss, and running
RL on the learned reward. Human labelling is the bottleneck.

\subsection{Active query selection}

Since annotation is expensive, query the pairs that are most \emph{informative}.
\emph{Uncertainty sampling} picks the pair whose preference the model is least sure about,
$(\tau_i,\tau_j)^* = \argmax \Var_\theta[\sigma(r_\theta(\tau_i) - r_\theta(\tau_j))]$.
\emph{Ensemble disagreement} trains $K$ reward models on the same data and queries the
pairs where they disagree most --- disagreement approximates epistemic uncertainty, hence
the most informative label. This can cut annotation cost by an order of magnitude.

\section{Modern Alignment: RLHF and DPO}

\textbf{RLHF} (Reinforcement Learning from Human Feedback) is the same preference-learning
idea applied at the scale of large language models. The pipeline:
\begin{enumerate}
  \item Pre-train a language model on large text corpora.
  \item Supervised fine-tuning (SFT) on high-quality (prompt, response) pairs.
  \item Collect preferences: for each prompt, humans pick the better of two responses.
  \item Train a reward model $r_\theta(x,y)$ via the Bradley--Terry loss.
  \item RL fine-tuning: maximize $r_\theta$ with PPO, plus a KL penalty from the SFT model.
\end{enumerate}
This is the InstructGPT / ChatGPT recipe --- algorithmically identical to Christiano et al.
(2017), just scaled. \textbf{RLAIF} replaces human annotators with a strong language model
(``which response is less harmful?''), exploiting that critique is easier than generation
even for an AI; Constitutional AI encodes the values as a written set of principles the
model uses to critique and revise its own outputs.

\subsection{Direct Preference Optimization (DPO)}

The reward-model-plus-PPO loop is expensive and sensitive. DPO observes that the
KL-regularized RLHF objective has a closed-form optimal policy,
\[
  \pi^*(y\mid x) \propto \pi_{\text{ref}}(y\mid x)\,\exp\!\left(\tfrac{1}{\beta} r^*(x,y)\right),
\]
which can be \emph{rearranged} to express the reward implicitly in terms of the policy:
\[
  r(x,y) = \beta \log \frac{\pi^*(y\mid x)}{\pi_{\text{ref}}(y\mid x)} + \beta\log Z(x).
\]
Substituting this into the Bradley--Terry likelihood makes the reward $r$ (and the
intractable $Z(x)$) cancel, leaving a pure supervised objective over preference pairs:
\[
  \loss_{\text{DPO}}(\theta) = -\E_{(\tau_w, \tau_l)}\!\left[
    \log \sigma\!\left(\beta \log \frac{\pi_\theta(\tau_w)}{\pi_{\text{ref}}(\tau_w)}
      - \beta \log \frac{\pi_\theta(\tau_l)}{\pi_{\text{ref}}(\tau_l)}\right)\right].
\]
No separate reward model, no RL loop --- just supervised learning on preferences, which is
simpler and more stable than PPO-based RLHF. Its limits: it assumes a fixed offline
preference set and cannot easily absorb a changing reward signal. The lesson is that the
RL loop can sometimes be \emph{bypassed entirely}.

\section{Reward Hacking and Goodhart's Law}

Every method in this chapter shares one vulnerability: \textbf{the learned reward is a
proxy for the true objective, and any proxy can be exploited.} The IRL reward depends on
feature design and the demonstrations; the preference reward is trained on finite,
possibly biased comparisons; a goal classifier is fooled by states outside its training
distribution; the RLHF reward model is over-optimized by the policy.

\begin{intbox}
\textbf{Goodhart's Law:} ``When a measure becomes a target, it ceases to be a good
measure.'' The agent optimizes the proxy $r_\psi$, not the true reward $r^*$. Classic
examples: the OpenAI boat-race agent that spins in circles collecting speed-boost tokens
instead of finishing; a robot that moves the camera so an object \emph{appears} grasped;
an RLHF model that learns verbose, sycophantic answers that score well. The common thread
is that the proxy has loopholes the designer never anticipated.
\end{intbox}

\textbf{Reward overoptimization} is the quantitative version: as $r_\theta$ is optimized
harder, true performance first rises, then \emph{degrades}, because the policy finds
adversarial inputs that score high without being genuinely good. The standard mitigation
is a KL penalty to a reference policy,
\[
  \max_\pi\; \E[r_\theta] - \beta\,\KL{\pi}{\pi_{\text{ref}}},
\]
which keeps the policy near the distribution on which the reward model was trained.
The core takeaway of data-centric RL: \emph{the reward is a hypothesis about what the
agent should achieve --- like any hypothesis, treat it with skepticism and validate it.}


% ==========================================================
%  CHAPTER 9 — PLANNING WITH A MODEL
% ==========================================================
\chapter{Planning with a Model}
\label{chap:planning}

Humans and animals \emph{simulate} future outcomes before acting: a chess player explores
candidate moves mentally, a robot plans a collision-free trajectory. Every method so far
has been \emph{model-free} --- the agent observes $(s,a,r,s')$ tuples and updates value or
policy estimates directly, with no explicit model of the world. This is unbiased but
sample-hungry, and each real interaction can be expensive in time, hardware, or safety.
What if the agent could think ahead using an internal model? This chapter assumes we
\emph{have} a model (given or learned) and asks how to \emph{plan} with it;
\cref{chap:mbrl} asks how to learn one.

\section{What Is a Model?}

A \textbf{model} of the environment is a function that predicts what happens when the
agent acts: a transition model $s_{t+1} \sim p_\phi(s_{t+1}\mid s_t, a_t)$ and a reward
model $\hat r = r_\phi(s_t, a_t)$. Together they define a \emph{simulated MDP}: given any
$(s,a)$, the agent can produce a prediction $(s', r)$ without touching the real
environment.

Models vary along several axes. Their \emph{origin} can be \emph{known} (given exactly,
e.g.\ the rules of chess or Go) or \emph{learned} from data. Their \emph{output} can be
deterministic, $s' = f(s,a)$, or stochastic, $s' \sim p(\cdot\mid s,a)$. They can operate
in the raw observation space or in a learned \emph{latent} space (as in MuZero). And they
can be full or only local/partial. A model unlocks four capabilities: planning (search
over action sequences), synthetic Dyna-style rollouts, policy gradients backpropagated
through the model, and uncertainty-guided exploration. This chapter focuses on planning.

\section{Planning vs.\ Learning}

\emph{Model-free learning} interacts with the real environment and updates value/policy
from real transitions: unbiased but sample-hungry. \emph{Planning} uses a model to
simulate transitions and search over imagined trajectories: fast once the model is known,
but model errors propagate. The two are \emph{not} mutually exclusive --- the most powerful
agents do both. Model-free methods often need millions of environment steps; model-based
methods can reach comparable performance with orders of magnitude fewer real interactions,
especially when the environment is expensive to simulate and the dynamics are accurate
(low-dimensional, smooth). The caveat: model errors can cause catastrophic policy updates.

\section{Dynamic Programming as Planning}

Dynamic programming (\cref{chap:tabular}) already \emph{is} planning: it assumes the MDP
$p(s'\mid s,a),\, r(s,a)$ is known and uses it to improve values or policies via Bellman
backups, with no new environment data. The two classic control algorithms:
\[
  \text{Value iteration:}\quad
  V_{k+1}(s) = \max_a\!\left[r(s,a) + \gamma \sum_{s'} p(s'\mid s,a) V_k(s')\right] \;\forall s,
\]
\[
  \text{Policy iteration:}\quad
  V^{\pi_k}(s) = \sum_a \pi_k(a\mid s)\!\left[r(s,a) + \gamma\sum_{s'} p(s'\mid s,a) V^{\pi_k}(s')\right],
\]
followed by greedy improvement $\pi_{k+1}(s) = \argmax_a[\,r(s,a) + \gamma\sum_{s'}p(s'\mid s,a)V^{\pi_k}(s')\,]$.

But DP plans over the \emph{entire} MDP: every state, every action, every possible next
state. It requires an explicit model, performs full-state sweeps, sums over all next
states at each backup, and offers no generalization in tabular form (similar states are
unrelated table entries). For large problems this is hopeless. The key shift for the rest
of the chapter: \emph{instead of solving the whole MDP, plan only from the current state.}

\section{Forward Search and Simulation-Based Search}

\textbf{Forward search} selects the best action by lookahead from the current state. It
builds a search tree with $s_t$ at the root, uses the model to look ahead, and solves only
the sub-MDP starting from \emph{now} --- there is no need to solve the whole MDP.

\textbf{Simulation-based search} is the sample-based variant: simulate trajectories from
$s_t$ using the model $\hat p_\phi$ and a policy $\pi$, then apply \emph{any} model-free
method (Monte Carlo, SARSA, Q-learning) to the simulated data. Simulation is using a model
to generate imaginary experience. Concretely, Monte Carlo simulation search rolls out $K$
full episodes from $s_t$ and estimates local values from the mean simulated return:
\[
  \hat v(s_t) = \frac{1}{K}\sum_{k=1}^K G_t^k \approx v^\pi(s_t),
  \qquad
  \hat q(s_t, a) = \frac{1}{K}\sum_{k=1}^K G_t^k \approx q^\pi(s_t, a),
  \quad A_t = \argmax_a \hat q(s_t, a).
\]
No global value function is needed --- estimates are \emph{local} to $s_t$ --- and the
model can be a black box that only needs to generate samples $(r, s')$.

\begin{notebox}[title={The branching problem}]
Full tree expansion explodes. At depth $d$ with branching factor $b = \abs{\Aset}$, the
tree has $b^d$ nodes. Chess: $b\approx 35$, $d\approx 80 \Rightarrow 35^{80}\approx 10^{123}$.
Go: $b\approx 250$, $d\approx 150$, far beyond any computer. Two classical remedies:
\emph{depth reduction} (replace deep rollouts with a value function $\hat V(s)$) and
\emph{breadth reduction} (only expand promising branches). MCTS addresses both.
\end{notebox}

\section{Monte Carlo Tree Search}

\textbf{MCTS} builds a \emph{partial} lookahead tree using Monte Carlo simulations, growing
it asymmetrically toward promising regions. Four phases repeat many times:
\begin{enumerate}
  \item \textbf{Selection:} from the root, traverse to a leaf using a \emph{tree policy} (UCT).
  \item \textbf{Expansion:} add one or more child nodes to the selected leaf.
  \item \textbf{Simulation:} run a rollout from the new node to a terminal state using a fast
    \emph{default policy} (e.g.\ random).
  \item \textbf{Backpropagation:} update visit counts $N(s,a)$ and value estimates $Q(s,a)$
    along the path back to the root.
\end{enumerate}
After the budget is exhausted, play the action with the \emph{highest visit count} at the
root. The tree grows asymmetrically (promising subtrees are visited more), simulations give
unbiased value estimates with no value function needed, and over time estimates converge to
the true minimax value.

\subsection{UCT: the tree policy}

How do we decide which node to expand? \textbf{UCT} (Upper Confidence Bounds for Trees)
applies UCB1 at each node, balancing exploitation against exploration:
\[
  a^* = \argmax_a\!\left[\, Q(s,a) + c\sqrt{\frac{\ln N(s)}{N(s,a)}}\,\right].
\]
Here $Q(s,a)$ is the mean return from simulations through $(s,a)$, $N(s)$ the total visits
to $s$, $N(s,a)$ the visits via action $a$, and $c$ an exploration constant. High $Q$
exploits promising branches; a low $N(s,a)$ relative to $N(s)$ explores under-visited ones.
UCT converges to the minimax-optimal action as the number of simulations $\to \infty$.
MCTS is \emph{anytime} (more simulations improve the decision), \emph{selective} (computation
focuses on promising branches), and \emph{black-box} (needs only a simulator).

\subsection{From rollouts to learned value functions}

MCTS's bottleneck is leaf evaluation. The classical Monte Carlo rollout
$\hat V(s_{\text{leaf}}) = \tfrac{1}{K}\sum_k G^k(s_{\text{leaf}})$ is poor in complex
domains: a random rollout policy is a weak evaluator, each rollout costs $\bigO{d}$ model
queries, and long rollouts have high variance. The fix is to replace expensive rollouts
with a \emph{learned value function} $V_\theta(s_{\text{leaf}})$, trained on game outcomes,
Monte Carlo returns, TD targets, or self-play. $V_\theta$ gives $\bigO{1}$ leaf evaluation,
generalizes across states, and has lower variance, while MCTS still provides local
decision-time search and an improved policy at the root. \textbf{Neural MCTS} combines
search for local improvement with function approximation for generalization.

\section{AlphaGo and AlphaZero}

Go has simple rules but $\sim 10^{170}$ legal positions, an enormous branching factor, and
no simple hand-designed evaluation function --- classical minimax search is infeasible. It
needs both selective search and learned heuristics.

\textbf{AlphaGo} (2016) combined MCTS with three learned networks: a large policy network
$p_\sigma(a\mid s)$ (trained by supervised learning on expert moves, then improved by
self-play policy gradients, used as a prior over promising actions in the tree), a fast
rollout policy $p_\pi$ (small, used to simulate from leaves), and a value network
$v_\theta(s)$ (trained on self-play outcomes to estimate the win probability and evaluate
leaves). Its tree policy mixed the policy prior into selection and evaluated leaves by
blending the value network with a rollout outcome.

\textbf{AlphaZero} (2017) keeps the core idea --- MCTS + policy/value networks + self-play
--- but removes the human expert games, the separate rollout policy, and the supervised
pre-training, and uses a \emph{single} network predicting both policy and value,
$f_\theta(s) = (\mathbf{p}, v)$.

\begin{intbox}[title={Search improves the policy; learning amortizes search}]
At each expanded node the network gives a prior $\mathbf{p}$ and a value $v$; MCTS uses them
to produce an improved search policy $\pi_{\text{MCTS}}(a\mid s) \propto N(s,a)^{1/\tau}$.
Self-play with MCTS at every state generates training targets $(s_t,\, \pi_{\text{MCTS}},\, z)$
where $z$ is the final game outcome. The network is then trained to predict both:
\[
  \loss(\theta) = \underbrace{(v_\theta(s_t) - z)^2}_{\text{value loss}}
  - \underbrace{\pi_{\text{MCTS}}\T \log \mathbf{p}_\theta}_{\text{policy loss}}
  + \underbrace{c\norm{\theta}^2}_{\text{regularization}}.
\]
\textbf{MCTS is a policy-improvement operator}: it turns the raw network policy
$\mathbf{p}_\theta$ into the stronger $\pi_{\text{MCTS}}$, and the network is trained to
\emph{amortize} that improved search. Self-play provides an automatic curriculum --- the
opponent improves alongside the agent --- and values are always read from the current
player's perspective, inducing minimax-like reasoning.
\end{intbox}

\section{Trajectory Optimisation}

Tree search works for discrete actions, but in continuous control, branching over all
actions is impossible. Instead, plan by directly optimizing an \emph{action sequence}:
\[
  \mathbf{a}^*_{0:H-1} = \argmax_{\mathbf{a}_{0:H-1}} \sum_{t=0}^{H-1} r(s_t, a_t),
  \qquad s_{t+1} = f(s_t, a_t),
\]
where $f$ is a queryable dynamics model and $H$ the planning horizon. Trajectory
optimization searches over action sequences, not tree nodes.

\textbf{Model Predictive Control (MPC)} plans over a finite horizon but executes only the
\emph{first} action, then re-plans:
\[
  \mathbf{a}^*_{t:t+H-1} = \argmax_{\mathbf{a}_{t:t+H-1}} \sum_{\tau=t}^{t+H-1} r(s_\tau, a_\tau),
  \qquad \text{execute } a_t = a^*_t,\ \text{observe } s_{t+1},\ \text{re-plan}.
\]
Re-planning after each real transition corrects for model errors, avoids committing to a
long open-loop plan, and naturally handles changing situations.

How do we solve the inner optimization? Two derivative-free methods:
\begin{itemize}
  \item \textbf{Random shooting:} sample $N$ action sequences uniformly, simulate each under
    the model, and execute the first action of the best one. Simple and trivially
    parallelizable, needs no gradients, but covers high-dimensional action spaces poorly.
  \item \textbf{Cross-Entropy Method (CEM):} iteratively fit the sampling distribution to the
    best sequences. Start from $\Normal(\mathbf{0}, \sigma^2\eye)$ over action sequences;
    each iteration sample $N$ sequences, score them, keep the top elites $\mathcal{E}$, and
    refit $\mu \leftarrow \E_\mathcal{E}[\mathbf{a}]$, $\Sigma \leftarrow \Cov_\mathcal{E}[\mathbf{a}]$.
    Robust, works with non-differentiable models, effective in low-to-medium action
    dimensions, but the Gaussian assumption can miss multimodal solutions and scales poorly
    to very high dimensions.
\end{itemize}

Both tree search and trajectory optimisation are \emph{online planners}: they spend
computation at decision time using a model. Tree search suits discrete actions and deep,
branching lookahead with UCT/value-net guidance; shooting and CEM suit continuous actions
and short, MPC-style horizons.

\begin{notebox}[title={What changes when the model is learned}]
Everything here assumed a \emph{queryable} model $s_{t+1} = f(s_t,a_t)$. When the model is
learned from data, $s_{t+1}\sim p_\phi(\cdot\mid s_t,a_t)$, four problems appear: model
error, compounding error over rollouts, distribution shift (the planner queries unfamiliar
states), and model exploitation (the policy exploits model mistakes). \Cref{chap:mbrl}
tackles these.
\end{notebox}


% ==========================================================
%  CHAPTER 10 — LEARNING MODELS (MODEL-BASED RL II)
% ==========================================================
\chapter{Learning Models}
\label{chap:mbrl}

\Cref{chap:planning} assumed a queryable model: dynamic programming, tree search,
AlphaZero, and CEM all needed $p(s'\mid s,a)$ and $r(s,a)$. In most real problems the
model is \emph{not} given --- it must be learned from data. This chapter is about
learning a dynamics model, the characteristic ways learned models fail, and the
algorithms (Dyna, MBPO, PETS, latent world models) that use a learned model while
controlling its errors.

\begin{intbox}[title={The model-based RL loop}]
Model-free RL goes \emph{experience $\to$ value/policy} directly. Model-based RL inserts
a model in the middle: \emph{experience $\to$ model}, then \emph{model $\to$ plan/imagine
$\to$ value/policy}. The agent collects transitions, fits a model
$\hat p_\phi(s'\mid s,a),\, \hat r_\phi(s,a)$, plans or imagines with it, and updates the
policy from real \emph{and} imagined data. The model is the extra source of training data
that buys sample efficiency.
\end{intbox}

\section{Learning a Model from Transitions}

The minimal model for one-step rollouts is a transition model (deterministic
$s_{t+1}\approx f_\phi(s_t,a_t)$ or stochastic $s_{t+1}\sim p_\phi(\cdot\mid s_t,a_t)$)
plus a reward model $\hat r_t = r_\phi(s_t,a_t)$. Both are trained by ordinary supervised
learning on a data set $\dataset = \{(s_t,a_t,r_t,s_{t+1})\}$:
\[
  \loss_{\text{trans}}(\phi) = \frac{1}{\abs{\dataset}}\sum \norm{s' - f_\phi(s,a)}^2
  \quad\text{(or }-\log p_\phi(s'\mid s,a)\text{)},
  \qquad
  \loss_{\text{rew}}(\phi) = \frac{1}{\abs{\dataset}}\sum \bigl(r - r_\phi(s,a)\bigr)^2.
\]
Architecture follows the observation type: an MLP for low-dimensional states, a CNN for
images, an RNN for partial observability. The reward model is often simpler (scalar
output) and sometimes known analytically.

The quality of a learned model depends critically on \emph{which} data is collected. A
random policy gives broad but mostly irrelevant coverage; the current policy gives data
where the agent actually operates (accurate where it matters); a replay-buffer mixture
keeps all past transitions; and active, uncertainty-driven collection queries the
environment where the model is least confident.

\section{Why Learned Models Fail}

\begin{notebox}[title={Three failure modes of learned models}]
\textbf{1. Distributional shift.} The model is trained on the data-collection policy, but
the planning policy visits states the model has never seen,
$(s,a)\sim\pi_{\text{plan}} \neq (s,a)\sim\pi_{\text{data}}$. Model errors are undefined
off-distribution and the planner confidently exploits artefacts --- the same covariate
shift as imitation learning (recall DAgger). \emph{Fix:} re-collect data under the
planning policy (iterative MBRL); uncertainty-aware planning.

\textbf{2. Compounding errors.} A per-step error $\epsilon$ accumulates over a rollout:
total error $\le H\epsilon$ (linear, or worse for divergent dynamics). With
$\epsilon\approx 0.01$ and $H=100$ the error reaches the scale of the state itself; value
estimates on such trajectories are meaningless. \emph{Fix:} short rollouts, mix with real
data, penalize uncertainty.

\textbf{3. Model exploitation.} A policy optimized against a learned model finds its
\emph{weaknesses}. The classic example is a robot that discovers a model glitch and
vibrates indefinitely to accumulate infinite simulated reward. \emph{Fix:} limit the
planning horizon, plan pessimistically (penalize uncertainty), frequently re-train on new
real data.
\end{notebox}

These pressures create an \textbf{accuracy--usability trade-off}. A short horizon ($H$
small) keeps compounding error low and predictions accurate but limits planning depth and
needs a good value function at the leaf. A long horizon gives a richer learning signal but
suffers high compounding error and likely model exploitation, demanding a very accurate
model. The practical recommendation: \emph{keep $H$ short (5--20 steps) and combine
imagined data with real data.}

\section{Re-planning at Every Step: MPC + CEM}

The simplest way to use a learned model is to plan with it at decision time and re-plan
continuously, exactly as in \cref{chap:planning}. \textbf{MPC} runs a receding-horizon
loop: at state $s_t$, optimize an action sequence under the learned model, execute only
the first action, observe $s_{t+1}$, and repeat:
\[
  \mathbf{a}^* = \argmax_{a_t,\dots,a_{t+H-1}} \sum_{k=0}^{H-1} \gamma^k\, \hat r(s_{t+k}, a_{t+k}).
\]
Re-planning at each step \emph{corrects for model errors} but is computationally expensive
(an inner optimization every step) and works with any black-box model. The inner optimizer
is usually \textbf{CEM}: represent each action as an independent Gaussian
$p(a_t) = \Normal(\mu_t, \sigma_t^2\eye)$; sample $K$ sequences, simulate each under
$\hat p_\phi$, keep the top-$N$ elites, refit $\mu_t,\sigma_t$ to the elites, iterate, and
execute $\mu_0$. CEM is derivative-free and parallelizable but uses a separate distribution
per time step (trajectories need not be smooth or physically consistent) and keeps no
memory between MPC steps (warm-starting helps).

\section{Dyna: Learning from Imagined Transitions}

\textbf{Dyna} (Sutton, 1991) is the early, influential framework for combining model
learning with model-free RL. Its key insight: \emph{imagined transitions from a learned
model can be treated like ordinary experience.} Real experience updates the model
\emph{and} the value/policy; imagined experience updates the value/policy only. Each real
transition can be followed by $K$ model-generated updates, so a good model greatly improves
sample efficiency --- and crucially, the model need \emph{not} be used for explicit
lookahead planning.

\begin{thmbox}[title={Dyna-Q}]
Given a $Q$-function, learned model $\hat p_\phi$, and replay buffer $\dataset$:
\begin{enumerate}
  \item \textbf{Act:} choose $a_t \leftarrow \epsilon$-greedy$(Q, s_t)$; observe $r_t, s_{t+1}$.
  \item \textbf{Learn from real experience} $(s_t,a_t,r_t,s_{t+1})$: update $Q$, update the
    model $\hat p_\phi$, store the transition in $\dataset$.
  \item \textbf{Learn from imagined experience}, repeat $K$ times: sample
    $(s,a)\sim\dataset$, predict $(r',s')\sim\hat p_\phi(\cdot\mid s,a)$, and update $Q$ with
    the imagined transition.
  \item Advance to $s_{t+1}$ and repeat.
\end{enumerate}
\end{thmbox}

With $K=0$ this is plain model-free Q-learning; with $K>0$ each real step is supplemented by
$K$ model-generated value updates. The caveats: a wrong model produces misleading imagined
transitions, the model can become outdated in non-stationary environments, and standard Dyna
only samples \emph{previously visited} $(s,a)$ pairs (limited exploration).
\textbf{Dyna-Q+} repurposes the model for exploration, adding a bonus for state-action pairs
not tried for a long time:
\[
  r^+(s,a) = r(s,a) + \kappa\sqrt{\tau(s,a)},
\]
where $\tau(s,a)$ is the time since $a$ was last tried in $s$ --- useful when the
environment may have changed. Dyna is the blueprint for modern MBRL: real data $\to$
learned model $\to$ synthetic data $\to$ policy.

\section{Model-Based Policy Optimization (MBPO)}

Long rollouts from an imperfect model drift away from realistic trajectories. \textbf{MBPO}
combines a learned dynamics model with an off-policy actor-critic (SAC) and makes two design
choices to keep the synthetic data trustworthy: \emph{short model rollouts starting from
real states} in the replay buffer, and \emph{mixed replay} (train SAC on a mixture of real
and imagined transitions). Start from states the model has actually seen, imagine only a
few steps, and use the synthetic transitions to improve the policy.

\begin{thmbox}[title={MBPO}]
Given policy $\pi_\theta$, model ensemble $\{\hat p_{\phi_i}\}$, real buffer
$\dataset_{\text{real}}$, model buffer $\dataset_{\text{model}}$:
\begin{enumerate}
  \item \textbf{Collect real data:} roll out $\pi_\theta$ in the environment; add to
    $\dataset_{\text{real}}$.
  \item \textbf{Train the model:} fit the ensemble on $\dataset_{\text{real}}$.
  \item \textbf{Generate imagined data:} sample $s_0\sim\dataset_{\text{real}}$ and roll out
    $\pi_\theta$ for $k$ \emph{short} model steps; add transitions to $\dataset_{\text{model}}$.
  \item \textbf{Update the policy:} train SAC on minibatches from
    $\dataset_{\text{real}}\cup\dataset_{\text{model}}$.
  \item Repeat.
\end{enumerate}
Key hyperparameter: rollout length $k$, typically very short (1--5 steps), sometimes even
$k=1$.
\end{thmbox}

The rollout length balances two effects: $k=0$ has no model bias but no synthetic data;
small $k$ gives useful extra data with limited model error; large $k$ gives more imagined
experience but compounding error. \emph{The goal is not to build a perfect simulator --- it
is to use the model only where it is accurate enough.} MBPO matches strong model-free SAC
with far fewer real steps on MuJoCo locomotion (often $\sim 20\times$ more efficient), and
the model \emph{ensemble} is essential for reducing bias and estimating uncertainty. Seen
this way, MBPO is a modern, error-controlled Dyna-style algorithm: where Dyna-Q uses
Q-learning, simple models, one-step samples, and replayed start states, MBPO uses SAC, a
neural ensemble, short rollouts, and real replay states.

\section{Probabilistic Models and Ensembles}

A deterministic model gives a single prediction $s' = f_\phi(s,a)$ and cannot say how
reliable it is. A useful model should predict not just \emph{what} may happen but \emph{how
reliable} the prediction is. There are two distinct sources of uncertainty:

\begin{center}
\begin{tabular}{@{}lll@{}}
\toprule
\textbf{Type} & \textbf{Source} & \textbf{Reducible with more data?} \\
\midrule
Aleatoric & inherent environment randomness (noisy sensors, stochastic transitions) & no \\
Epistemic & limited knowledge (model queried far from training data) & yes \\
\bottomrule
\end{tabular}
\end{center}
Aleatoric uncertainty is part of the \emph{environment}; epistemic uncertainty is part of
the \emph{model}.

A \textbf{Gaussian neural network} outputs a mean and variance,
$p_\phi(s'\mid s,a) = \Normal(\mu_\phi(s,a), \Sigma_\phi(s,a))$, trained by maximizing the
likelihood of observed transitions. The mean predicts the most likely next state and the
variance captures noisy or ambiguous outcomes --- it models \emph{aleatoric} uncertainty.
But a large predictive variance is not the same as knowing the model is outside its
training data: predictive variance does not capture \emph{epistemic} uncertainty. For that
we need model \emph{disagreement}.

\textbf{Deep ensembles} train $M$ independent models $\{p_{\phi_i}\}_{i=1}^M$ on the same
data with different seeds, initializations, or bootstraps. If they agree, the prediction is
reliable; if they disagree, the query is likely out-of-distribution. For Gaussian members,
the total predictive variance decomposes cleanly:
\[
  \Var_{\text{ens}}[s'] =
  \underbrace{\frac{1}{M}\sum_i \Sigma_{\phi_i}}_{\text{within-model: aleatoric}}
  + \underbrace{\frac{1}{M}\sum_i (\mu_{\phi_i} - \bar\mu)^2}_{\text{between-model: epistemic}}.
\]
In MBRL epistemic uncertainty is the important one --- it tells the planner where the model
\emph{should not be trusted}. Ensembles are then used for model selection (sample one model
per imagined rollout), uncertainty-aware planning (avoid high-disagreement actions),
pessimistic value estimates (penalize uncertain regions), and active data collection.

\subsection{PETS: planning with an ensemble}

\textbf{PETS} (Probabilistic Ensembles with Trajectory Sampling) combines all the pieces for
continuous control: a probabilistic ensemble for the dynamics, CEM for action-sequence
search, and \emph{trajectory sampling} to propagate uncertainty. During CEM planning, each
candidate action sequence is evaluated by rolling out several \emph{particles} through the
ensemble. Two schemes: TS1 keeps one model fixed per rollout (a single coherent hypothesis),
while TS$\infty$ samples a fresh ensemble member at each step. Trajectory sampling makes
disagreement between models actually affect planning. In short, PETS is MPC with CEM where
candidate sequences are scored by averaging the returns of particles rolled out through the
probabilistic ensemble.

\section{Latent World Models}

So far the model was a one-step transition predictor over the raw state,
$(s_{t+1}, r_t)\sim p_\phi(\cdot\mid s_t,a_t)$, assuming $s_t$ is already a good
representation. A \textbf{world model} instead learns an internal state representation and
predicts future experience from it:
\[
  z_t = e_\phi(o_{\le t}, a_{<t}), \qquad (z_{t+1}, r_t)\sim p_\phi(\cdot\mid z_t, a_t).
\]
A world model is not just a next-state predictor --- it learns \emph{what state information
the agent should carry forward} for imagination and decision making.

Raw observations are often the wrong state for control: high-dimensional (images, sensors,
text), partial (one frame need not reveal the full state), full of irrelevant variation
(background, lighting, noise), or missing history-dependent information (velocity, memory,
hidden variables). A latent state $z_t$ summarizes the useful information from the history.
\emph{The goal is not perfect reconstruction --- it is a representation useful for control.}
Examples: \textbf{World Models} (Ha \& Schmidhuber, 2018) learn latent dynamics and optimize
a controller in imagined rollouts; \textbf{Dreamer} learns recurrent latent dynamics and
trains an actor-critic purely in imagination; \textbf{MuZero} learns \emph{value-equivalent}
latent dynamics and runs MCTS in latent space. Same high-level idea: learn an internal
model, then use imagined futures to improve decisions.

\section{When Does Model-Based RL Help?}

MBRL works best when real data is expensive (robotics, chemistry, clinical trials), the
dynamics are smooth and learnable (continuous locomotion, manipulation), short planning
horizons suffice (so model errors stay bounded), and the reward is dense (easy to fit a
reward model). It is harder when observations are high-dimensional and complex (raw pixels
need a latent model), dynamics are discontinuous or stochastic (contacts, multi-agent),
the reward is sparse, or distribution shift is severe (rapidly changing environment).
Across all of these, the recurring strategy is the same: \emph{trust the model only where
it is accurate, and keep correcting it with real data.}


% ==========================================================
%  CHAPTER 11 — RL FOR FOUNDATION MODELS
% ==========================================================
\chapter{RL for Foundation Models}
\label{chap:fm}

Two trends are reshaping AI. \textbf{Foundation models} (LLMs, vision-language, multimodal)
are trained on large-scale data, carry strong general-purpose representations, and are
adapted by prompting, supervised fine-tuning, or post-training. \textbf{Agentic AI} systems
act, use tools, and interact with environments through multi-step reasoning and feedback.
This chapter shows that the RL vocabulary of the whole course --- behavior cloning,
preference learning, policy gradients, KL constraints, reward hacking, model-based thinking
--- reappears almost unchanged when a foundation model becomes a decision-making system.

\section{Autoregressive Generation as Sequential Decision-Making}

An autoregressive language model factorizes a response as
$p_\theta(y\mid x) = \prod_{t=1}^T p_\theta(y_t\mid x, y_{<t})$ and is trained by maximum
likelihood. This is \emph{exactly} a policy: reading it as RL,
\[
  \pi_\theta(a_t\mid s_t) = p_\theta(y_t\mid x, y_{<t}),
\]
the \emph{state} is the prompt plus the generated prefix, the \emph{action} is the next
token (or a tool call), and the \emph{reward} is a preference, a verifier, or a task-success
signal. Temperature, top-$p$, and decoding rules are part of the effective policy. The same
network can act as a text simulator or as an environment actor.

\begin{notebox}[title={What makes language RL difficult}]
\textbf{Huge discrete action space} --- tens or hundreds of thousands of tokens.
\textbf{Long horizons} --- early words shape later possibilities.
\textbf{Sparse sequence-level rewards} --- the answer is judged only after generation.
\textbf{Partial, conflicting objectives} --- helpfulness, truthfulness, style, safety, and
brevity can clash.
\textbf{Distribution shift} --- optimizing the model changes the very text people see and
rate. The good news: we already studied every ingredient needed to cope.
\end{notebox}

\section{LLM Training: Pretraining, SFT, Preferences}

\textbf{Pretraining is behavior cloning from text.} Maximizing
$\sum_{(x,y)\in\dataset_{\text{text}}} \sum_t \log \pi_\theta(y_t\mid x, y_{<t})$ imitates
the conditional distribution of internet, book, and code text. This is powerful because the
data contains many latent skills --- but it is \emph{not} the same as optimizing for being a
helpful assistant. Pretraining learns plausible continuations and many \emph{conflicting}
behaviors; users want instruction-following, calibrated answers, refusal in unsafe cases,
and consistency with preferences.

\textbf{Supervised fine-tuning (SFT)} is instruction-following imitation: collect
demonstrations $(x, y^*)\sim\dataset_{\text{demo}}$ and fine-tune with the same token-level
likelihood objective. SFT is simple, stable, cheap relative to RL, and usually the first
step toward an assistant --- but it says ``copy this behavior,'' not ``search for outputs
humans prefer.''

To go further we need a preference signal. Humans find \emph{comparison} easier than writing
the ideal answer, so we collect pairs and train a reward model with the Bradley--Terry
objective (\cref{chap:irl}):
\[
  P_\phi(y_w \succ y_l\mid x) = \sigma\bigl(r_\phi(x, y_w) - r_\phi(x, y_l)\bigr),
  \qquad
  \loss_{\text{RM}}(\phi) = -\E_{(x,y_w,y_l)}\bigl[\log\sigma\bigl(r_\phi(x,y_w) - r_\phi(x,y_l)\bigr)\bigr].
\]
The reward model scores complete responses; it is a learned \emph{proxy} for human
preference, not ground truth.

\section{RLHF Objective and PPO}

RLHF optimizes the policy against the learned reward while staying close to the SFT
reference, via a KL-regularized objective:
\[
  \max_\theta\; \E_{y\sim\pi_\theta(\cdot\mid x)}\!\left[\,
    r_\phi(x, y) - \beta \log \frac{\pi_\theta(y\mid x)}{\pi_{\text{ref}}(y\mid x)}\right].
\]
Here $r_\phi$ is the reward model and $\pi_{\text{ref}}$ is usually the SFT policy; the KL
penalty keeps the policy near the reference. This is the \emph{target} objective; PPO is one
way to optimize it approximately, with the clipped policy-gradient surrogate
\[
  L^{\text{PPO}}(\theta) = \E\bigl[\min\bigl(\rho_t(\theta) A_t,\, \mathrm{clip}(\rho_t(\theta), 1-\epsilon, 1+\epsilon) A_t\bigr)\bigr],
  \qquad
  \rho_t(\theta) = \frac{\pi_\theta(y_t\mid x, y_{<t})}{\pi_{\theta_{\text{old}}}(y_t\mid x, y_{<t})},
\]
where the advantage $A_t$ is estimated from rewards and a learned value function.

\begin{intbox}[title={Two anchors in PPO-style RLHF}]
PPO-RLHF holds the policy between \emph{two} reference points that do different jobs. The
\textbf{KL penalty} compares $\pi_\theta$ to $\pi_{\text{ref}}$ (the SFT model) --- it keeps
the \emph{content} close to the trusted assistant and curbs reward overoptimization. The
\textbf{PPO clipping} compares $\pi_\theta$ to $\pi_{\theta_{\text{old}}}$ (the previous
policy) --- it keeps each \emph{update step} stable. Don't conflate them: one anchors where
you end up, the other anchors how far you move per step.
\end{intbox}

\subsection{Group Relative Policy Optimization (GRPO)}

PPO-RLHF needs a separate value network for $A_t$, which is costly at LLM scale.
\textbf{GRPO} removes it. Sample a \emph{group} of $G$ outputs for the same prompt, score
them all, and convert each reward into a \emph{group-relative advantage} by standardizing
within the group:
\[
  \hat A_i = \frac{r_i - \mathrm{mean}(r_{1:G})}{\mathrm{std}(r_{1:G}) + \delta}.
\]
The single sequence-level advantage $\hat A_i$ is applied to all tokens of output $y_i$, and
the update keeps the PPO-style clip plus a token-level KL penalty to the reference:
\[
  L^{\text{GRPO}}(\theta) = \E\!\left[\frac{1}{G}\sum_{i=1}^G \frac{1}{\abs{y_i}}
    \sum_t \Bigl(\min\bigl(\rho_{i,t}\hat A_i,\, \mathrm{clip}(\rho_{i,t}, 1-\epsilon, 1+\epsilon)\hat A_i\bigr)
    - \beta\, D_{\mathrm{KL},i,t}\Bigr)\right].
\]
Using the group as its own baseline avoids a learned value function, reducing memory and
compute --- which is why GRPO is central to recent reasoning-model training.

\section{RLHF Failure Modes and Verifiable Rewards}

RLHF inherits every reward-proxy pathology from \cref{chap:irl}: \emph{reward hacking}
(optimize what the reward model likes, not human intent), \emph{overoptimization} (quality
drops after too many RL steps), \emph{preference misspecification} (style or confidence
rewarded over correctness), \emph{mode collapse} (generic safe-sounding answers become
attractive), and \emph{distribution shift} (the reward model scores outputs unlike its
training data). Learned reward models are useful but imperfect proxies.

When the answer can be checked automatically, we can sidestep the learned proxy.
\textbf{RLVR} (Reinforcement Learning with Verifiable Rewards) replaces preference rewards
with \emph{automatic checkers}: math (final answer matches a verifier), code (tests pass),
tool use (an API result satisfies a check), games (win/loss observed), or structured domains
(symbolic constraints satisfied). Cheap verifier rewards allow many sampled attempts per
prompt and have been central to recent reasoning models and specialized agents. Reward
design matters: an answer must be both \emph{checkable} and \emph{correct}, so a composite
reward teaches both format and correctness,
\[
  r(y) = \lambda_{\text{format}}\, r_{\text{format}}(y) + \lambda_{\text{correct}}\, r_{\text{correct}}(y),
\]
where the format reward rewards a parseable answer (JSON, tagged fields) and the correctness
reward rewards passing the verifier. Format failures otherwise cause us to underestimate the
model's true capability.

\begin{center}
\begin{tabular}{@{}llll@{}}
\toprule
\textbf{Method} & \textbf{Signal} & \textbf{Strength} & \textbf{Main risk} \\
\midrule
SFT  & demonstrations     & stable imitation         & limited by examples \\
RLHF & human preferences  & aligns subjective quality & proxy reward hacking \\
DPO  & preference pairs   & simple and stable        & offline preference limits \\
RLVR & automatic verifier & scalable objective signal & narrow rewards, hacking the verifier \\
\bottomrule
\end{tabular}
\end{center}

\section{Generalist Agents}

Once everything --- instruction text, visual observations, state observations, previous
actions --- is tokenized into a shared sequence representation, language modeling and action
prediction become structurally identical. \textbf{Gato} (Reed et al., 2022) exploits this:
one multimodal, multi-task, multi-embodiment Transformer with the \emph{same weights}
playing Atari, controlling robots, captioning images, and holding dialogue, with inputs and
outputs serialized into token sequences. Gato is deployed as a policy over observations and
actions but trained primarily by supervised sequence modeling on mixed data --- ``a policy
can be general without being generally intelligent.''

Modern agentic systems extend a generalist policy with tools, memory/retrieval, evaluators,
and control loops: the foundation model becomes a \emph{controller} that reads a task
instruction, queries memory and tools, observes environment state, and receives feedback from
tests, rewards, or users. Learning can target the model, tool use, planning, or task success,
drawing on the full menu: imitation (copy traces or stronger models), offline RL (logs of
tool use), preference optimization, verifiable rewards, and online RL (live interaction, with
strong safety filters). Pure online trial-and-error is often too expensive or risky, so
practical systems lean heavily on data, simulators, and verifiers.

\begin{intbox}[title={Text as computation: planning before acting}]
A reasoning model can \emph{spend tokens computing} before committing to an action ---
``\texttt{<think>} plan, reason, check \texttt{</think>}'' then ``\texttt{<answer>} action
\texttt{</answer>}.'' The final answer, tool call, or command is the action the environment
sees; feedback from a verifier, tool, or user can trigger another round of planning. This is
\cref{chap:planning}'s lookahead idea --- spend computation at decision time --- expressed in
natural language instead of a search tree.
\end{intbox}

The mapping back to classic RL is exact but broadened. The \emph{policy} is a foundation
model rather than a task-specific network; the \emph{state} is context, memory, screen,
files, and history rather than a fixed observation; the \emph{action} is tokens, tool calls,
clicks, code, or commands; the \emph{reward} comes from tests, users, preferences, and task
success; the \emph{environment} is software, the web, tools, and people; and \emph{training}
is a mixture of SFT, RLHF, RLVR, tools, and logs rather than RL in a single task
distribution. Generalist agents still fit the RL vocabulary --- the state, action, reward,
and environment simply become broader and less cleanly specified.

\end{document}
